% Primitives d’une fonction continue sur un intervalle — Mathématiques Tle Tech — Cours signé OnBuch+
% Programme officiel MINESEC. Compiler avec : tectonic N10-primitives-fonction-continue-intervalle.tex
\newcommand{\DOCMATIERE}{Mathématiques}
\newcommand{\DOCNIVEAU}{Tle Tech}
\newcommand{\DOCMODULE}{Mathématiques – classe de Terminale technique}
\newcommand{\DOCLECON}{N10}
\newcommand{\DOCTITRE}{Primitives d’une fonction continue sur un intervalle}
% OnBuch+ — préambule « cours signés » ENRICHI (compilé avec tectonic / XeLaTeX)
% Système visuel « Pop » d'OnBuch+ : fond crème (jamais blanc), bordures encre
% épaisses, ombres « dures » décalées, titres Archivo Black en capitales.
% Version MATHÉMATIQUES : graphes (pgfplots/tikz), boîtes pédagogiques
% (définition, propriété, méthode, attention, le savais-tu, exercice résolu,
% exercice, TP, bilan), page de garde et signature OnBuch+ ancrée en bas.
%
% Macros attendues dans main.tex AVANT \input{preamble} :
%   \DOCMATIERE  \DOCNIVEAU  \DOCMODULE  \DOCLECON (numéro)  \DOCTITRE
\documentclass[11pt]{article}

\usepackage{fontspec}
\usepackage[french]{babel}
\usepackage[a4paper,margin=2cm,top=3.4cm,bottom=2.8cm,headheight=1.4cm,footskip=1.3cm]{geometry}
\usepackage{amsmath,amssymb,amsfonts}
\usepackage{mathtools} % \xmapsto, \coloneqq... souvent utilisés par les modèles
\usepackage{cancel} % simplifications barrées
% \abs{z} (module, valeur absolue) et \norm{u} (norme), utilisés par les modèles
\providecommand{\abs}{}\let\abs\relax\DeclarePairedDelimiter{\abs}{\lvert}{\rvert}
\providecommand{\norm}{}\let\norm\relax\DeclarePairedDelimiter{\norm}{\lVert}{\rVert}
\usepackage[table]{xcolor} % [table] : fournit aussi \rowcolors (alternance de lignes)
\usepackage{pagecolor}
\usepackage{tikz}
\usetikzlibrary{arrows.meta,positioning,calc,shapes.geometric,shapes.misc,shapes.arrows,decorations.pathmorphing,decorations.pathreplacing,decorations.markings,patterns,shadows,mindmap,trees,fit,backgrounds,matrix,intersections,angles,quotes}
\usepackage{pgfplots}
\usepackage{pgf-pie} % diagrammes circulaires \pie (statistiques)
\pgfplotsset{compat=1.18}
\usepgfplotslibrary{fillbetween,groupplots} % aires entre courbes (calcul intégral)
\usepackage{enumitem}
\usepackage{fancyhdr}
% [most] charge toutes les bibliothèques tcolorbox (listings, minted, raster,
% documentation...) et gonfle inutilement la mémoire TeX sur un document long
% (~300 boîtes) : on ne charge que ce qui est réellement utilisé.
\usepackage[skins,breakable]{tcolorbox}
\usepackage{graphicx}
\usepackage{colortbl}
\usepackage{booktabs}
\usepackage{tabularx}
\usepackage{diagbox} % case d'en-tête diagonale des tableaux à double entrée
% Colonnes extensibles centrée / gauche / droite (C, L, R), souvent utilisées par les modèles
\newcolumntype{C}{>{\centering\arraybackslash}X}
\newcolumntype{L}{>{\raggedright\arraybackslash}X}
\newcolumntype{R}{>{\raggedleft\arraybackslash}X}
\usepackage{array}
\usepackage{multicol}
\usepackage{siunitx}
% parse-numbers=false : accepte aussi \SI{2,5 \times 10^{-3}}{...}
\sisetup{locale=FR,output-decimal-marker={,},group-minimum-digits=4,parse-numbers=false}
\DeclareSIUnit\ohm{\ensuremath{\symup{\Omega}}} % U+2126 absent de la police
\usepackage{qrcode}
% \degree : commande fréquemment utilisée par les modèles pour un angle ou une
% température en dehors de \SI{}{} (non fournie par les paquets chargés).
\providecommand{\degree}{^{\circ}}
% \qed (fin de démonstration), utilisé par les modèles sans amsthm
\providecommand{\qed}{\hfill$\square$}
% \barre : double barre des valeurs interdites dans un tableau de variations (tkz-tab)
\providecommand{\barre}{\ensuremath{\|}}
% environnement proof (amsthm non chargé) utilisé par les modèles
\ifdefined\proof\else\newenvironment{proof}[1][Démonstration]{\par\noindent\textit{#1.}\ }{\qed\par}\fi
\usepackage{lastpage}
\usepackage{hyperref}
\hypersetup{hidelinks}

% ============================================================
% Palette « Pop » OnBuch+ — mêmes hex que la classe P (plus_ui.dart)
% ============================================================
\definecolor{poporange}{HTML}{F59321}
\definecolor{popdark}{HTML}{E07A0C}
\definecolor{popcream}{HTML}{FFF6E8}
\definecolor{popink}{HTML}{1C1714}
\definecolor{poppurple}{HTML}{7A5AE0}
\definecolor{popgreen}{HTML}{2E9E6A}
\definecolor{poppink}{HTML}{E0457B}
\definecolor{popblue}{HTML}{2F7DE1}
\definecolor{popgold}{HTML}{C98A12}
\definecolor{popmuted}{HTML}{8A7F76}
\definecolor{popline}{HTML}{EADFD2}
\definecolor{popgray}{HTML}{8A7F76}
\colorlet{popgrey}{popgray}
% Alias tolérants (le modèle nomme parfois les couleurs différemment)
\colorlet{popwhite}{white}
\colorlet{popblack}{popink}
\colorlet{popred}{poppink}
\colorlet{popyellow}{popgold}
% Teintes claires pour fonds de tableaux / zones
\colorlet{popblueL}{popblue!10!white}
\colorlet{poporangeL}{poporange!14!white}
\colorlet{popgreenL}{popgreen!12!white}
\colorlet{poppurpleL}{poppurple!12!white}
\colorlet{poppinkL}{poppink!12!white}
\colorlet{popgoldL}{popgold!14!white}

\pagecolor{popcream}

% ============================================================
% Typographie
% ============================================================
\defaultfontfeatures{Ligatures=TeX}
\setmainfont{PlusJakartaSans-Medium.ttf}[Path=../../../fonts/,BoldFont=PlusJakartaSans-Bold.ttf]
% Maths en sans-serif (Fira Math) avec chiffres et lettres droites de Plus
% Jakarta, pour que les formules se fondent dans le texte.
\usepackage[math-style=ISO,bold-style=ISO]{unicode-math}
\setmathfont{FiraMath-Regular.otf}[Path=../../../fonts/]
\setmathfont{PlusJakartaSans-Medium.ttf}[Path=../../../fonts/,range={up/num,up/{Latin,latin},\mathup}]
\usepackage{icomma} % virgule décimale sans espace en mode math
% Caractères mathématiques Unicode absents des polices de texte (Plus Jakarta,
% Archivo Black) : rendus par leur équivalent mathématique, en texte comme en
% formule — sinon ils s'affichent en carré vide (ex. « Arithmétique dans ℤ »).
\usepackage{newunicodechar}
\newunicodechar{ℝ}{\ensuremath{\mathbb{R}}}
\newunicodechar{ℤ}{\ensuremath{\mathbb{Z}}}
\newunicodechar{ℂ}{\ensuremath{\mathbb{C}}}
\newunicodechar{ℕ}{\ensuremath{\mathbb{N}}}
\newunicodechar{ℚ}{\ensuremath{\mathbb{Q}}}
\newunicodechar{𝔻}{\ensuremath{\mathbb{D}}}
\newunicodechar{α}{\ensuremath{\alpha}}
\newunicodechar{β}{\ensuremath{\beta}}
\newunicodechar{γ}{\ensuremath{\gamma}}
\newunicodechar{δ}{\ensuremath{\delta}}
\newunicodechar{ε}{\ensuremath{\varepsilon}}
\newunicodechar{θ}{\ensuremath{\theta}}
\newunicodechar{λ}{\ensuremath{\lambda}}
\newunicodechar{μ}{\ensuremath{\mu}}
\newunicodechar{σ}{\ensuremath{\sigma}}
\newunicodechar{τ}{\ensuremath{\tau}}
\newunicodechar{φ}{\ensuremath{\varphi}}
\newunicodechar{ψ}{\ensuremath{\psi}}
\newunicodechar{ω}{\ensuremath{\omega}}
\newunicodechar{Σ}{\ensuremath{\Sigma}}
% majuscules produites par \MakeUppercase dans les titres (α-aminés -> Α)
\newunicodechar{Α}{\ensuremath{\alpha}}
\newunicodechar{Β}{\ensuremath{\beta}}
\newunicodechar{Γ}{\ensuremath{\Gamma}}
\newunicodechar{Δ}{\ensuremath{\Delta}}
\newunicodechar{Ε}{\ensuremath{\varepsilon}}
\newunicodechar{Θ}{\ensuremath{\Theta}}
\newunicodechar{Λ}{\ensuremath{\Lambda}}
\newunicodechar{Μ}{\ensuremath{\mu}}
\newunicodechar{Τ}{\ensuremath{\tau}}
\newunicodechar{Φ}{\ensuremath{\Phi}}
\newunicodechar{Ψ}{\ensuremath{\Psi}}
\newunicodechar{Ω}{\ensuremath{\Omega}}
\newunicodechar{↦}{\ensuremath{\mapsto}}
\newunicodechar{∈}{\ensuremath{\in}}
\newunicodechar{∉}{\ensuremath{\notin}}
\newunicodechar{∧}{\ensuremath{\wedge}}
\newunicodechar{⇔}{\ensuremath{\Leftrightarrow}}
\newunicodechar{⟺}{\ensuremath{\Longleftrightarrow}}
\newunicodechar{⇒}{\ensuremath{\Rightarrow}}
\newunicodechar{⟹}{\ensuremath{\Longrightarrow}}
\newunicodechar{∘}{\ensuremath{\circ}}
\newunicodechar{∪}{\ensuremath{\cup}}
\newunicodechar{∩}{\ensuremath{\cap}}
\newunicodechar{≡}{\ensuremath{\equiv}}
\newunicodechar{⊂}{\ensuremath{\subset}}
\newunicodechar{⊥}{\ensuremath{\perp}}
\newunicodechar{∃}{\ensuremath{\exists}}
\newunicodechar{∀}{\ensuremath{\forall}}
\newunicodechar{∛}{\ensuremath{\sqrt[3]{\,}}}
\newunicodechar{∜}{\ensuremath{\sqrt[4]{\,}}}
\newunicodechar{⃗}{\ensuremath{{}^{\to}}}
\newunicodechar{ⁿ}{\ifmmode^{n}\else\textsuperscript{\ensuremath{n}}\fi}
\newunicodechar{ˣ}{\ifmmode^{x}\else\textsuperscript{\ensuremath{x}}\fi}
\newunicodechar{ᵇ}{\ifmmode^{b}\else\textsuperscript{\ensuremath{b}}\fi}
\newunicodechar{ʳ}{\ifmmode^{r}\else\textsuperscript{\ensuremath{r}}\fi}
\newunicodechar{ᵘ}{\ifmmode^{u}\else\textsuperscript{\ensuremath{u}}\fi}
\newunicodechar{ʸ}{\ifmmode^{y}\else\textsuperscript{\ensuremath{y}}\fi}
\newunicodechar{ᵃ}{\ifmmode^{a}\else\textsuperscript{\ensuremath{a}}\fi}
\newunicodechar{ᵉ}{\ifmmode^{e}\else\textsuperscript{\ensuremath{e}}\fi}
\newunicodechar{ᵏ}{\ifmmode^{k}\else\textsuperscript{\ensuremath{k}}\fi}
\newunicodechar{ᵖ}{\ifmmode^{p}\else\textsuperscript{\ensuremath{p}}\fi}
\newunicodechar{ᴺ}{\ifmmode^{N}\else\textsuperscript{\ensuremath{N}}\fi}
\newunicodechar{ᶜ}{\ifmmode^{c}\else\textsuperscript{\ensuremath{c}}\fi}
\newunicodechar{ʰ}{\ifmmode^{h}\else\textsuperscript{\ensuremath{h}}\fi}
\newunicodechar{ᵠ}{\ifmmode^{q}\else\textsuperscript{\ensuremath{q}}\fi}
\newunicodechar{ₙ}{\ifmmode_{n}\else\textsubscript{\ensuremath{n}}\fi}
\newunicodechar{ₐ}{\ifmmode_{a}\else\textsubscript{\ensuremath{a}}\fi}
\newunicodechar{ᵢ}{\ifmmode_{i}\else\textsubscript{\ensuremath{i}}\fi}
\newunicodechar{ᵣ}{\ifmmode_{r}\else\textsubscript{\ensuremath{r}}\fi}
\newunicodechar{ₖ}{\ifmmode_{k}\else\textsubscript{\ensuremath{k}}\fi}
\newunicodechar{ᵦ}{\ifmmode_{\beta}\else\textsubscript{\ensuremath{\beta}}\fi}
\newunicodechar{ₓ}{\ifmmode_{x}\else\textsubscript{\ensuremath{x}}\fi}
\newfontfamily\popdisplay{ArchivoBlack-Regular.ttf}[Path=../../../fonts/]
\newfontfamily\poplogo{SpaceGrotesk-Bold.ttf}[Path=../../../fonts/]
\newfontfamily\popsemi{PlusJakartaSans-SemiBold.ttf}[Path=../../../fonts/]
\newfontfamily\popxbold{PlusJakartaSans-ExtraBold.ttf}[Path=../../../fonts/]
\newfontfamily\popxboldls{PlusJakartaSans-ExtraBold.ttf}[Path=../../../fonts/,LetterSpace=3.0]

\setlist[enumerate]{leftmargin=1.7em,itemsep=5pt,topsep=4pt}
\setlist[itemize]{leftmargin=1.7em,itemsep=4pt,topsep=4pt,label={\color{poporange}\textbullet}}
\setlength{\parindent}{0pt}
\setlength{\parskip}{5pt}
\renewcommand{\arraystretch}{1.25}

% pgfplots : style Pop par défaut
\pgfplotsset{
  every axis/.append style={axis line style={popink,line width=1pt},tick style={popink},
    grid style={popline,line width=0.5pt},label style={font=\small},tick label style={font=\footnotesize}},
  popcurve/.style={poporange,line width=1.8pt,no markers,smooth},
}
% Même style disponible dans un \draw TikZ simple (hors pgfplots)
\tikzset{popcurve/.style={poporange,line width=1.8pt}}

% ============================================================
% Pill / squiggle
% ============================================================
\newcommand{\poppill}[3]{%
  \tikz[baseline=-0.55ex]\node[draw=popink,line width=1.2pt,rounded corners=2.6pt,%
    fill=#2,inner xsep=6.5pt,inner ysep=3pt]{%
    \popxboldls\fontsize{7}{7}\selectfont\color{#3}\MakeUppercase{#1}};%
}
\newcommand{\popsquiggle}[1]{%
  \tikz[baseline=-0.4ex]{
    \draw[#1,line width=2.6pt,line cap=round]
      plot [smooth] coordinates {(0,0.09) (0.34,0.01) (0.68,0.21) (1.02,0.01) (1.36,0.10)};
  }%
}
% Mot-clé surligné (marqueur orange sous le texte)
\newcommand{\cle}[1]{{\popxbold\color{popdark}#1}}

% ============================================================
% En-tête / pied
% ============================================================
\pagestyle{fancy}
\fancyhf{}
\renewcommand{\headrulewidth}{0pt}
\renewcommand{\footrulewidth}{0pt}
\lhead{\raisebox{-0.15ex}{\poplogo\fontsize{13}{13}\selectfont\color{popink}OnBuch\color{poporange}+}}
\rhead{\poppill{\DOCMATIERE\ \textbullet\ \DOCNIVEAU\ \textbullet\ Leçon \DOCLECON}{white}{popink}}
\lfoot{\popxbold\fontsize{7.6}{7.6}\selectfont\color{popmuted}\MakeUppercase{Cours signé OnBuch+}\ \textbullet\ {\popsemi\DOCTITRE}}
\rfoot{\tikz[baseline=-0.6ex]\node[rounded corners=8.5pt,fill=popink,minimum height=17pt,minimum width=17pt,inner xsep=5pt,inner sep=0]{\popxbold\fontsize{8}{8}\selectfont\color{popcream}\thepage\,/\,\pageref*{LastPage}};}

% ============================================================
% Titres
% ============================================================
\newcommand{\chaptitre}[2]{%
  \begin{tcolorbox}[enhanced,colback=poporange,colframe=popink,boxrule=2.6pt,arc=11pt,
      drop shadow=popink,left=15pt,right=15pt,top=12pt,bottom=11pt,before skip=3pt,after skip=18pt]
    \poppill{#2}{popink}{popcream}\par\vspace{9pt}
    {\raggedright\hyphenpenalty=10000\exhyphenpenalty=10000\popdisplay\fontsize{22}{24}\selectfont\color{popcream}\MakeUppercase{#1}\par}
  \end{tcolorbox}
}
% Section numérotée : pastille orange avec le numéro + titre Archivo
\newcounter{popsec}
\newcommand{\coursec}[1]{%
  \stepcounter{popsec}\setcounter{popsub}{0}%
  \par\vspace{16pt}\noindent
  \begin{minipage}{\linewidth}
  \tikz[baseline=-0.7ex]\node[draw=popink,line width=1.6pt,fill=poporange,rounded corners=4pt,minimum size=22pt,inner sep=0]{\popdisplay\fontsize{12}{12}\selectfont\color{popcream}\thepopsec};%
  \hspace{8pt}{\popdisplay\fontsize{15}{17}\selectfont\color{popink}\MakeUppercase{#1}}\par
  \vspace{4pt}\noindent{\color{poporange}\rule{36pt}{3.6pt}}
  \end{minipage}\par\nopagebreak\vspace{8pt}
}
\newcounter{popsub}
\newcommand{\courssub}[1]{%
  \stepcounter{popsub}%
  \par\vspace{9pt}\noindent{\popxbold\fontsize{11.5}{13}\selectfont\color{popink}{\color{poporange}\thepopsec.\thepopsub}\hspace{6pt}#1}\par\nopagebreak\vspace{4pt}
}

% ============================================================
% Boîtes pédagogiques — même recette : fond blanc, bordure épaisse colorée,
% ombre dure encre, badge plein. Titre optionnel [#1].
% ============================================================
\tcbset{popbase/.style={breakable,enhanced,colback=white,boxrule=2.3pt,arc=9pt,
  shadow={3pt}{-3pt}{0pt}{popink},left=14pt,right=14pt,top=11pt,bottom=11pt,before skip=11pt,after skip=15pt,
  fontupper=\popsemi\fontsize{10.6}{14.5}\selectfont\color{popink},
  fontlower=\popsemi\fontsize{10.6}{14.5}\selectfont\color{popink}}}
% Coche dessinée (le glyphe ✓ n'existe pas dans les polices du document)
\newcommand{\popcheck}[1]{\tikz[baseline=-0.5ex]\draw[#1,line width=1.6pt,line cap=round,line join=round](-0.13,0.0)--(-0.04,-0.09)--(0.13,0.1);}
\newcommand{\popboxhead}[3]{\poppill{#1}{#2}{white}\ifx\relax#3\relax\else\hspace{6pt}{\popxbold\fontsize{10.6}{12}\selectfont\color{popink}#3}\fi\par\vspace{7pt}}

\newtcolorbox{aretenir}[1][]{popbase,colframe=popgreen,before upper={\popboxhead{À retenir}{popgreen}{#1}}}
\newtcolorbox{exemplebox}[1][]{popbase,colframe=poporange,before upper={\popboxhead{Exemple}{poporange}{#1}}}
\newtcolorbox{definition}[1][]{popbase,colframe=popblue,before upper={\popboxhead{Définition}{popblue}{#1}}}
\newtcolorbox{propriete}[1][]{popbase,colframe=poppurple,before upper={\popboxhead{Propriété}{poppurple}{#1}}}
\newtcolorbox{methode}[1][]{popbase,colframe=popgold,before upper={\popboxhead{Méthode}{popgold}{#1}}}
\newtcolorbox{attention}[1][]{popbase,colframe=poppink,before upper={\popboxhead{Attention}{poppink}{#1}}}
\newtcolorbox{experience}[1][]{popbase,colframe=popblue,colback=popblueL,before upper={\popboxhead{Expérience}{popblue}{#1}}}
% « Le savais-tu ? » : carte violette pleine (contexte, culture, Cameroun)
\newtcolorbox{savaistu}[1][]{popbase,colback=poppurple,colframe=popink,
  fontupper=\popsemi\fontsize{10.4}{14.2}\selectfont\color{white},
  before upper={\poppill{Le savais-tu\,?}{white}{poppurple}\ifx\relax#1\relax\else\hspace{6pt}{\popxbold\color{white}#1}\fi\par\vspace{7pt}}}
% Exercice résolu : énoncé en haut, \tcblower puis la solution sur fond crème
\newcounter{popexo}\newcounter{popexr}
\newtcolorbox{exoresolu}[1][]{popbase,colframe=poporange,colbacklower=popcream,
  segmentation style={popink,line width=1.2pt,dashed},
  before upper={\stepcounter{popexr}\popboxhead{Exercice résolu \thepopexr}{poporange}{#1}},
  before lower={\poppill{Solution}{popink}{popcream}\par\vspace{6pt}}}
% Exercice (non résolu en place) — la correction est dans la partie Corrigés
\newtcolorbox{exercice}[1][]{popbase,colframe=popink,
  before upper={\stepcounter{popexo}\popboxhead{Exercice \thepopexo}{popink}{#1}}}
% Corrigé
\newtcolorbox{corrige}[1][]{popbase,colframe=popgreen,colback=popgreenL,
  before upper={\popboxhead{Corrigé}{popgreen}{#1}}}
% Objectifs / prérequis / bilan
\newtcolorbox{objectifs}{popbase,colframe=popink,colback=poporangeL,
  before upper={\popboxhead{Objectifs de la leçon}{popink}{}}}
\newtcolorbox{prerequis}{popbase,colframe=popink,colback=popblueL,
  before upper={\popboxhead{Prérequis}{popblue}{}}}
\newtcolorbox{bilan}{popbase,colframe=popink,colback=white,boxrule=2.8pt,
  before upper={\popboxhead{Fiche bilan}{poporange}{L'essentiel à connaître pour l'examen}}}
% Figure encadrée avec légende
\newtcolorbox{popfigure}[1][]{enhanced,breakable=false,colback=white,colframe=popline,boxrule=1.4pt,arc=8pt,
  left=10pt,right=10pt,top=10pt,bottom=8pt,before skip=10pt,after skip=12pt,halign=center,
  lower separated=false,halign lower=center,
  fontlower=\popsemi\fontsize{9}{11.5}\selectfont\color{popmuted},
  #1}
\newcommand{\legende}[1]{\par\vspace{6pt}{\popsemi\fontsize{9}{11.5}\selectfont\color{popmuted}#1}}

% ============================================================
% Signature OnBuch+
% ============================================================
\newcommand{\poplogoinline}[1]{{\poplogo\fontsize{#1}{#1}\selectfont\color{popink}OnBuch\color{poporange}+}}
\newcommand{\signatureonbuch}{%
  \begin{tcolorbox}[enhanced,colback=popink,colframe=popink,boxrule=2.2pt,arc=10pt,
      drop shadow=poporange,left=14pt,right=14pt,top=10pt,bottom=10pt,before skip=0pt,after skip=0pt]
    \tikz[baseline=-0.7ex]\node[circle,fill=popgreen,draw=popcream,line width=1.3pt,minimum size=22pt,inner sep=0]{\popcheck{white}};%
    \hspace{9pt}\begin{minipage}[c]{0.62\linewidth}
      {\popxbold\fontsize{12}{13}\selectfont\color{popcream}Signé\ \textbullet\ {\poplogo OnBuch\color{poporange}+}}\par\vspace{2pt}
      {\raggedright\popsemi\fontsize{8}{10.5}\selectfont\color{popline}\MakeUppercase{Équipe pédagogique OnBuch+}\\\MakeUppercase{Conforme au programme officiel MINESEC}\par}
    \end{minipage}\hfill
    {\poplogo\fontsize{22}{22}\selectfont\color{popcream}OnBuch\color{poporange}+}
  \end{tcolorbox}%
}

% ============================================================
% Page de garde
% #1 titre · #2 matière · #3 niveau · #4 module · #5 n° leçon · #6 durée ·
% #7 description / objectifs (texte ou itemize)
% ============================================================
\newcommand{\pagegarde}[7]{%
  \thispagestyle{empty}
  \newgeometry{a4paper,margin=1.7cm,top=1.6cm,bottom=1.4cm}%
  \begin{tikzpicture}[remember picture,overlay]
    \fill[poporange!18] ([xshift=-3.2cm,yshift=-3.6cm]current page.north east) circle (4.6cm);
    \fill[poppurple!14] ([xshift=2.4cm,yshift=7.6cm]current page.south west) circle (3.8cm);
  \end{tikzpicture}%
  {\poplogo\fontsize{30}{30}\selectfont\color{popink}OnBuch\color{poporange}+}\hfill
  \raisebox{0.6ex}{\poppill{Cours signé \textbullet\ Premium}{popink}{popcream}}\par
  \vspace{1pt}\popsquiggle{poppurple}\par
  \vspace{8pt}
  \begin{tcolorbox}[enhanced,colback=poporange,colframe=popink,boxrule=2.8pt,arc=13pt,
      drop shadow=popink,left=18pt,right=18pt,top=12pt,bottom=13pt,after skip=10pt]
    \poppill{#2 \textbullet\ #3}{popink}{popcream}\hspace{5pt}\poppill{#4}{white}{popink}\par\vspace{8pt}
    {\popdisplay\fontsize{14}{14}\selectfont\color{popink}LEÇON #5}\par\vspace{4pt}
    {\raggedright\hyphenpenalty=10000\exhyphenpenalty=10000\popdisplay\fontsize{25}{27}\selectfont\color{popcream}\MakeUppercase{#1}\par}
    \vspace{8pt}
    {\popxbold\fontsize{9.5}{9.5}\selectfont\color{popink}\MakeUppercase{Durée conseillée : #6}}
  \end{tcolorbox}
  \begin{tcolorbox}[enhanced,colback=white,colframe=popink,boxrule=2.2pt,arc=10pt,
      drop shadow=popink,left=16pt,right=16pt,top=10pt,bottom=10pt,after skip=8pt]
    \poppill{Dans cette leçon}{poppurple}{white}\par\vspace{5pt}
    \popsemi\fontsize{9.3}{12.6}\selectfont\color{popink}#7
  \end{tcolorbox}
  \noindent\begin{minipage}[t]{0.487\linewidth}
    \begin{tcolorbox}[enhanced,colback=popgreen,colframe=popink,boxrule=2.2pt,arc=10pt,
        drop shadow=popink,left=11pt,right=11pt,top=10pt,bottom=10pt,height=3.1cm,valign=center]
      \tcbox[colback=white,colframe=popink,boxrule=1.3pt,arc=5pt,boxsep=0pt,nobeforeafter,
        left=3pt,right=3pt,top=3pt,bottom=3pt,valign=center]{\qrcode[height=1.9cm]{https://play.google.com/store/apps/details?id=com.luvvix.onbuch&hl=fr}}%
      \hspace{8pt}\begin{minipage}[c]{0.5\linewidth}
        {\popdisplay\fontsize{11}{12.5}\selectfont\color{white}\MakeUppercase{Télécharge OnBuch}}\par\vspace{4pt}
        {\raggedright\popsemi\fontsize{8.4}{11}\selectfont\color{white}Gratuit \textbullet\ Google Play\\Tuteur IA, annales, concours, résultats\par}
      \end{minipage}
    \end{tcolorbox}
  \end{minipage}\hfill
  \begin{minipage}[t]{0.487\linewidth}
    \begin{tcolorbox}[enhanced,colback=poppurple,colframe=popink,boxrule=2.2pt,arc=10pt,
        drop shadow=popink,left=11pt,right=11pt,top=10pt,bottom=10pt,height=3.1cm,valign=center]
      \tikz[baseline=-0.7ex]\node[circle,fill=white,draw=popink,line width=1.3pt,minimum size=1.6cm,inner sep=0]{\popdisplay\fontsize{24}{24}\selectfont\color{poppurple}+};%
      \hspace{8pt}\begin{minipage}[c]{0.6\linewidth}
        {\popdisplay\fontsize{11}{12.5}\selectfont\color{white}\MakeUppercase{Passe à OnBuch+}}\par\vspace{4pt}
        {\raggedright\popsemi\fontsize{8.4}{11}\selectfont\color{white}Tous les cours signés, Tuteur IA illimité, fiches PDF — onglet OnBuch+ de l'app\par}
      \end{minipage}
    \end{tcolorbox}
  \end{minipage}
  \vspace{8pt}
  \popcontenu
  \vfill
  \signatureonbuch
  \restoregeometry
  \newpage
}

% Rangée « Ce cours contient » (page de garde)
\newcommand{\poptile}[3]{% #1 couleur #2 gros texte #3 légende
  \begin{tcolorbox}[enhanced,colback=#1,colframe=popink,boxrule=2pt,arc=9pt,shadow={2.5pt}{-2.5pt}{0pt}{popink},
      left=6pt,right=6pt,top=9pt,bottom=9pt,halign=center,valign=center,height=1.75cm,width=\linewidth,nobeforeafter]
    {\popdisplay\fontsize{12.5}{13}\selectfont\color{white}\MakeUppercase{#2}}\par\vspace{3pt}
    {\centering\popsemi\fontsize{7.8}{9.5}\selectfont\color{white}#3\par}
  \end{tcolorbox}}
\newcommand{\popcontenu}{%
  \par\noindent{\popxboldls\fontsize{7.5}{7.5}\selectfont\color{popmuted}CE COURS SIGNÉ CONTIENT}\par\vspace{7pt}
  \noindent\begin{minipage}[t]{0.235\linewidth}\poptile{popblue}{Cours}{complet, illustré, pas à pas}\end{minipage}\hfill
  \begin{minipage}[t]{0.235\linewidth}\poptile{poporange}{Méthodes}{et exercices résolus}\end{minipage}\hfill
  \begin{minipage}[t]{0.235\linewidth}\poptile{poppink}{Exercices}{type examen + corrigés}\end{minipage}\hfill
  \begin{minipage}[t]{0.235\linewidth}\poptile{popgreen}{Fiche bilan}{l'essentiel à retenir}\end{minipage}\par}

% Sommaire
\newtcolorbox{sommaire}{enhanced,colback=white,colframe=popink,boxrule=2.2pt,arc=10pt,
  drop shadow=popink,left=18pt,right=18pt,top=13pt,bottom=13pt,before skip=6pt,after skip=20pt,
  before upper={\poppill{Sommaire}{poppurple}{white}\par\vspace{9pt}\popsemi\fontsize{10.6}{14.5}\selectfont\color{popink}}}

% Signature de fin de cours — poussée tout en bas de la dernière page
\newcommand{\findecours}{%
  \par\vfill\nopagebreak
  \begin{center}\popsquiggle{poporange}\end{center}\vspace{4pt}
  \signatureonbuch
}
\AtBeginDocument{\def\llbracket{\mathopen{[\mkern-3.5mu[}}\def\rrbracket{\mathclose{]\mkern-3.5mu]}}} % intervalles d'entiers (glyphe ⟦ absent de la police)
% Glyphes absents de Fira Math : redéfinis à partir de glyphes disponibles
\AtBeginDocument{%
  \def\setminus{\mathbin{\backslash}}\let\smallsetminus\setminus
  \def\vdots{\mathord{\vbox{\baselineskip3.2pt\lineskiplimit0pt\kern4pt\hbox{.}\hbox{.}\hbox{.}}}}%
  \def\ddots{\mathinner{\mkern1mu\raise6.5pt\hbox{.}\mkern2mu\raise3.6pt\hbox{.}\mkern2mu\raise0.7pt\hbox{.}\mkern1mu}}%
  \def\triangle{\mathord{\scalebox{0.9}{$\symup{\Delta}$}}}%
  \def\nabla{\mathord{\raisebox{\depth}{\scalebox{1}[-1]{$\symup{\Delta}$}}}}%
  \def\checkmark{\popcheck{popdark}}%
  \def\star{\mathbin{\ast}}\def\bigstar{\mathord{\ast}}%
  \def\diamond{\mathbin{\scalebox{0.6}{$\Diamond$}}}%
  \def\bigcup{\mathop{\mathchoice{\raisebox{-0.3em}{\scalebox{1.7}{$\cup$}}}{\scalebox{1.25}{$\cup$}}{\cup}{\cup}}\displaylimits}%
  \def\bigcap{\mathop{\mathchoice{\raisebox{-0.3em}{\scalebox{1.7}{$\cap$}}}{\scalebox{1.25}{$\cap$}}{\cap}{\cap}}\displaylimits}%
  \def\lesseqqgtr{\mathrel{\lessgtr}}\def\bigoplus{\mathop{\oplus}\displaylimits}%
}
\providecommand{\ssi}{si et seulement si} % abréviation fréquente
\AtBeginDocument{\providecommand{\wideparen}{\overparen}} % arc de cercle

\begin{document}

\pagegarde{Primitives d’une fonction continue sur un intervalle}{Mathématiques}{Tle Tech}{Mathématiques – classe de Terminale technique}{N10}{1 séance (fiche de progression harmonisée)}{Cette leçon introduit la notion de primitive d'une fonction continue sur un intervalle, montre son existence et donne les outils pour la déterminer. Elle prépare les élèves à résoudre des problèmes concrets de calcul de quantités accumulées.
\vspace{4pt}
\begin{multicols}{2}\small
\begin{itemize}
\item À la fin de cette leçon, je sais définir une primitive d'une fonction continue sur un intervalle.
\item Je sais énoncer et expliquer le théorème d'existence des primitives pour les fonctions continues.
\item Je sais utiliser les propriétés de linéarité et de translation par une constante des primitives.
\item Je sais déterminer une primitive à partir du tableau des primitives usuelles.
\item Je sais calculer une primitive en appliquant la méthode de l'intégration par parties ou par substitution quand c'est nécessaire.
\item Je sais modéliser une situation réelle (distance, aire, coût) à l'aide d'une primitive et interpréter le résultat.
\end{itemize}\end{multicols}}

\begin{sommaire}
\begin{multicols}{2}
\begin{enumerate}
\item Rappels : dérivée et notion de primitive
\item Théorème d'existence des primitives pour les fonctions continues
\item Propriétés algébriques des primitives
\item Détermination de primitives : exercices types et astuces
\item Applications concrètes : modélisation de situations réelles
\item Synthèse, entraînement et évaluation formative
\item Activité d'intégration
\item Méthodes et exercices résolus
\item Exercices
\item Corrigés des exercices
\item Fiche bilan
\end{enumerate}
\end{multicols}
\end{sommaire}

\begin{objectifs}
\begin{itemize}
\item À la fin de cette leçon, je sais définir une primitive d'une fonction continue sur un intervalle.
\item Je sais énoncer et expliquer le théorème d'existence des primitives pour les fonctions continues.
\item Je sais utiliser les propriétés de linéarité et de translation par une constante des primitives.
\item Je sais déterminer une primitive à partir du tableau des primitives usuelles.
\item Je sais calculer une primitive en appliquant la méthode de l'intégration par parties ou par substitution quand c'est nécessaire.
\item Je sais modéliser une situation réelle (distance, aire, coût) à l'aide d'une primitive et interpréter le résultat.
\item Je sais rédiger une démonstration guidée du théorème d'existence et justifier chaque étape.
\item Je sais identifier et éviter les erreurs fréquentes (oubli de la constante, confusion dérivée/primitive).
\end{itemize}
\end{objectifs}

\begin{prerequis}
\begin{itemize}
\item Définition de la dérivée d'une fonction en un point et sur un intervalle (classe de 1ère).
\item Règles de dérivation usuelles (somme, produit, quotient, chaîne).
\item Notion de fonction continue sur un intervalle (classe de 1ère).
\item Calcul de limites simples et notion de continuité à une frontière.
\item Manipulation algébrique de base (factorisation, développement).
\end{itemize}
\end{prerequis}

\newpage

\coursec{Leçon 10 : Primitives d'une fonction continue sur un intervalle}

\coursec{Rappels : dérivée et notion de primitive}

\textbf{Situation d'accroche.} À Douala, un chauffeur de moto-taxi (un \emph{benskin}) note sa vitesse en \SI{}{km/h} toutes les 10 minutes pendant une course entre Akwa et Bonabéri. À la fin de la journée, il veut connaître la distance totale parcourue, mais il n'a mesuré aucun trajet : il ne possède que sa \emph{fonction vitesse}. Comment retrouver la distance à partir de la vitesse ?

Ce problème est l'inverse de celui que nous connaissons déjà : si l'on connaît la distance parcourue en fonction du temps, on obtient la vitesse en \emph{dérivant}. Ici, on connaît la vitesse et l'on cherche la distance : il faut donc faire le \emph{chemin inverse de la dérivation}. Cette opération inverse s'appelle la \cle{recherche de primitives}. C'est l'objet de toute cette leçon, et cette première section pose les fondations : rappeler ce qu'est une dérivée, puis définir rigoureusement ce qu'est une primitive.

\textbf{Problématique.} Étant donnée une fonction $f$, comment trouver une fonction $F$ dont la dérivée est exactement $f$ ?

\courssub{Rappel : la dérivée d'une fonction en un point}

\begin{definition}[Nombre dérivé]
Soit $f$ une fonction définie sur un intervalle $I$ et $a$ un réel de $I$. On dit que $f$ est \cle{dérivable en} $a$ si le taux d'accroissement
\[ \frac{f(a+h)-f(a)}{h} \]
admet une limite finie lorsque $h$ tend vers $0$. Cette limite est alors appelée \cle{nombre dérivé} de $f$ en $a$ et notée $f'(a)$ :
\[ f'(a)=\lim_{h\to 0}\frac{f(a+h)-f(a)}{h}. \]
\end{definition}

\begin{aretenir}[Interprétations du nombre dérivé]
\begin{itemize}
\item \textbf{Géométrique :} $f'(a)$ est le \emph{coefficient directeur de la tangente} à la courbe de $f$ au point d'abscisse $a$. L'équation de cette tangente est $y=f'(a)(x-a)+f(a)$.
\item \textbf{Physique :} si $x(t)$ est la position d'un mobile à l'instant $t$, alors $x'(t)$ est sa \emph{vitesse instantanée}. C'est exactement le lien entre distance et vitesse de notre moto-taxi.
\end{itemize}
\end{aretenir}

Si $f$ est dérivable en tout point de $I$, on définit la \cle{fonction dérivée} $f'$, qui à tout $x$ de $I$ associe $f'(x)$.

\courssub{Tableau des dérivées usuelles}

La recherche de primitives consistant à « remonter » la dérivation, il est indispensable de connaître parfaitement le tableau des dérivées usuelles : c'est lui que nous lirons \emph{à l'envers} dans toute la leçon.

\begin{center}
\begin{tabularx}{\linewidth}{|l|l|X|}
\hline
\rowcolor{popblueL}
Fonction $f$ & Dérivée $f'$ & Ensemble de dérivabilité \\
\hline
$k$ (constante) & $0$ & $\mathbb{R}$ \\
\hline
$x$ & $1$ & $\mathbb{R}$ \\
\hline
$x^n$ ($n\in\mathbb{N}^*$) & $n x^{n-1}$ & $\mathbb{R}$ \\
\hline
$x^{\alpha}$ ($\alpha\in\mathbb{R}$) & $\alpha x^{\alpha-1}$ & Dépend de $\alpha$ ($]0,+\infty[$ si $\alpha\notin\mathbb{Z}$) \\
\hline
$\dfrac{1}{x}$ & $-\dfrac{1}{x^2}$ & $\mathbb{R}^*$ \\
\hline
$\sqrt{x}$ & $\dfrac{1}{2\sqrt{x}}$ & $]0\,;+\infty[$ \\
\hline
$e^{x}$ & $e^{x}$ & $\mathbb{R}$ \\
\hline
$\ln x$ & $\dfrac{1}{x}$ & $]0\,;+\infty[$ \\
\hline
$\cos x$ & $-\sin x$ & $\mathbb{R}$ \\
\hline
$\sin x$ & $\cos x$ & $\mathbb{R}$ \\
\hline
\end{tabularx}
\end{center}

On rappelle aussi les règles opératoires de la dérivation :
\begin{itemize}
\item $(u+v)' = u' + v'$ \quad (somme)
\item $(ku)' = k u'$ \quad ($k\in\mathbb{R}$, multiple)
\item $(uv)' = u'v + uv'$ \quad (produit)
\item $\left(\dfrac{u}{v}\right)' = \dfrac{u'v - uv'}{v^2}$ \quad (quotient, $v\neq 0$)
\item \textbf{(Règle de la chaîne)} $(f \circ g)' = g' \cdot (f' \circ g)$ \quad (composition)
\end{itemize}

\begin{attention}[La dérivation détruit les constantes]
La dérivée d'une constante est nulle. Ainsi les fonctions $x^2$, $x^2+5$ et $x^2-13$ ont \emph{toutes} la même dérivée : $2x$. Conséquence capitale : quand on « remonte » de $2x$ vers une fonction dont elle est la dérivée, on ne peut pas savoir quelle constante avait été ajoutée. C'est toute la subtilité des primitives.
\end{attention}

\courssub{Introduction intuitive : remonter la dérivation}

Revenons au chauffeur de moto-taxi. Supposons, pour simplifier, que sa vitesse (en unité adaptée) soit donnée par $v(t)=2t$. On sait que la vitesse est la dérivée de la distance : $v(t)=d'(t)$. Trouver la distance revient donc à chercher une fonction $d$ telle que $d'(t)=2t$. Or on sait dériver : la dérivée de $t^2$ est $2t$. Donc $d(t)=t^2$ convient. On dit que $t^2$ est une \emph{primitive} de $2t$.

L'idée est donc très simple : \textbf{une primitive de $f$ est une fonction $F$ dont la dérivée redonne $f$}. La primitivation est l'opération inverse de la dérivation, comme la division est l'opération inverse de la multiplication.

\begin{definition}[Primitive d'une fonction sur un intervalle]
Soit $f$ une fonction définie sur un intervalle $I$. On appelle \cle{primitive} de $f$ sur $I$ toute fonction $F$ définie et dérivable sur $I$ telle que :
\[ \text{pour tout } x\in I,\quad F'(x)=f(x). \]
\end{definition}

\begin{exemplebox}[Un premier exemple fondamental]
Vérifions que $F(x)=x^2$ est une primitive de $f(x)=2x$ sur $\mathbb{R}$.

$F$ est dérivable sur $\mathbb{R}$ et, pour tout réel $x$ :
\[ F'(x)=2x=f(x). \]
Donc $F$ est bien une primitive de $f$ sur $\mathbb{R}$.

Mais remarquons que $G(x)=x^2+7$ vérifie aussi $G'(x)=2x=f(x)$ : $G$ est \emph{aussi} une primitive de $f$. Une fonction qui admet une primitive en admet donc une infinité.
\end{exemplebox}

La figure suivante illustre le lien entre $f(x)=2x$ et sa primitive $F(x)=x^2$ : en chaque point de la parabole, le coefficient directeur de la tangente est donné par la valeur de $f$.

\begin{popfigure}
\begin{center}
\begin{tikzpicture}
\begin{axis}[
  width=0.85\linewidth,
  axis lines=middle,
  xlabel={$x$}, ylabel={$y$},
  xmin=-2.6, xmax=2.6, ymin=-1.5, ymax=5.2,
  xtick={-2,-1,1,2}, ytick={1,2,3,4,5},
  legend style={at={(0.02,0.98)},anchor=north west,font=\small},
  grid=both, grid style={popline!40},
]
\addplot[popcurve,domain=-2.2:2.2,samples=200]{x^2};
\addlegendentry{$F(x)=x^2$}
\addplot[poporange,thick,domain=-1.2:2.4,samples=2]{2*x};
\addlegendentry{$f(x)=2x$}
% tangente en x=1 : y = 2(x-1)+1 = 2x-1
\addplot[poppurple,thick,dashed,domain=0:2.2,samples=2]{2*x-1};
\addlegendentry{tangente à $F$ en $x=1$}
% tangente en x=-1 : y = -2(x+1)+1 = -2x-1
\addplot[popgreen,thick,dashed,domain=-2.2:0,samples=2]{-2*x-1};
\addlegendentry{tangente à $F$ en $x=-1$}
\addplot[only marks,mark=*,popdark] coordinates {(1,1) (-1,1)};
\end{axis}
\end{tikzpicture}
\end{center}
\legende{La courbe de $F(x)=x^2$ et la droite $f(x)=2x$. Au point d'abscisse $1$, la tangente à la parabole a pour coefficient directeur $2=f(1)$ ; au point d'abscisse $-1$, ce coefficient vaut $-2=f(-1)$. La fonction $f$ donne la \emph{pente} de la tangente à la courbe de $F$ en tout point.}
\end{popfigure}

\courssub{Une infinité de primitives : le rôle de la constante}

\begin{propriete}[Primitives et constante additive]
Soit $f$ une fonction admettant une primitive $F$ sur un intervalle $I$. Alors, pour toute constante réelle $C$, la fonction $F+C$ (définie par $x\mapsto F(x)+C$) est aussi une primitive de $f$ sur $I$.
\end{propriete}

\begin{experience}[Démonstration guidée]
Montrons cette propriété pas à pas.
\begin{enumerate}
\item Soit $G$ la fonction définie sur $I$ par $G(x)=F(x)+C$, où $C$ est un réel fixé.
\item $F$ est dérivable sur $I$ (c'est une primitive de $f$) et la fonction constante $x\mapsto C$ est dérivable sur $I$, de dérivée nulle.
\item La somme de deux fonctions dérivables est dérivable, donc $G$ est dérivable sur $I$ et :
\[ G'(x)=F'(x)+0=f(x). \]
\item Conclusion : $G'=f$ sur $I$, donc $G=F+C$ est bien une primitive de $f$ sur $I$. \hfill$\square$
\end{enumerate}
\end{experience}

\begin{aretenir}[À retenir absolument]
Si $F$ est une primitive de $f$ sur $I$, alors \textbf{toutes} les fonctions $F+C$ (où $C$ est une constante réelle quelconque) sont des primitives de $f$. On écrit souvent : « les primitives de $f$ sont les fonctions $x\mapsto F(x)+C$, $C\in\mathbb{R}$ ». On verra dans la section « Propriétés algébriques des primitives » que, sur un intervalle, il n'y en a pas d'autres : deux primitives d'une même fonction diffèrent toujours d'une constante.
\end{aretenir}

\begin{attention}[Erreur fréquente : oublier la constante]
L'erreur la plus classique au Probatoire et au Baccalauréat technique est d'écrire « la primitive de $2x$ est $x^2$ » en oubliant la constante. Rédaction correcte : « les primitives de $f(x)=2x$ sur $\mathbb{R}$ sont les fonctions $F(x)=x^2+C$, où $C\in\mathbb{R}$ ». L'oubli de la constante coûte des points, surtout lorsque l'énoncé demande ensuite la primitive vérifiant une condition donnée (par exemple $F(0)=3$), car c'est précisément la constante qu'il faut alors ajuster.
\end{attention}

\courssub{Exemple chiffré détaillé}

\begin{exoresolu}[Primitive de $f(x)=3x^2$ sur $\mathbb{R}$]
Déterminer les primitives de la fonction $f$ définie sur $\mathbb{R}$ par $f(x)=3x^2$, puis celle qui s'annule en $x=1$.
\tcblower
\textbf{Étape 1 : lecture « à l'envers » du tableau des dérivées.}

On sait que $(x^3)'=3x^2$. La fonction $F_0(x)=x^3$ convient donc.

\textbf{Étape 2 : vérification (réflexe obligatoire).}

$F_0$ est dérivable sur $\mathbb{R}$ et $F_0'(x)=3x^2=f(x)$ pour tout réel $x$. $F_0$ est bien une primitive de $f$.

\textbf{Étape 3 : ensemble de toutes les primitives.}

D'après la propriété de la constante additive, les primitives de $f$ sur $\mathbb{R}$ sont les fonctions :
\[ F(x)=x^3+C,\quad C\in\mathbb{R}. \]

\textbf{Étape 4 : primitive qui s'annule en $1$.}

On cherche $C$ tel que $F(1)=0$ :
\[ F(1)=1^3+C=0 \iff C=-1. \]
La primitive cherchée est donc $F(x)=x^3-1$.
\end{exoresolu}

\begin{methode}[Méthode générale pour trouver une primitive simple]
\begin{enumerate}
\item \textbf{Identifier} la forme de $f$ (puissance, exponentielle, logarithme, trigonométrique, composition $u(ax+b)$).
\item \textbf{Lire le tableau des dérivées à l'envers} : quelle fonction usuelle a pour dérivée $f$ (ou un multiple de $f$) ?
\item \textbf{Ajuster le coefficient} si nécessaire : par exemple pour $f(x)=6x$, on pense à $x^2$ dont la dérivée est $2x$, donc on multiplie par $3$ : $F(x)=3x^2$. Pour une forme $f(ax+b)$, on divise par $a$ : primitive de $\cos(2x+1)$ est $\frac{1}{2}\sin(2x+1)$.
\item \textbf{Vérifier en dérivant} : on doit retrouver exactement $f$.
\item \textbf{Ajouter la constante} $+C$ si l'on demande toutes les primitives.
\end{enumerate}
\end{methode}

\courssub{Primitive et intégrale indéfinie : une distinction importante}

Dans la pratique (et dans certains ouvrages), on rencontre la notation
\[ \int f(x)\,\mathrm{d}x \]
appelée \cle{intégrale indéfinie} de $f$. Elle désigne \emph{l'ensemble de toutes les primitives} de $f$ : on écrit par exemple
\[ \int 2x\,\mathrm{d}x = x^2 + C,\quad C\in\mathbb{R}. \]

\begin{attention}[Ne pas confondre]
\begin{itemize}
\item Une \textbf{primitive} de $f$ est \emph{une} fonction $F$ précise telle que $F'=f$ (par exemple $x^2$).
\item L'\textbf{intégrale indéfinie} désigne la \emph{famille de toutes} les primitives, d'où la présence systématique du $+C$.
\item Il ne faut pas non plus confondre avec l'\emph{intégrale définie} $\displaystyle\int_a^b f(x)\,\mathrm{d}x$ (un nombre, lié à une aire), qui sera étudiée dans une leçon ultérieure et qui répondra pleinement au problème de la distance parcourue par notre moto-taxi.
\end{itemize}
\end{attention}

\begin{propriete}[Théorème d'existence des primitives (annoncé)]
\emph{Toute fonction continue sur un intervalle $I$ admet au moins une primitive sur $I$.}
Ce théorème fondamental garantit que pour toutes les fonctions continues que nous rencontrerons (polynômes, rationnelles sur leur ensemble de définition, exponentielle, logarithme, trigonométriques), la recherche de primitives est toujours possible.
\end{propriete}

\begin{savaistu}[Pourquoi « primitive » ?]
Le mot vient du latin \emph{primitivus}, « qui vient en premier » : la primitive est la fonction qui existait « avant » la dérivation. C'est Newton et Leibniz, au XVII\textsuperscript{e} siècle, qui ont mis en évidence ce lien entre dérivation et calcul d'aires, connu sous le nom de \emph{théorème fondamental de l'analyse}. Au Cameroun, ce même principe est utilisé chaque jour : le compteur d'électricité d'ENEO « remonte » de la puissance consommée (instantanée) à l'énergie totale consommée, exactement comme notre chauffeur remonte de la vitesse à la distance.
\end{savaistu}

\begin{exercice}[Pour s'entraîner]
\begin{enumerate}
\item Vérifier que $F(x)=e^x+4$ est une primitive de $f(x)=e^x$ sur $\mathbb{R}$.
\item Déterminer une primitive de $f(x)=5x^4$ sur $\mathbb{R}$, puis toutes ses primitives.
\item Déterminer la primitive de $f(x)=\cos x$ sur $\mathbb{R}$ qui prend la valeur $2$ en $x=0$.
\item Déterminer les primitives de $f(x)=\frac{1}{2\sqrt{x}}$ sur $]0\,;+\infty[$.
\item Trouver la primitive de $f(x)=3\cos(2x)$ sur $\mathbb{R}$ qui s'annule en $0$.
\end{enumerate}
\end{exercice}

\begin{corrige}[Exercices d'entraînement]
\begin{enumerate}
\item $F$ est dérivable sur $\mathbb{R}$ et $F'(x)=e^x+0=e^x=f(x)$ : $F$ est bien une primitive de $f$.
\item $(x^5)'=5x^4$, donc $F_0(x)=x^5$ convient. Toutes les primitives sont $F(x)=x^5+C$, $C\in\mathbb{R}$.
\item Une primitive de $\cos x$ est $\sin x$ car $(\sin x)'=\cos x$. Les primitives sont $F(x)=\sin x + C$. La condition $F(0)=2$ donne $\sin 0 + C = 2$, soit $C=2$. D'où $F(x)=\sin x + 2$.
\item On reconnaît la dérivée de $\sqrt{x}$ : $(\sqrt{x})' = \frac{1}{2\sqrt{x}}$. Les primitives sont $F(x)=\sqrt{x}+C$, $C\in\mathbb{R}$ (sur $]0,+\infty[$).
\item Une primitive de $\cos(2x)$ est $\frac{1}{2}\sin(2x)$ (règle de la chaîne : on divise par le coefficient de $x$, ici $2$). Donc une primitive de $3\cos(2x)$ est $F_0(x)=\frac{3}{2}\sin(2x)$. Les primitives sont $F(x)=\frac{3}{2}\sin(2x)+C$. La condition $F(0)=0$ donne $\frac{3}{2}\sin(0)+C=0 \iff C=0$. D'où $F(x)=\frac{3}{2}\sin(2x)$.
\end{enumerate}
\end{corrige}

\coursec{Théorème d'existence des primitives pour les fonctions continues}

Dans la section précédente, nous avons rappelé la notion de \cle{dérivée} et défini ce qu'est une \cle{primitive} : une fonction $F$ est une primitive de $f$ sur un intervalle $I$ lorsque $F$ est dérivable sur $I$ et que $F'(x) = f(x)$ pour tout $x$ de $I$. Nous avons aussi constaté qu'une fonction qui admet une primitive en admet en réalité une infinité, toutes différentes d'une constante.

Mais une question fondamentale reste en suspens : \textbf{à quelle condition une fonction $f$ admet-elle une primitive ?} Autrement dit, peut-on toujours « remonter » de $f$ vers une fonction $F$ dont $f$ serait la dérivée ? C'est à cette question que répond le théorème central de cette leçon, et la réponse est remarquablement simple : il suffit que $f$ soit \emph{continue}.

\courssub{Intuition : l'aire sous la courbe comme machine à fabriquer des primitives}

Imaginons une situation concrète d'atelier. Une pompe hydraulique débite un liquide avec un débit variable $f(t)$ (en litres par minute) au cours du temps $t$. Le débit $f$ est une fonction continue du temps : il ne saute pas brutalement d'une valeur à une autre. On se pose la question : quelle quantité totale de liquide $F(x)$ a été pompée entre l'instant initial $a$ (mise en route) et l'instant $x$ ?

Intuitivement, cette quantité est l'\cle{aire} comprise entre la courbe du débit, l'axe des abscisses et les droites verticales d'abscisses $a$ et $x$. Notons-la :
\[ F(x) = \int_a^x f(t)\, dt. \]

Voici l'idée géniale : si l'on avance le temps d'un tout petit instant $h$, la quantité pompée augmente d'environ $f(x) \times h$ (débit quasiment constant sur un intervalle très court, multiplié par la durée). Autrement dit :
\[ F(x+h) - F(x) \approx f(x) \times h \qquad\text{donc}\qquad \frac{F(x+h)-F(x)}{h} \approx f(x). \]
En faisant tendre $h$ vers $0$, le taux d'accroissement de $F$ tend exactement vers $f(x)$ : cela signifie précisément que $F$ est dérivable et que $F'(x) = f(x)$. \textbf{L'aire sous la courbe, vue comme fonction de la borne $x$, est une primitive de $f$.} C'est toute la substance du théorème.

\begin{popfigure}
\begin{center}
\begin{tikzpicture}
\begin{axis}[
  width=0.92\linewidth, height=6.2cm,
  axis lines=middle,
  xlabel={$t$}, ylabel={$f(t)$},
  xlabel style={at={(ticklabel* cs:1)}, anchor=west},
  ylabel style={at={(ticklabel* cs:1)}, anchor=south},
  xmin=-0.3, xmax=6.4, ymin=-0.4, ymax=3.4,
  xtick={1,4.2}, xticklabels={$a$,$x$},
  ytick=\empty,
  clip=false
]
\addplot[popcurve, domain=0.2:5.8, samples=200] {0.35*(x-1.2)^2 + 0.6};
\addplot[fill=popblueL, draw=none, domain=1:4.2, samples=100] {0.35*(x-1.2)^2 + 0.6} \closedcycle;
\draw[popdark, dashed] (axis cs:1,0) -- (axis cs:1,{0.35*(1-1.2)^2+0.6});
\draw[popdark, dashed] (axis cs:4.2,0) -- (axis cs:4.2,{0.35*(4.2-1.2)^2+0.6});
\node[popblue, font=\small] at (axis cs:2.6,1.1) {$F(x)=\displaystyle\int_a^x f(t)\,dt$};
\node[popink, font=\small, anchor=south] at (axis cs:5.3,3.0) {$f$};
\end{axis}
\end{tikzpicture}
\end{center}
\legende{La fonction $F(x)$ est l'aire (colorée) sous la courbe de $f$ entre $a$ et $x$. Quand $x$ varie, cette aire varie, et sa vitesse de variation est exactement $f(x)$.}
\end{popfigure}

\begin{popfigure}
\begin{center}
\begin{tikzpicture}
\begin{axis}[
  width=0.92\linewidth, height=6.2cm,
  axis lines=middle,
  xlabel={$t$}, ylabel={$f(t)$},
  xlabel style={at={(ticklabel* cs:1)}, anchor=west},
  ylabel style={at={(ticklabel* cs:1)}, anchor=south},
  xmin=-0.3, xmax=6.4, ymin=-0.4, ymax=3.4,
  xtick={1,3.2,4.2}, xticklabels={$a$,$x_1$,$x_2$},
  ytick=\empty,
  clip=false
]
\addplot[popcurve, domain=0.2:5.8, samples=200] {0.35*(x-1.2)^2 + 0.6};
\addplot[fill=poporangeL, draw=none, domain=1:3.2, samples=100] {0.35*(x-1.2)^2 + 0.6} \closedcycle;
\addplot[fill=popgreenL, draw=none, domain=3.2:4.2, samples=60] {0.35*(x-1.2)^2 + 0.6} \closedcycle;
\draw[popdark, dashed] (axis cs:1,0) -- (axis cs:1,{0.35*(1-1.2)^2+0.6});
\draw[popdark, dashed] (axis cs:3.2,0) -- (axis cs:3.2,{0.35*(3.2-1.2)^2+0.6});
\draw[popdark, dashed] (axis cs:4.2,0) -- (axis cs:4.2,{0.35*(4.2-1.2)^2+0.6});
\node[poporange, font=\small] at (axis cs:2.0,1.0) {$F(x_1)$};
\node[popgreen, font=\small] at (axis cs:3.7,1.4) {$F(x_2)-F(x_1)$};
\end{axis}
\end{tikzpicture}
\end{center}
\legende{Quand $x$ augmente de $x_1$ à $x_2$, l'aire augmente de la bande verte. Plus $f$ est grande, plus cette bande est haute, donc plus $F$ croît vite : c'est pourquoi $F' = f$.}
\end{popfigure}

\courssub{Énoncé rigoureux du théorème}

\begin{propriete}[Théorème d'existence des primitives]
Toute fonction \cle{continue} sur un intervalle $I$ admet au moins une \cle{primitive} sur $I$.

Plus précisément, si $f$ est continue sur $I$ et si $a$ est un point fixé de $I$, alors la fonction $F$ définie sur $I$ par
\[ F(x) = \int_a^x f(t)\, dt \]
est dérivable sur $I$, et pour tout $x$ de $I$ : $F'(x) = f(x)$. Autrement dit, $F$ est \textbf{la primitive de $f$ sur $I$ qui s'annule en $a$}.
\end{propriete}

Ce théorème est un pilier de l'analyse. Il garantit que dès qu'une fonction est continue — ce qui est le cas de toutes les fonctions usuelles rencontrées au programme (polynômes, fractions rationnelles sur leur ensemble de définition, exponentielle, logarithme, fonctions trigonométriques) — on est \emph{certain} qu'une primitive existe, même si l'on ne sait pas toujours l'écrire avec des formules simples.

\begin{attention}[Continuité : une hypothèse à ne pas oublier]
Le théorème exige la continuité sur un \textbf{intervalle}. Deux pièges classiques :
\begin{itemize}
\item Une fonction \emph{discontinue} peut ne pas admettre de primitive. Contre-exemple célèbre : la \textbf{fonction de Dirichlet}, qui vaut $1$ si $x$ est rationnel et $0$ si $x$ est irrationnel, n'est la dérivée d'aucune fonction.
\item Une fonction définie sur une \emph{réunion d'intervalles disjoints} (comme $x \mapsto \frac{1}{x}$ sur $\mathbb{R}^* = ]-\infty\,;\,0[ \cup ]0\,;\,+\infty[$) admet des primitives sur \textbf{chaque} intervalle séparément, avec des constantes \emph{différentes} possibles sur chacun. On ne peut pas écrire une seule constante $C$ valable partout.
\end{itemize}
Ne concluez donc jamais « $f$ admet une primitive » sans avoir vérifié (ou rappelé) que $f$ est continue sur l'intervalle considéré.
\end{attention}

\courssub{Preuve guidée : pourquoi $F$ est dérivable et $F' = f$}

La démonstration complète est formatrice : elle repose sur la définition même de la dérivée comme limite du taux d'accroissement. Suivons-la pas à pas.

\begin{experience}[Démonstration guidée du théorème fondamental]
Soit $f$ une fonction continue sur un intervalle $I$, $a$ un point fixé de $I$, et $F(x) = \displaystyle\int_a^x f(t)\,dt$. Fixons un point $x_0$ de $I$ et montrons que $F$ est dérivable en $x_0$, de dérivée $f(x_0)$.

\textbf{Étape 1 : exprimer l'accroissement de $F$.} Pour $h \neq 0$ assez petit pour que $x_0 + h$ reste dans $I$, la relation de Chasles pour les intégrales donne :
\[ F(x_0 + h) - F(x_0) = \int_a^{x_0+h} f(t)\,dt - \int_a^{x_0} f(t)\,dt = \int_{x_0}^{x_0+h} f(t)\,dt. \]
Géométriquement, c'est l'aire de la fine bande verte de la seconde figure.

\textbf{Étape 2 : encadrer l'aire de la bande.} Supposons $h > 0$ (le cas $h < 0$ se traite de même). Comme $f$ est continue sur le segment $[x_0\,;\,x_0+h]$, elle y atteint un minimum $m_h$ et un maximum $M_h$. L'aire de la bande est alors comprise entre celle du petit rectangle et celle du grand rectangle de même base $h$ :
\[ m_h \times h \;\le\; \int_{x_0}^{x_0+h} f(t)\,dt \;\le\; M_h \times h. \]

\textbf{Étape 3 : diviser par $h$ pour faire apparaître le taux d'accroissement.}
\[ m_h \;\le\; \frac{F(x_0+h)-F(x_0)}{h} \;\le\; M_h. \]

\textbf{Étape 4 : passer à la limite grâce à la continuité.} Quand $h$ tend vers $0$, l'intervalle $[x_0\,;\,x_0+h]$ se réduit au point $x_0$. Par continuité de $f$ en $x_0$, le minimum $m_h$ et le maximum $M_h$ tendent tous les deux vers $f(x_0)$. Par le théorème des gendarmes :
\[ \lim_{h \to 0} \frac{F(x_0+h)-F(x_0)}{h} = f(x_0). \]

\textbf{Conclusion.} $F$ est dérivable en $x_0$ et $F'(x_0) = f(x_0)$. Le point $x_0$ étant quelconque dans $I$, on a bien $F' = f$ sur $I$ tout entier. $\blacksquare$
\end{experience}

\begin{methode}[Construire une primitive grâce à l'intégrale à borne variable]
Pour exhiber une primitive d'une fonction $f$ continue sur un intervalle $I$ :
\begin{enumerate}
\item Vérifier (ou justifier) que $f$ est \textbf{continue} sur $I$ : c'est l'hypothèse du théorème.
\item Choisir un point $a$ de $I$ (souvent $0$, ou la borne gauche de $I$).
\item Poser $F(x) = \displaystyle\int_a^x f(t)\,dt$ : c'est une primitive de $f$, et c'est même la \textbf{seule} qui vérifie $F(a) = 0$.
\item Si l'on cherche la primitive vérifiant une condition $F(x_0) = y_0$, ajuster la constante : $G(x) = F(x) - F(x_0) + y_0$.
\end{enumerate}
\end{methode}

\courssub{Le choix de l'origine $a$ et les primitives à une constante près}

Que se passe-t-il si l'on change le point de départ $a$ ? Prenons un autre point $b$ de $I$ et posons $G(x) = \displaystyle\int_b^x f(t)\,dt$. Par la relation de Chasles :
\[ G(x) = \int_b^x f(t)\,dt = \int_b^a f(t)\,dt + \int_a^x f(t)\,dt = F(x) + \underbrace{\int_b^a f(t)\,dt}_{\text{constante}}. \]
Changer l'origine $a$ revient donc simplement à ajouter une \textbf{constante} à $F$. On retrouve ainsi, par une voie nouvelle, le résultat de la section précédente :

\begin{aretenir}[Toutes les primitives d'une fonction continue]
Si $f$ est continue sur un intervalle $I$ et si $F$ est une primitive de $f$ sur $I$, alors l'ensemble de \textbf{toutes} les primitives de $f$ sur $I$ est exactement l'ensemble des fonctions
\[ x \longmapsto F(x) + C, \qquad C \in \mathbb{R}. \]
Deux primitives de $f$ sur $I$ diffèrent donc toujours d'une constante. C'est pourquoi on écrit la « primitive générale » sous la forme $F(x) + C$.
\end{aretenir}

\begin{savaistu}[Pourquoi écrit-on $\int f(x)\,dx = F(x) + C$ ?]
La notation $\displaystyle\int f(x)\,dx = F(x) + C$, appelée « intégrale indéfinie », n'a de sens rigoureux que grâce au théorème d'existence : c'est lui qui garantit qu'une telle fonction $F$ existe dès que $f$ est continue. Le symbole $\int$, un « S » allongé signifiant \emph{somme}, fut introduit par Leibniz à la fin du XVII\textsuperscript{e} siècle. Le « $+C$ » rappelle que la primitive n'est jamais unique : toute la famille des primitives s'obtient en translatant verticalement le graphe de l'une d'entre elles. Le lien entre aire et primitive, appelé \textbf{théorème fondamental du calcul intégral}, a transformé le calcul des aires — problème vieux de deux mille ans, depuis Archimède — en un simple calcul de primitives.
\end{savaistu}

\courssub{Lien avec le théorème fondamental du calcul intégral}

Le résultat que nous venons d'établir possède un corollaire immédiat et d'une utilité pratique immense, que nous admettons en Terminale technique sous la forme suivante :

\begin{propriete}[Théorème fondamental du calcul intégral (admis)]
Soit $f$ une fonction continue sur un intervalle $I$, $a$ et $b$ deux points de $I$, et $F$ \textbf{une primitive quelconque} de $f$ sur $I$. Alors :
\[ \int_a^b f(t)\,dt = F(b) - F(a) = \bigl[F(x)\bigr]_a^b. \]
\end{propriete}

En effet, si $G(x) = \int_a^x f(t)\,dt$ est la primitive qui s'annule en $a$ et si $F$ est une autre primitive, alors $F = G + C$, donc $F(b) - F(a) = G(b) - G(a) = G(b) = \int_a^b f(t)\,dt$. Ce théorème sera l'outil central de la leçon sur le calcul intégral : pour calculer une aire, il suffit de connaître une primitive.

\begin{exoresolu}[Existence et détermination d'une primitive avec condition initiale]
\textbf{Énoncé.} Sur un chantier de Bafoussam, le débit d'une pompe (en litres par minute) à l'instant $t$ (en minutes) est modélisé par $f(t) = 3t^2 + 2$, pour $t \in [0\,;\,10]$.
\begin{enumerate}
\item Justifier que $f$ admet une primitive sur $[0\,;\,10]$.
\item Déterminer la primitive $F$ de $f$ qui s'annule en $0$. Que représente $F(4)$ ?
\end{enumerate}
\tcblower
\textbf{Solution détaillée.}
\begin{enumerate}
\item La fonction $f$ est une fonction \textbf{polynôme}, donc elle est \textbf{continue} sur $\mathbb{R}$, et en particulier sur l'intervalle $[0\,;\,10]$. D'après le \textbf{théorème d'existence des primitives}, $f$ admet au moins une primitive sur $[0\,;\,10]$.
\item Cherchons $F$ telle que $F'(t) = f(t)$. Une primitive de $3t^2$ est $t^3$ (car $(t^3)' = 3t^2$) et une primitive de $2$ est $2t$. Donc les primitives de $f$ sont les fonctions
\[ F(t) = t^3 + 2t + C, \qquad C \in \mathbb{R}. \]
La condition $F(0) = 0$ donne $0 + 0 + C = 0$, soit $C = 0$. Ainsi :
\[ F(t) = t^3 + 2t. \]
Vérifié : $F'(t) = 3t^2 + 2 = f(t)$.

$F(4) = 4^3 + 2 \times 4 = 64 + 8 = 72$. Concrètement, $F(4) = \displaystyle\int_0^4 f(t)\,dt$ représente la \textbf{quantité totale de liquide pompée} pendant les 4 premières minutes, soit 72 litres.
\end{enumerate}
\end{exoresolu}

\begin{attention}[Erreurs fréquentes à éviter]
\begin{itemize}
\item \textbf{Croire que toute fonction admet une primitive.} C'est faux : la continuité est indispensable. La fonction de Dirichlet (valant $1$ sur les rationnels, $0$ sur les irrationnels) n'admet aucune primitive. De même, une fonction en escalier discontinue n'est la dérivée d'aucune fonction dérivable.
\item \textbf{Oublier la constante $C$.} Écrire « la primitive de $f$ est $F$ » sans préciser la condition initiale est un abus : il y a une infinité de primitives. À l'examen, si une condition du type $F(x_0) = y_0$ est donnée, elle sert précisément à déterminer $C$.
\item \textbf{Confondre la variable d'intégration et la borne.} Dans $F(x) = \int_a^x f(t)\,dt$, la variable de $F$ est la borne $x$ ; la lettre $t$ est une variable « muette » : on aurait pu écrire $\int_a^x f(u)\,du$. Écrire $\int_a^x f(x)\,dx$ avec la même lettre aux deux rôles est une faute de rédaction classique.
\item \textbf{Appliquer le théorème sur une réunion d'intervalles.} Sur $\mathbb{R}^* = ]-\infty\,;\,0[ \cup ]0\,;\,+\infty[$, les constantes peuvent différer de part et d'autre de $0$.
\end{itemize}
\end{attention}

\begin{exercice}[Pour s'entraîner]
\begin{enumerate}
\item Soit $g$ la fonction définie sur $\mathbb{R}$ par $g(x) = 2x + \cos x$. Justifier que $g$ admet une primitive sur $\mathbb{R}$, puis déterminer la primitive $G$ vérifiant $G(0) = 3$.
\item Soit $f$ continue sur $[0\,;\,5]$ et $F(x) = \displaystyle\int_1^x f(t)\,dt$. Que vaut $F(1)$ ? Exprimer $\displaystyle\int_1^3 f(t)\,dt$ à l'aide de $F$.
\end{enumerate}
\end{exercice}

\begin{corrige}[Exercice]
\begin{enumerate}
\item $g$ est la somme d'une fonction affine et de la fonction cosinus, toutes deux continues sur $\mathbb{R}$ : $g$ est donc continue sur $\mathbb{R}$, et le théorème d'existence garantit qu'elle admet une primitive. Les primitives de $g$ sont $G(x) = x^2 + \sin x + C$. La condition $G(0) = 3$ donne $0 + 0 + C = 3$, soit $C = 3$. Donc $G(x) = x^2 + \sin x + 3$.
\item $F(1) = \displaystyle\int_1^1 f(t)\,dt = 0$ (intégrale sur un intervalle réduit à un point). Par définition, $\displaystyle\int_1^3 f(t)\,dt = F(3)$, et plus généralement $F(3) - F(1) = F(3)$.
\end{enumerate}
\end{corrige}

En résumé, la continuité est la clé qui ouvre la porte des primitives : toute fonction continue sur un intervalle possède une primitive, construite concrètement comme l'aire sous sa courbe à partir d'un point fixé, et toutes ses primitives s'en déduisent par ajout d'une constante. Dans la section suivante, nous exploiterons ce résultat pour établir les règles algébriques de calcul des primitives (linéarité, primitives des fonctions usuelles, formes composées).

\coursec{Primitives d'une fonction continue sur un intervalle}

\courssub{Rappels : dérivée et notion de primitive}

\begin{definition}[Dérivée d'une fonction]
Soit $f$ une fonction définie sur un intervalle $I$ et $a \in I$. Si le taux d'accroissement $\dfrac{f(x)-f(a)}{x-a}$ admet une limite finie lorsque $x \to a$, on dit que $f$ est \textbf{dérivable} en $a$. Cette limite est appelée \textbf{dérivée de $f$ en $a$} et se note $f'(a)$.
\[
f'(a) = \lim_{x \to a} \frac{f(x)-f(a)}{x-a}.
\]
La fonction $f'$ qui associe à $x \in I$ la dérivée $f'(x)$ (lorsqu'elle existe) est la \textbf{fonction dérivée} de $f$.
\end{definition}

\begin{definition}[Primitive d'une fonction]
Soit $f$ une fonction définie sur un intervalle $I$. Une fonction $F$ dérivable sur $I$ est une \textbf{primitive} de $f$ sur $I$ si et seulement si :
\[
\forall x \in I,\quad F'(x) = f(x).
\]
L'ensemble des primitives de $f$ sur $I$ se note $\displaystyle \int f(x)\,dx$ et s'écrit $F(x) + C$, où $C \in \mathbb{R}$ est une constante arbitraire.
\end{definition}

\begin{exemplebox}[Exemples élémentaires]
\begin{itemize}
    \item $f(x)=2x$ sur $\mathbb{R}$ : $F(x)=x^2$ est une primitive car $F'(x)=2x$. L'ensemble des primitives est $x^2 + C$.
    \item $f(x)=\cos x$ sur $\mathbb{R}$ : $F(x)=\sin x$ est une primitive. L'ensemble est $\sin x + C$.
    \item $f(x)=\dfrac{1}{x}$ sur $]0,+\infty[$ : $F(x)=\ln x$ est une primitive. L'ensemble est $\ln x + C$.
\end{itemize}
\end{exemplebox}

\begin{attention}[Notation et constante]
\begin{itemize}
    \item L'écriture $\int f(x)\,dx = F(x) + C$ signifie : « toute primitive de $f $ est de la forme $F(x)+C$ ».
    \item La constante $C$ est \textbf{indispensable} : deux primitives d'une même fonction sur un intervalle ne diffèrent que d'une constante.
    \item On ne doit jamais oublier le $dx$ qui indique la variable d'intégration.
\end{itemize}
\end{attention}

\courssub{Théorème d'existence des primitives pour les fonctions continues}

\begin{propriete}[Théorème d'existence (admis au Probatoire)]
Toute fonction \textbf{continue} sur un intervalle $I$ admet au moins une primitive sur $I$.
\end{propriete}

\begin{savaistu}[Pourquoi la continuité est-elle nécessaire ?]
Si $f$ n'est pas continue sur $I$, elle peut ne pas avoir de primitive. Par exemple, la fonction signe $f(x) = 1$ si $x \ge 0$ et $f(x) = -1$ si $x < 0$ n'a pas de primitive sur $\mathbb{R}$ car une fonction dérivable a la propriété des valeurs intermédiaires (théorème de Darboux), ce que $f$ n'a pas. En Terminale technique, on travaille presque exclusivement avec des fonctions continues (polynômes, rationnelles sur leur domaine, trigonométriques, exponentielles, logarithmes), donc l'existence est garantie.
\end{savaistu}

\begin{propriete}[Primitive nulle et constante]
Si $f$ est continue sur $I$ et $F$ est une primitive de $f$ sur $I$, alors :
\begin{itemize}
    \item $F$ est constante sur $I$ $\iff$ $f(x)=0$ pour tout $x \in I$.
    \item Deux primitives de $f$ sur $I$ diffèrent d'une constante.
\end{itemize}
\end{propriete}

\courssub{Propriétés algébriques des primitives}

\noindent L'opération de primitivation hérite de la linéarité de la dérivation. C'est la propriété fondamentale qui permet de décomposer une fonction complexe en somme de fonctions simples.

\begin{propriete}[Linéarité de la primitive]
Soient $f$ et $g$ deux fonctions continues sur un intervalle $I$, et $k \in \mathbb{R}$. Si $F$ et $G$ sont des primitives respectives de $f$ et $g$ sur $I$, alors :
\begin{itemize}
    \item $F + G$ est une primitive de $f + g$ sur $I$.
    \item $kF$ est une primitive de $kf$ sur $I$.
\end{itemize}
En notation intégrale :
\[
\int \bigl( f(x) + g(x) \bigr) \, dx = \int f(x) \, dx + \int g(x) \, dx
\qquad\text{et}\qquad
\int k\,f(x) \, dx = k \int f(x) \, dx .
\]
\emph{Preuve (exigible au Probatoire) :} Il suffit de dériver : $(F+G)' = F' + G' = f+g$ et $(kF)' = kF' = kf$.
\end{propriete}

\begin{exemplebox}[Application de la linéarité]
Déterminer une primitive de $f(x) = 3x^2 - 5\cos x + \dfrac{2}{x}$ sur $]0, +\infty[$.
\tcblower
On utilise la linéarité et le tableau des primitives usuelles (vu plus bas) :
\[
\int \left( 3x^2 - 5\cos x + \frac{2}{x} \right) dx
= 3\int x^2 dx - 5\int \cos x \, dx + 2\int \frac{1}{x} dx .
\]
\[
= 3 \cdot \frac{x^3}{3} - 5 \sin x + 2 \ln x + C
= x^3 - 5\sin x + 2\ln x + C .
\]
\emph{Vérification :} On dérive $F(x) = x^3 - 5\sin x + 2\ln x$ : $F'(x) = 3x^2 - 5\cos x + \frac{2}{x} = f(x)$.
\end{exemplebox}

\begin{attention}[Piège classique : oublier la constante ou la factoriser]
\begin{itemize}
    \item Ne pas écrire $\int 3x^2 dx = x^3$ sans le $+C$ final (une seule constante pour toute la somme).
    \item Ne pas factoriser une constante \emph{après} l'intégration : $\int 3x^2 dx = 3 \int x^2 dx = 3 \cdot \frac{x^3}{3} = x^3$, et non $3 \cdot \frac{x^3}{3} + C = x^3 + C$ (c'est correct mais lourd). La constante $C$ s'ajoute \emph{une seule fois} à la fin du calcul complet.
\end{itemize}
\end{attention}

\courssub{Détermination de primitives : méthodes et exercices types}

\noindent Au-delà de la linéarité, deux techniques majeures permettent de « faire entrer » une intégrale dans le tableau des primitives usuelles : la substitution (changement de variable) et l'intégration par parties.

\courssub{Primitive d'une fonction composée : la méthode de substitution}

\begin{methode}[Substitution (changement de variable) pour déterminer une primitive]
Pour calculer $\int f(g(x))\,g'(x) \, dx$ :
\begin{enumerate}
    \item \textbf{Identifier} la fonction intérieure $g(x)$ et vérifier que sa dérivée $g'(x)$ est présente (éventuellement à une constante près).
    \item \textbf{Poser} $u = g(x)$, puis calculer $du = g'(x)\,dx$.
    \item \textbf{Remplacer} $x$ et $dx$ par $u$ et $du$ : l'intégrale devient $\int f(u) \, du$.
    \item \textbf{Calculer} la primitive par rapport à $u$ : $F(u) + C$.
    \item \textbf{Revenir} à la variable $x$ en remplaçant $u$ par $g(x)$ : $F(g(x)) + C$.
\end{enumerate}
\emph{Variante :} Si l'intégrale est $\int k\,f(g(x))\,g'(x)\,dx$ avec $k$ constante, on factorise $k$ et on applique la méthode.
\end{methode}

\begin{exemplebox}[Substitution simple : forme $u^n u'$]
Calculer $\displaystyle \int x\,(x^2+1)^3 \, dx$.
\tcblower
\underline{Étape 1 :} On repère $g(x) = x^2+1$, dont la dérivée $g'(x) = 2x$ apparaît à facteur près ($x = \frac{1}{2}\cdot 2x$).
\underline{Étape 2 :} Posons $u = x^2+1$, alors $du = 2x\,dx$, donc $x\,dx = \frac{1}{2}du$.
\underline{Étape 3 :} Substitution :
\[
\int x\,(x^2+1)^3 \, dx = \int u^3 \cdot \frac{1}{2}du = \frac{1}{2}\int u^3 \, du .
\]
\underline{Étape 4 :} Intégration : $\frac{1}{2} \cdot \frac{u^4}{4} = \frac{u^4}{8}$.
\underline{Étape 5 :} Retour à $x$ : $\dfrac{(x^2+1)^4}{8} + C$.
\end{exemplebox}

\begin{exemplebox}[Substitution trigonométrique]
Calculer $\displaystyle \int \sin(3x-1) \, dx$.
\tcblower
Ici $g(x) = 3x-1$, $g'(x) = 3$. On a $\sin(3x-1) = \frac{1}{3} \cdot 3 \sin(3x-1)$.
Posons $u = 3x-1$, $du = 3\,dx \Rightarrow dx = \frac{1}{3}du$.
\[
\int \sin(3x-1) \, dx = \int \sin u \cdot \frac{1}{3}du = \frac{1}{3}(-\cos u) + C = -\frac{1}{3}\cos(3x-1) + C .
\]
\end{exemplebox}

\begin{exemplebox}[Forme $\frac{u'}{u}$ : logarithme]
Calculer $\displaystyle \int \frac{2x+1}{x^2+x+3} \, dx$ sur $\mathbb{R}$.
\tcblower
Le numérateur $2x+1$ est la dérivée du dénominateur $x^2+x+3$.
Posons $u = x^2+x+3$, $du = (2x+1)dx$.
\[
\int \frac{2x+1}{x^2+x+3} dx = \int \frac{du}{u} = \ln|u| + C = \ln(x^2+x+3) + C
\]
(le trinôme est toujours $>0$ car $\Delta = 1-12 = -11 < 0$).
\end{exemplebox}

\begin{attention}[Erreurs fréquentes en substitution]
\begin{itemize}
    \item \textbf{Oublier le retour à $x$ :} La réponse finale doit être exprimée en $x$, pas en $u$.
    \item \textbf{Mauvaise gestion du différentiel :} Écrire $du = 2x$ au lieu de $du = 2x\,dx$. Le $dx$ est indispensable pour la substitution.
    \item \textbf{Ne pas vérifier la présence de $g'(x)$ :} Si l'intégrale est $\int x^2 (x^3+1)^3 dx$, la substitution $u=x^3+1$ marche car $du=3x^2 dx$. Mais pour $\int x (x^3+1)^3 dx$, $x\,dx$ n'est pas le différentiel de $x^3+1$ : la substitution directe échoue.
\end{itemize}
\end{attention}

\courssub{Primitive par parties : formule $\int u\,dv = uv - \int v\,du$}

\noindent La dérivation d'un produit $(uv)' = u'v + uv'$ donne, en intégrant : $uv = \int u'v + \int uv'$. On en déduit la formule d'\textbf{intégration par parties}, qui permet d'échanger les rôles des facteurs : on dérive l'un et on intègre l'autre. C'est l'outil indispensable pour les produits « polynôme $\times$ exponentielle/trigonométrique/logarithme ».

\begin{methode}[Intégration par parties]
Pour calculer $\int u(x)\,v'(x) \, dx$ (ou $\int f(x)g'(x)\,dx$) :
\begin{enumerate}
    \item \textbf{Choisir} $u$ et $dv$ (ou $f$ et $g'$) selon la règle \textbf{LIATE} (voir encadré \texttt{Attention}).
    \item \textbf{Calculer} $du = u'(x)\,dx$ (en dérivant $u$) et $v = \int dv$ (en intégrant $dv$).
    \item \textbf{Appliquer} la formule : $\displaystyle \int u\,dv = uv - \int v\,du$.
    \item \textbf{Calculer} la nouvelle intégrale $\int v\,du$ (qui doit être plus simple que la première).
    \item \textbf{Ajouter} la constante $C$ à la fin.
\end{enumerate}
\end{methode}

\begin{exoresolu}[Calculer une primitive de $f(x)=x\cdot e^x$]
\underline{Choix de $u$ et $dv$ :} Selon LIATE, $x$ (Algébrique) se dérive, $e^x$ (Exponentielle) s'intègre.
Posons $u = x \Rightarrow du = dx$ et $dv = e^x dx \Rightarrow v = e^x$.
\tcblower
\underline{Application de la formule :}
\[
\int x e^x \, dx = x e^x - \int e^x \, dx = x e^x - e^x + C = e^x(x-1) + C .
\]
\underline{Vérification :} On dérive $F(x) = e^x(x-1)$ : $F'(x) = e^x(x-1) + e^x = xe^x - e^x + e^x = xe^x$. Correct.
\end{exoresolu}

\begin{attention}[Erreur fréquente – Mauvais choix de $u$ et $dv$]
La réussite de l'intégration par parties repose sur le choix judicieux des fonctions à dériver et à intégrer. Utilisez la mnémotechnique \textbf{LIATE} (ordre de priorité pour choisir $u$, celui qu'on \textbf{dérive}) :
\begin{center}
\begin{tabular}{|c|l|}
\hline
\textbf{L} & \textbf{L}ogarithme ($\ln x$, $\log x$) \\
\textbf{I} & \textbf{I}nverse trigonométrique ($\arcsin x$, $\arctan x$) \\
\textbf{A} & \textbf{A}lgébrique ($x$, $x^2$, $\sqrt{x}$, polynômes) \\
\textbf{T} & \textbf{T}rigonométrique ($\sin x$, $\cos x$) \\
\textbf{E} & \textbf{E}xponentielle ($e^x$, $a^x$) \\
\hline
\end{tabular}
\end{center}
\emph{Exemple d'erreur :} Pour $\int x e^x dx$, si on choisit $u=e^x$ (à dériver) et $dv=x\,dx$ (à intégrer), on obtient $du=e^x dx$, $v=x^2/2$, et la nouvelle intégrale $\int \frac{x^2}{2}e^x dx$ est \emph{plus compliquée} que la première. Le bon choix est $u=x$ (Algébrique avant Exponentielle).
\end{attention}

\begin{exemplebox}[Intégration par parties itérée : $\int x^2 e^x dx$]
Il faut appliquer la méthode deux fois.
\tcblower
1\textsuperscript{re} fois : $u=x^2 \Rightarrow du=2x\,dx$ ; $dv=e^x dx \Rightarrow v=e^x$.
$\int x^2 e^x dx = x^2 e^x - \int 2x e^x dx = x^2 e^x - 2\int x e^x dx$.
2\textsuperscript{e} fois : On sait que $\int x e^x dx = e^x(x-1)$ (exemple précédent).
$\int x^2 e^x dx = x^2 e^x - 2\bigl[ e^x(x-1) \bigr] + C = e^x(x^2 - 2x + 2) + C$.
\end{exemplebox}

\begin{exemplebox}[{Produit polynôme $\times$ logarithme : $\int x^2 \ln x \, dx$ sur $]0,+\infty[$}]
LIATE : $\ln x$ (Logarithme) avant $x^2$ (Algébrique) $\rightarrow$ on dérive $\ln x$.
$u = \ln x \Rightarrow du = \frac{1}{x}dx$ ; $dv = x^2 dx \Rightarrow v = \frac{x^3}{3}$.
\tcblower
\[
\int x^2 \ln x \, dx = \frac{x^3}{3}\ln x - \int \frac{x^3}{3}\cdot\frac{1}{x}dx = \frac{x^3}{3}\ln x - \frac{1}{3}\int x^2 dx = \frac{x^3}{3}\ln x - \frac{x^3}{9} + C .
\]
\end{exemplebox}

\courssub{Tableau récapitulatif des primitives usuelles}

\noindent Ce tableau est la « boîte à outils » de base. Tout élève de Terminale technique doit le connaître par cœur (sauf la colonne « Condition » qui rappelle le domaine de validité). Les primitives s'entendent sur tout intervalle où la fonction est continue.

\begin{center}
\begin{tabularx}{\linewidth}{|l|X|l|}\hline
\rowcolor{popblueL}
\textbf{Fonction $f(x)$} & \textbf{Primitive $F(x)$} & \textbf{Condition / Intervalle} \\ \hline
$k$ (constante) & $kx + C$ & $\mathbb{R}$ \\ \hline
$x^n$ ($n \in \mathbb{N}$) & $\dfrac{x^{n+1}}{n+1} + C$ & $\mathbb{R}$ \\ \hline
$x^{\alpha}$ ($\alpha \in \mathbb{Q} \setminus \{-1\}$) & $\dfrac{x^{\alpha+1}}{\alpha+1} + C$ & $]0,+\infty[$ (ou $\mathbb{R}^*$ si $\alpha = p/q$, $q$ impair) \\ \hline
$\dfrac{1}{x}$ & $\ln|x| + C$ & $\mathbb{R}^*$ (intervalle ne contenant pas 0) \\ \hline
$e^x$ & $e^x + C$ & $\mathbb{R}$ \\ \hline
$a^x$ ($a>0, a\neq 1$) & $\dfrac{a^x}{\ln a} + C$ & $\mathbb{R}$ \\ \hline
$\cos x$ & $\sin x + C$ & $\mathbb{R}$ \\ \hline
$\sin x$ & $-\cos x + C$ & $\mathbb{R}$ \\ \hline
$\dfrac{1}{\cos^2 x}$ & $\tan x + C$ & $]-\frac{\pi}{2}+k\pi, \frac{\pi}{2}+k\pi[$ \\ \hline
$\dfrac{1}{\sin^2 x}$ & $-\cot x + C$ & $]k\pi, (k+1)\pi[$ \\ \hline
$\dfrac{1}{1+x^2}$ & $\arctan x + C$ & $\mathbb{R}$ \\ \hline
$\dfrac{1}{\sqrt{1-x^2}}$ & $\arcsin x + C$ & $]-1, 1[$ \\ \hline
\end{tabularx}
\end{center}

\begin{savaistu}[Le logarithme népérien et l'intégrale de 1/x]
La primitive de $1/x$ sur $]0,+\infty[$ est $\ln x$. Sur $]-\infty,0[$, c'est $\ln(-x)$. On regroupe les deux cas en écrivant $\ln|x| + C$ sur $\mathbb{R}^*$. Attention : $\ln|x|$ n'est \textbf{pas} une primitive de $1/x$ sur $\mathbb{R}^*$ tout entier (qui n'est pas un intervalle) car la constante $C$ peut être différente de part et d'autre de 0. Au Probatoire, on précise toujours l'intervalle considéré.
\end{savaistu}

\courssub{Arbre de décision : quelle méthode choisir ?}

\noindent Face à une intégrale $\int f(x) dx$, l'élève doit adopter une démarche systématique. L'arbre de décision ci-dessous synthétise la stratégie à suivre. Il ne remplace pas l'intuition et l'entraînement, mais il structure la réflexion.

\begin{popfigure}
\begin{tikzpicture}[
    node distance=1.2cm and 1.5cm,
    decision/.style={diamond, draw=popdark, thick, fill=popblue!15, text width=3.5cm, align=center, inner sep=2pt},
    process/.style={rectangle, draw=popdark, thick, fill=popgreen!15, text width=3.5cm, align=center, inner sep=2pt, rounded corners},
    start/.style={ellipse, draw=popdark, thick, fill=poporange!20, text width=2.5cm, align=center, inner sep=2pt},
    arrow/.style={-Stealth, thick, draw=popdark},
    label/.style={font=\small, text=popdark, midway, fill=white, inner sep=1pt}
]
% Nœuds
\node[start] (start) {Intégrale $\int f(x)dx$};
\node[decision, below of=start] (tab) {Est-elle dans le \textbf{tableau des primitives usuelles} ?};
\node[process, below left=1.5cm and 2cm of tab, align=center] (direct){Application directe \\ du tableau + Linéarité};
\node[decision, below right=1.5cm and 2cm of tab, align=center] (comp){Est-ce une forme $f(g(x))g'(x)$ \\ (fonction composée $\times$ dérivée intérieure) ?};
\node[process, below left=1.5cm and 2cm of comp, align=center] (sub){\textbf{Substitution} $u=g(x)$ \\ puis tableau};
\node[decision, below right=1.5cm and 2cm of comp, align=center] (prod){Est-ce un \textbf{produit} de deux fonctions \\ (ex: poly $\times$ exp, poly $\times$ ln, poly $\times$ trig) ?};
\node[process, below left=1.5cm and 2cm of prod, align=center] (part){\textbf{Intégration par parties} \\ (règle LIATE pour choisir $u$)};
\node[process, below right=1.5cm and 2cm of prod, align=center] (autre){Autres techniques : \\ décomposition en éléments simples, \\ changement de variable trigonométrique, \\ identification (non au programme Tle Tech)};
\node[process, below=2cm of direct, align=center] (fin){Résultat $F(x)+C$ \\ \textbf{Vérifier par dérivation}};
% Flèches
\draw[arrow] (start) -- (tab);
\draw[arrow] (tab) -- node[label, left] {OUI} (direct);
\draw[arrow] (tab) -- node[label, above right] {NON} (comp);
\draw[arrow] (comp) -- node[label, left] {OUI} (sub);
\draw[arrow] (comp) -- node[label, above right] {NON} (prod);
\draw[arrow] (prod) -- node[label, left] {OUI} (part);
\draw[arrow] (prod) -- node[label, right] {NON} (autre);
\draw[arrow] (direct) -- (fin);
\draw[arrow] (sub) |- (fin);
\draw[arrow] (part) |- (fin);
\draw[arrow] (autre) |- (fin);
\end{tikzpicture}
\legende{Arbre de décision pour le choix de la méthode de primitivation.}
\end{popfigure}

\begin{methode}[Stratégie globale de recherche de primitive]
\begin{enumerate}
    \item \textbf{Simplifier} l'expression (développer, factoriser, utiliser les identités trigonométriques, réécrire les puissances/racines en exposants fractionnaires).
    \item \textbf{Regarder le tableau} : est-ce une forme immédiate ou une somme de formes immédiates (linéarité) ?
    \item \textbf{Chercher une composition} : voyez-vous une fonction et sa dérivée ? $\rightarrow$ \textbf{Substitution}.
    \item \textbf{Chercher un produit} : polynôme $\times$ (exponentielle, sinus, cosinus, logarithme) ? $\rightarrow$ \textbf{Intégration par parties} (LIATE).
    \item \textbf{Vérifier} systématiquement le résultat en dérivant.
\end{enumerate}
\end{methode}

\begin{exoresolu}[Synthèse : déterminer les primitives sur l'intervalle indiqué]
\begin{enumerate}
    \item $f(x) = \dfrac{2x+1}{x^2+x+3}$ sur $\mathbb{R}$.
    \item $g(x) = x^2 \ln x$ sur $]0, +\infty[$.
    \item $h(x) = \sin^2 x \cos x$ sur $\mathbb{R}$.
\end{enumerate}
\tcblower
\underline{1.} Le numérateur $2x+1$ est la dérivée du dénominateur $x^2+x+3$. Forme $\frac{u'}{u}$.
$u = x^2+x+3$, $du = (2x+1)dx$. $\int \frac{2x+1}{x^2+x+3} dx = \int \frac{du}{u} = \ln|u| + C = \ln(x^2+x+3) + C$ (le trinôme est toujours $>0$).

\underline{2.} Produit $x^2 \times \ln x$. LIATE : $\ln x$ (Logarithme) avant $x^2$ (Algébrique) $\rightarrow$ on dérive $\ln x$.
$u = \ln x \Rightarrow du = \frac{1}{x}dx$ ; $dv = x^2 dx \Rightarrow v = \frac{x^3}{3}$.
$\int x^2 \ln x \, dx = \frac{x^3}{3}\ln x - \int \frac{x^3}{3}\cdot\frac{1}{x}dx = \frac{x^3}{3}\ln x - \frac{1}{3}\int x^2 dx = \frac{x^3}{3}\ln x - \frac{x^3}{9} + C$.

\underline{3.} Forme $\sin^2 x \cos x = (\sin x)^2 \cdot \cos x$. $g(x)=\sin x$, $g'(x)=\cos x$. Substitution $u = \sin x$, $du = \cos x\,dx$.
$\int \sin^2 x \cos x \, dx = \int u^2 du = \frac{u^3}{3} + C = \frac{\sin^3 x}{3} + C$.
\end{exoresolu}

\courssub{Applications concrètes : modélisation de situations réelles}

\noindent La primitive permet de retrouver une grandeur à partir de sa variation instantanée (sa dérivée). C'est le cœur de la modélisation en technique.

\begin{exemplebox}[Électricité : charge d'un condensateur]
Dans un circuit RC série alimenté par une tension constante $E$, le courant $i(t)$ vérifie $E = Ri(t) + \frac{1}{C}q(t)$ où $q(t)$ est la charge du condensateur. On sait que $i(t) = q'(t)$. L'équation différentielle devient :
\[
R q'(t) + \frac{1}{C} q(t) = E.
\]
En cherchant une solution particulière constante $q_p = CE$ et en résolvant l'équation homogène, on trouve :
\[
q(t) = CE \left(1 - e^{-t/(RC)}\right) \quad \text{(pour } q(0)=0\text{)}.
\]
Ici, la primitive de l'exponentielle a permis d'exprimer la charge en fonction du temps. Le produit $RC$ (en secondes) est la \textbf{constante de temps} du circuit.
\end{exemplebox}

\begin{exemplebox}[Mécanique : du mouvement à la position]
Un mobile se déplace sur un axe. Son accélération est $a(t) = 6t - 4$ (en m/s²). On sait que $v(0) = 10$ m/s et $x(0) = 0$ m.
\tcblower
\underline{Vitesse :} $v(t) = \int a(t)\,dt = \int (6t-4)\,dt = 3t^2 - 4t + C_1$.
Condition $v(0)=10 \Rightarrow C_1 = 10$. Donc $v(t) = 3t^2 - 4t + 10$.
\underline{Position :} $x(t) = \int v(t)\,dt = \int (3t^2 - 4t + 10)\,dt = t^3 - 2t^2 + 10t + C_2$.
Condition $x(0)=0 \Rightarrow C_2 = 0$. Donc $x(t) = t^3 - 2t^2 + 10t$.
\underline{Interprétation :} À $t=2$ s, $v(2) = 12 - 8 + 10 = 14$ m/s, $x(2) = 8 - 8 + 20 = 20$ m.
\end{exemplebox}

\begin{exemplebox}[Gestion : coût total à partir du coût marginal]
Dans une entreprise, le coût marginal (coût de production d'une unité supplémentaire) est modélisé par $C_m(q) = 0,02q^2 - 3q + 150$ (en FCFA/unité), pour une production $q \in [0, 200]$. Les coûts fixes sont $CF = 50\,000$ FCFA.
\tcblower
Le coût total $CT(q)$ est une primitive du coût marginal : $CT'(q) = C_m(q)$.
\[
CT(q) = \int (0,02q^2 - 3q + 150)\,dq = 0,02\frac{q^3}{3} - 3\frac{q^2}{2} + 150q + K.
\]
La condition $CT(0) = CF = 50\,000$ donne $K = 50\,000$.
\[
CT(q) = \frac{0,02}{3}q^3 - 1,5q^2 + 150q + 50\,000.
\]
Pour $q=100$ : $CT(100) \approx 66,67 - 15\,000 + 15\,000 + 50\,000 = 50\,066,67$ FCFA.
\end{exemplebox}

\courssub{Synthèse, entraînement et évaluation formative}

\begin{aretenir}[Résumé des points clés]
\begin{itemize}
    \item \textbf{Définition :} $F$ est primitive de $f$ sur $I$ $\iff$ $F' = f$ sur $I$. Ensemble noté $\int f(x)dx = F(x) + C$.
    \item \textbf{Existence :} Toute fonction continue sur un intervalle admet des primitives (théorème admis).
    \item \textbf{Linéarité :} $\int (f+g) = \int f + \int g$ ; $\int kf = k\int f$.
    \item \textbf{Substitution :} $\int f(g(x))g'(x)dx = \int f(u)du$ avec $u=g(x)$. Repérer la forme « composée $\times$ dérivée intérieure ».
    \item \textbf{Intégration par parties :} $\int u\,dv = uv - \int v\,du$. Choisir $u$ selon LIATE (Log, InvTrig, Algébrique, Trig, Exp).
    \item \textbf{Tableau de référence :} Maîtriser les 12 primitives usuelles (puissances, $1/x$, exp, log, trig, trig inverses).
    \item \textbf{Vérification :} Dériver la primitive trouvée pour retrouver la fonction de départ.
    \item \textbf{Applications :} Reconstituer une grandeur (position, charge, coût total) à partir de sa dérivée (vitesse, courant, coût marginal) + conditions initiales.
\end{itemize}
\end{aretenir}

\begin{exercice}[Entraînement 1 : Primitives immédiates et linéarité]
Déterminer une primitive sur l'intervalle indiqué :
\begin{enumerate}
    \item $f(x) = 4x^3 - 3\sqrt{x} + \frac{5}{x}$ sur $]0,+\infty[$.
    \item $g(x) = 2\sin x - 3\cos x + e^x$ sur $\mathbb{R}$.
    \item $h(x) = \frac{1}{1+x^2} + \frac{1}{\sqrt{1-x^2}}$ sur $]-1,1[$.
\end{enumerate}
\end{exercice}

\begin{exercice}[Entraînement 2 : Substitution]
Calculer les primitives suivantes :
\begin{enumerate}
    \item $\displaystyle \int (2x+3)(x^2+3x+1)^4 \, dx$ sur $\mathbb{R}$.
    \item $\displaystyle \int \frac{\cos x}{\sin x} \, dx$ sur $]0,\pi[$.
    \item $\displaystyle \int x \sqrt{x^2-1} \, dx$ sur $]1,+\infty[$.
    \item $\displaystyle \int \frac{e^x}{1+e^{2x}} \, dx$ sur $\mathbb{R}$.
\end{enumerate}
\end{exercice}

\begin{exercice}[Entraînement 3 : Intégration par parties]
Calculer les primitives suivantes :
\begin{enumerate}
    \item $\displaystyle \int x \cos x \, dx$ sur $\mathbb{R}$.
    \item $\displaystyle \int \ln x \, dx$ sur $]0,+\infty[$ (astuce : écrire $\ln x = 1 \cdot \ln x$).
    \item $\displaystyle \int x^2 \sin x \, dx$ sur $\mathbb{R}$.
    \item $\displaystyle \int \arctan x \, dx$ sur $\mathbb{R}$.
\end{enumerate}
\end{exercice}

\begin{exercice}[Entraînement 4 : Problèmes concrets]
\begin{enumerate}
    \item \textbf{Chimie / Cinétique :} La vitesse de réaction $v(t)$ (en mol/L/s) suit la loi $v(t) = \frac{k}{(1+kt)^2}$ avec $k>0$. La concentration en réactif à $t=0$ est $C_0$. Exprimer la concentration $C(t)$ en fonction du temps.
    \item \textbf{Économie :} Le revenu marginal d'une entreprise est $R_m(q) = 100 - 0,5q$ (en kFCFA/unité). Le revenu pour $q=0$ est nul. Déterminer la fonction revenu total $R(q)$ et la production qui maximise le revenu.
    \item \textbf{Électricité :} La tension aux bornes d'une bobine d'inductance $L$ est $u(t) = U_0 \sin(\omega t)$. Sachant que $u(t) = L \frac{di}{dt}$ et $i(0)=0$, déterminer l'intensité $i(t)$.
\end{enumerate}
\end{exercice}

\begin{corrige}[Entraînement 1]
\begin{enumerate}
    \item $\int (4x^3 - 3x^{1/2} + 5/x) dx = x^4 - 2x^{3/2} + 5\ln x + C$.
    \item $\int (2\sin x - 3\cos x + e^x) dx = -2\cos x - 3\sin x + e^x + C$.
    \item $\int \left(\frac{1}{1+x^2} + \frac{1}{\sqrt{1-x^2}}\right) dx = \arctan x + \arcsin x + C$.
\end{enumerate}
\end{corrige}

\begin{corrige}[Entraînement 2]
\begin{enumerate}
    \item $u=x^2+3x+1$, $du=(2x+3)dx$ $\Rightarrow$ $\frac{1}{5}(x^2+3x+1)^5 + C$.
    \item $u=\sin x$, $du=\cos x\,dx$ $\Rightarrow$ $\ln|\sin x| + C = \ln(\sin x) + C$ (car $\sin x>0$ sur $]0,\pi[$).
    \item $u=x^2-1$, $du=2x\,dx$ $\Rightarrow$ $\frac{1}{3}(x^2-1)^{3/2} + C$.
    \item $u=e^x$, $du=e^x dx$ $\Rightarrow$ $\int \frac{du}{1+u^2} = \arctan(e^x) + C$.
\end{enumerate}
\end{corrige}

\begin{corrige}[Entraînement 3]
\begin{enumerate}
    \item $u=x$, $dv=\cos x\,dx$ $\Rightarrow$ $x\sin x + \cos x + C$.
    \item $u=\ln x$, $dv=dx$ $\Rightarrow$ $x\ln x - x + C$.
    \item Deux fois par parties : $-x^2\cos x + 2x\sin x + 2\cos x + C$.
    \item $u=\arctan x$, $dv=dx$ $\Rightarrow$ $x\arctan x - \frac{1}{2}\ln(1+x^2) + C$.
\end{enumerate}
\end{corrige}

\begin{corrige}[Entraînement 4]
\begin{enumerate}
    \item $C(t) = \int v(t)\,dt = \int \frac{k}{(1+kt)^2} dt$. Poser $u=1+kt$, $du=k\,dt$ $\Rightarrow$ $\int u^{-2} du = -u^{-1} + K = -\frac{1}{1+kt} + K$. $C(0)=C_0 \Rightarrow K = C_0 + 1$. Donc $C(t) = C_0 + 1 - \frac{1}{1+kt}$.
    \item $R(q) = \int (100 - 0,5q) dq = 100q - 0,25q^2 + K$. $R(0)=0 \Rightarrow K=0$. $R(q) = 100q - 0,25q^2$. Maximum pour $R'(q)=0 \Rightarrow 100 - 0,5q=0 \Rightarrow q=200$ unités.
    \item $i(t) = \frac{1}{L}\int U_0 \sin(\omega t) dt = -\frac{U_0}{L\omega}\cos(\omega t) + K$. $i(0)=0 \Rightarrow K = \frac{U_0}{L\omega}$. Donc $i(t) = \frac{U_0}{L\omega}(1 - \cos(\omega t))$.
\end{enumerate}
\end{corrige}

\begin{attention}[Auto-évaluation formative — Suis-je prêt pour le Probatoire ?]
Cochez les compétences que vous maîtrisez :
\begin{itemize}
    \item[$\square$] Je sais définir une primitive et écrire l'ensemble des primitives avec $+C$.
    \item[$\square$] J'applique correctement la linéarité (somme, factorisation par une constante).
    \item[$\square$] Je reconnais et traite les formes $\int u^n u'$, $\int \frac{u'}{u}$, $\int f(g(x))g'(x)$ par substitution.
    \item[$\square$] Je choisis correctement $u$ et $dv$ avec LIATE pour l'intégration par parties.
    \item[$\square$] Je connais par cœur le tableau des 12 primitives usuelles (y compris domaines).
    \item[$\square$] Je vérifie systématiquement mes primitives par dérivation.
    \item[$\square$] Je sais modéliser un problème concret (mouvement, électricité, économie) : écrire l'équation, intégrer, utiliser les conditions initiales.
    \item[$\square$] Je rédige mes réponses avec une justification claire (étapes, choix de méthode, conclusion).
\end{itemize}
Si une case n'est pas cochée : retravaillez les exemples correspondants et refaites les exercices d'entraînement.
\end{attention}

\coursec{Détermination de primitives : exercices types et astuces}

Dans la section précédente, nous avons établi les \cle{propriétés algébriques} des primitives : une primitive d'une somme est la somme des primitives, et l'on peut sortir un coefficient constant. Ces propriétés sont précieuses, mais elles ne suffisent pas toujours : face à une fonction comme $f(x)=\dfrac{2x+3}{x^2+3x+2}$ ou $g(x)=\cos^2 x$, aucune ligne du tableau de primitives ne semble convenir directement. Tout l'art de cette section consiste à \textbf{transformer l'expression} de la fonction pour la ramener à une ou plusieurs \cle{formes standards} du tableau. C'est exactement ce que fait un technicien en atelier : avant d'utiliser une machine, il prépare et ajuste la pièce brute.

\courssub{Reconnaître les formes standards}

\begin{aretenir}[Réflexe numéro 1 : lire le tableau à l'envers]
Chercher une primitive, c'est lire le tableau des dérivées \emph{de droite à gauche} : on se demande « de quelle fonction $f$ est-elle la dérivée ? ». Avant tout calcul, on compare la fonction donnée aux formes du tableau : puissance $x^n$, $\dfrac{1}{x}$, $\dfrac{1}{x^2}$, $\dfrac{1}{\sqrt{x}}$, $\cos x$, $\sin x$, $1+\tan^2 x = \dfrac{1}{\cos^2 x}$.
\end{aretenir}

Le tableau comparatif suivant résume les formes à reconnaître instantanément et montre la \cle{vérification} par dérivation, qui doit devenir un automatisme.

\begin{center}
\begin{tabularx}{\linewidth}{|X|X|X|}
\hline
\rowcolor{popblue!40} \textbf{Fonction $f(x)$} & \textbf{Primitive $F(x)$} & \textbf{Vérification $F'(x)=f(x)$} \\
\hline
\rowcolor{popblueL} $x^3$ & $\dfrac{x^4}{4}$ & $\left(\dfrac{x^4}{4}\right)'=\dfrac{4x^3}{4}=x^3$ \\
\hline
\rowcolor{popblueL} $\dfrac{1}{x^2}$ & $-\dfrac{1}{x}$ & $\left(-\dfrac{1}{x}\right)'=\dfrac{1}{x^2}$ \\
\hline
\rowcolor{popblueL} $\dfrac{1}{x}$ (pour $x>0$) & $\ln x$ & $(\ln x)'=\dfrac{1}{x}$ \\
\hline
\rowcolor{popblueL} $\cos x$ & $\sin x$ & $(\sin x)'=\cos x$ \\
\hline
\rowcolor{popblueL} $\sin x$ & $-\cos x$ & $(-\cos x)'=\sin x$ \\
\hline
\rowcolor{popblueL} $2x$ & $x^2$ & $(x^2)'=2x$ \\
\hline
\end{tabularx}
\end{center}

\begin{exemplebox}[Reconnaissance directe]
Déterminons une primitive de $f(x)=x^2-3x+\dfrac{2}{x^2}$ sur $]0\,;\,+\infty[$.

On reconnaît une somme de trois formes standards. Grâce aux propriétés algébriques :
\[ F(x)=\frac{x^3}{3}-3\times\frac{x^2}{2}+2\times\left(-\frac{1}{x}\right)=\frac{x^3}{3}-\frac{3x^2}{2}-\frac{2}{x}. \]
Vérification : $F'(x)=x^2-3x+\dfrac{2}{x^2}=f(x)$. \checkmark
\end{exemplebox}

\courssub{Se ramener à une forme connue par manipulation algébrique}

Très souvent, la fonction n'apparaît pas \emph{telle quelle} dans le tableau. Il faut alors la \textbf{réécrire}. Deux techniques sont au programme : la \cle{division} (ou réécriture du numérateur) et la \cle{décomposition en éléments simples}.

\textbf{Technique 1 : réécrire le numérateur.}\par
Quand le degré du numérateur est supérieur ou égal à celui du dénominateur, on effectue la division ou on « casse » la fraction.

\begin{exemplebox}[Casser une fraction]
Soit $f(x)=\dfrac{x^2+1}{x^2}$ sur $]0\,;\,+\infty[$. On écrit :
\[ f(x)=\frac{x^2}{x^2}+\frac{1}{x^2}=1+\frac{1}{x^2}. \]
On reconnaît deux formes standards, d'où $F(x)=x-\dfrac{1}{x}$. Vérification : $F'(x)=1+\dfrac{1}{x^2}=f(x)$. \checkmark
\end{exemplebox}

\textbf{Technique 2 : décomposition en éléments simples.}\par
Pour une fraction rationnelle dont le dénominateur se factorise en facteurs du premier degré distincts, on écrit la fraction comme somme de fractions plus simples, du type $\dfrac{a}{x-\alpha}+\dfrac{b}{x-\beta}$, dont les primitives sont des logarithmes.

\begin{methode}[Décomposer une fraction rationnelle]
Pour décomposer $f(x)=\dfrac{P(x)}{(x-\alpha)(x-\beta)}$ avec $P$ de degré $1$ et $\alpha \neq \beta$ :
\begin{enumerate}
\item Factoriser le dénominateur et identifier les racines $\alpha$ et $\beta$.
\item Écrire $f(x)=\dfrac{a}{x-\alpha}+\dfrac{b}{x-\beta}$ où $a$ et $b$ sont des constantes à déterminer.
\item Réduire au même dénominateur et identifier les numérateurs pour trouver $a$ et $b$ (système de deux équations).
\item Primitiver chaque morceau : une primitive de $\dfrac{a}{x-\alpha}$ est $a\ln|x-\alpha|$ (ou $a\ln(x-\alpha)$ si $x-\alpha>0$ sur l'intervalle considéré).
\item Vérifier en dérivant le résultat.
\end{enumerate}
\end{methode}

\begin{exoresolu}[Primitive de $f(x)=\dfrac{2x+3}{x^2+3x+2}$]
Énoncé : déterminer une primitive de $f(x)=\dfrac{2x+3}{x^2+3x+2}$ sur l'intervalle $]-1\,;\,+\infty[$.

\tcblower
\textbf{Solution détaillée.}

\textbf{Étape 1 : factoriser le dénominateur.} Le trinôme $x^2+3x+2$ a pour discriminant $\Delta=9-8=1$, donc pour racines $x_1=\dfrac{-3-1}{2}=-2$ et $x_2=\dfrac{-3+1}{2}=-1$. Ainsi $x^2+3x+2=(x+1)(x+2)$.

\textbf{Étape 2 : poser la décomposition.}
\[ \frac{2x+3}{(x+1)(x+2)}=\frac{a}{x+1}+\frac{b}{x+2}. \]

\textbf{Étape 3 : identifier.} En réduisant au même dénominateur :
\[ \frac{a(x+2)+b(x+1)}{(x+1)(x+2)}=\frac{(a+b)x+(2a+b)}{(x+1)(x+2)}. \]
On identifie avec $2x+3$ :
\[
\begin{cases}
a+b=2\\
2a+b=3
\end{cases}
\quad\Longrightarrow\quad a=1,\ b=1.
\]
Donc $f(x)=\dfrac{1}{x+1}+\dfrac{1}{x+2}$.

\textbf{Étape 4 : primitiver.} Sur $]-1\,;\,+\infty[$, on a $x+1>0$ et $x+2>0$, donc :
\[ F(x)=\ln(x+1)+\ln(x+2). \]

\textbf{Étape 5 : vérifier.}
\[ F'(x)=\frac{1}{x+1}+\frac{1}{x+2}=\frac{(x+2)+(x+1)}{(x+1)(x+2)}=\frac{2x+3}{x^2+3x+2}=f(x).\ \checkmark \]
\end{exoresolu}

\begin{attention}[Erreur fréquente : primitiver une fraction sans la décomposer]
Beaucoup d'élèves, pressés, écrivent qu'une primitive de $\dfrac{2x+3}{x^2+3x+2}$ serait $\ln(x^2+3x+2)$ « parce que la primitive de $1/u$ est $\ln u$ ». C'est \textbf{faux} en général ! En dérivant : $\big(\ln(x^2+3x+2)\big)'=\dfrac{2x+3}{x^2+3x+2}$... ici cela fonctionne \emph{uniquement parce que le numérateur est exactement la dérivée du dénominateur}. C'est un coup de chance, pas une méthode. La démarche correcte et toujours valable est la \cle{décomposition en éléments simples}. Au Baccalauréat technique, une primitive non vérifiée ou obtenue par un raccourci faux fait perdre des points même si le résultat est juste.
\end{attention}

\courssub{Utiliser les identités trigonométriques}

Les fonctions $\sin^2 x$ et $\cos^2 x$ ne figurent pas dans le tableau des primitives. L'astuce consiste à les \textbf{linéariser}, c'est-à-dire à les exprimer à l'aide de $\cos(2x)$, dont on connaît une primitive.

\begin{propriete}[Identités de linéarisation]
Pour tout réel $x$ :
\[ \cos^2 x=\frac{1+\cos(2x)}{2} \qquad\text{et}\qquad \sin^2 x=\frac{1-\cos(2x)}{2}. \]
Ces formules se déduisent de $\cos(2x)=2\cos^2 x - 1 = 1-2\sin^2 x$.
\end{propriete}

\begin{exoresolu}[Primitive de $\cos^2 x$]
Énoncé : déterminer une primitive de $g(x)=\cos^2 x$ sur $\mathbb{R}$.

\tcblower
\textbf{Solution.} On linéarise : $g(x)=\dfrac{1}{2}+\dfrac{1}{2}\cos(2x)$.

Une primitive de $\dfrac{1}{2}$ est $\dfrac{x}{2}$. Pour $\dfrac{1}{2}\cos(2x)$ : comme la dérivée de $\sin(2x)$ est $2\cos(2x)$, une primitive de $\cos(2x)$ est $\dfrac{\sin(2x)}{2}$. Donc :
\[ G(x)=\frac{x}{2}+\frac{1}{2}\times\frac{\sin(2x)}{2}=\frac{x}{2}+\frac{\sin(2x)}{4}. \]
Vérification : $G'(x)=\dfrac{1}{2}+\dfrac{2\cos(2x)}{4}=\dfrac{1+\cos(2x)}{2}=\cos^2 x$. \checkmark
\end{exoresolu}

\begin{attention}[Piège : la fonction composée]
Une primitive de $\cos(2x)$ n'est \textbf{pas} $\sin(2x)$ ! En dérivant $\sin(2x)$, on obtient $2\cos(2x)$ : il y a un facteur $2$ en trop. Il faut donc diviser par $2$ : une primitive de $\cos(2x)$ est $\dfrac{\sin(2x)}{2}$. Règle générale : une primitive de $\cos(ax+b)$ est $\dfrac{1}{a}\sin(ax+b)$, et une primitive de $\sin(ax+b)$ est $-\dfrac{1}{a}\cos(ax+b)$, pour $a\neq 0$.
\end{attention}

\courssub{Vérification systématique par dérivation}

\begin{methode}[Vérifier une primitive : le réflexe qui sauve des points]
Après avoir trouvé un candidat $F$ :
\begin{enumerate}
\item Dériver $F$ soigneusement, en appliquant les règles de dérivation (somme, produit, composée).
\item Simplifier $F'(x)$ (réduction au même dénominateur si nécessaire).
\item Comparer avec $f(x)$ : si $F'(x)=f(x)$ pour tout $x$ de l'intervalle, la primitive est correcte ; sinon, reprendre le calcul.
\end{enumerate}
Cette vérification ne coûte qu'une minute et garantit la justesse du résultat : c'est l'équivalent mathématique du contrôle qualité en bout de chaîne de production.
\end{methode}

\begin{popfigure}
\begin{center}
\begin{tikzpicture}[scale=0.95, every node/.style={font=\small}]
\node[draw=popblue, fill=popblueL, rounded corners, thick, minimum width=3.4cm, minimum height=1cm] (f) at (0,0) {Fonction $f$};
\node[draw=popgreen, fill=popgreenL, rounded corners, thick, minimum width=3.4cm, minimum height=1cm] (F) at (7,0) {Primitive $F$};
\draw[-{Stealth[length=3mm]}, very thick, popgreen] (f) -- node[above, popdark]{on cherche : $F'=f$} (F);
\draw[-{Stealth[length=3mm]}, very thick, poporange] (F.south) to[bend left=25] node[below, popdark]{on vérifie : $F'\stackrel{?}{=}f$} (f.south);
\end{tikzpicture}
\end{center}
\legende{Aller-retour fondamental : on \emph{cherche} $F$ en lisant le tableau à l'envers, puis on \emph{vérifie} en dérivant.}
\end{popfigure}

\begin{exercice}[Entraînement guidé]
\begin{enumerate}
\item Déterminer une primitive de $f(x)=\dfrac{3x-1}{x^2}$ sur $]0\,;\,+\infty[$ (penser à casser la fraction).
\item Déterminer une primitive de $g(x)=\sin^2 x$ sur $\mathbb{R}$.
\item Déterminer une primitive de $h(x)=\dfrac{x}{x+1}$ sur $]-1\,;\,+\infty[$ (écrire $x=(x+1)-1$).
\end{enumerate}
\end{exercice}

\begin{corrige}[Exercice : Entraînement guidé]
\textbf{1.} $f(x)=\dfrac{3x}{x^2}-\dfrac{1}{x^2}=\dfrac{3}{x}-\dfrac{1}{x^2}$, donc $F(x)=3\ln x+\dfrac{1}{x}$. Vérification : $F'(x)=\dfrac{3}{x}-\dfrac{1}{x^2}=f(x)$. \checkmark

\textbf{2.} $\sin^2 x=\dfrac{1-\cos(2x)}{2}$, donc $G(x)=\dfrac{x}{2}-\dfrac{\sin(2x)}{4}$. Vérification : $G'(x)=\dfrac{1}{2}-\dfrac{\cos(2x)}{2}=\sin^2 x$. \checkmark

\textbf{3.} $h(x)=\dfrac{(x+1)-1}{x+1}=1-\dfrac{1}{x+1}$, donc $H(x)=x-\ln(x+1)$. Vérification : $H'(x)=1-\dfrac{1}{x+1}=\dfrac{x}{x+1}=h(x)$. \checkmark
\end{corrige}

\begin{savaistu}[La primitive n'est jamais unique]
Si $F$ est une primitive de $f$ sur un intervalle, alors pour \emph{toute} constante réelle $C$, la fonction $F+C$ est aussi une primitive de $f$, car la dérivée d'une constante est nulle : $(F+C)'=F'+0=f$. Réciproquement, deux primitives d'une même fonction sur un intervalle diffèrent toujours d'une constante. C'est pourquoi on écrit souvent $F(x)+C$. Conséquence pratique : pour choisir \emph{la} primitive qui correspond à un problème concret (position d'un véhicule, niveau d'un réservoir, coût de production), il faut une information supplémentaire, appelée \cle{condition initiale}. Par exemple, si le réservoir contenait $50$ litres à l'instant $t=0$, on choisit la constante $C$ pour que $F(0)=50$.
\end{savaistu}

\begin{aretenir}[Bilan des astuces de cette section]
\begin{itemize}
\item \textbf{Somme de formes standards} : primitiver terme à terme.
\item \textbf{Fraction avec dénominateur simple} : casser la fraction ($\dfrac{P(x)}{x^2}$, $\dfrac{x}{x+1}$).
\item \textbf{Fraction rationnelle} : factoriser le dénominateur, décomposer en éléments simples, primitiver avec des logarithmes.
\item \textbf{$\cos^2 x$ et $\sin^2 x$} : linéariser avec $\cos(2x)$.
\item \textbf{$\cos(ax+b)$, $\sin(ax+b)$} : ne pas oublier le facteur $\dfrac{1}{a}$.
\item \textbf{Toujours} vérifier le résultat en dérivant.
\end{itemize}
\end{aretenir}

\coursec{Applications concrètes : modélisation de situations réelles}

Dans la section précédente, nous avons appris à \emph{déterminer} des primitives : lecture inverse du tableau des dérivées, linéarité, formes composées. Mais à quoi cela sert-il concrètement ? La réponse tient en une idée simple et puissante : dans la vie réelle, on mesure souvent un \cle{taux} (une vitesse, un débit, un coût marginal), c'est-à-dire une \emph{dérivée}. Pour retrouver la grandeur elle-même (distance, quantité produite, coût total), il faut \og remonter la dérivation \fg{} : c'est exactement le rôle des primitives. Cette section montre, sur des situations tirées des métiers techniques et de la vie courante au Cameroun, comment la recherche de primitives devient un outil de modélisation.

\courssub{Principe général : remonter du taux à la grandeur}

\begin{aretenir}[Idée directrice]
Si une grandeur $G$ varie et que l'on connaît son \cle{taux de variation} $g(t) = G'(t)$, alors $G$ est une primitive de $g$. La variation totale de $G$ entre deux instants $a$ et $b$ s'obtient par :
\[ G(b) - G(a) = F(b) - F(a) \quad \text{où } F \text{ est une primitive de } g. \]
On note cette quantité $\displaystyle \int_a^b g(t)\,\mathrm{d}t = F(b) - F(a)$ : c'est \og l'accumulation \fg{} du taux entre $a$ et $b$.
\end{aretenir}

L'intuition est celle du compteur d'eau ou d'électricité : le compteur affiche la \emph{consommation totale}, mais ce que l'on mesure instantanément, c'est le \emph{débit}. Consommation totale = accumulation du débit dans le temps. De même : distance = accumulation de la vitesse ; coût total = accumulation du coût marginal.

\begin{attention}[La constante d'intégration et la condition initiale]
Une primitive n'est jamais unique : elle est définie \emph{à une constante près}. Dans un problème concret, cette constante se fixe grâce à une \cle{condition initiale}. Par exemple, si la position est $x(t) = F(t) + C$ et que le mobile part du point $x(0) = x_0$, alors $C = x_0 - F(0)$. Oublier la condition initiale, c'est donner une réponse définie \og à un décalage près \fg{}, ce qui est insuffisant dans une situation réelle.
\end{attention}

\courssub{Distance parcourue à partir de la vitesse}

\begin{propriete}[Vitesse et distance]
Si un mobile se déplace avec une vitesse $v(t)$ (fonction continue du temps), alors la \cle{distance parcourue} entre les instants $t = a$ et $t = b$ est :
\[ D = \int_a^b v(t)\,\mathrm{d}t = V(b) - V(a), \]
où $V$ est une primitive de $v$. Autrement dit : la position $x(t)$ est une primitive de la vitesse, et la distance parcourue est la variation de la position.
\end{propriete}

\textbf{Justification (formatrice).} Par définition physique, la vitesse est la dérivée de la position : $v(t) = x'(t)$. Donc $x$ est une primitive de $v$, et la distance parcourue entre $a$ et $b$ est $x(b) - x(a)$, c'est-à-dire la variation d'une primitive de $v$ entre $a$ et $b$. \hfill $\blacksquare$

\begin{exemplebox}[Chauffeur de bus Yaoundé--Douala]
Un bus roule à la vitesse constante $v(t) = 80$ km/h pendant 3 heures. Une primitive de $v$ est $V(t) = 80t$. La distance parcourue est :
\[ D = V(3) - V(0) = 80 \times 3 - 0 = 240 \text{ km}. \]
On retrouve la formule bien connue $D = v \times t$ : le calcul par primitive \emph{généralise} cette formule au cas où la vitesse n'est pas constante.
\end{exemplebox}

\courssub{Le problème du moto-taxi}

Voici la situation centrale de cette leçon. Un \emph{moto-taxi} (benskin) circule en ville ; à cause des ralentissements et des accélérations, sa vitesse varie selon :
\[ v(t) = 10 + 2\sin(t) \quad \text{(en km/h, avec } t \text{ en heures, } t \in [0\,;\,2]\text{)}. \]
Quelle distance parcourt-il pendant ces deux heures ?

\begin{methode}[De la vitesse à la distance : démarche en 4 étapes]
\begin{enumerate}
\item \textbf{Identifier le taux} : la vitesse $v(t)$ est la dérivée de la position. La distance cherchée est la variation d'une primitive de $v$ entre les bornes.
\item \textbf{Poser une primitive} : chercher $V(t)$ telle que $V'(t) = v(t)$, terme à terme.
\item \textbf{Appliquer les bornes} : calculer $D = V(b) - V(a)$.
\item \textbf{Interpréter} : donner le résultat avec son unité et vérifier sa cohérence (ordre de grandeur).
\end{enumerate}
\end{methode}

\begin{exoresolu}[Distance parcourue par le moto-taxi]
Énoncé. La vitesse d'un moto-taxi est $v(t) = 10 + 2\sin(t)$ km/h pour $t \in [0\,;\,2]$ (heures, angles en radians). Calculer la distance totale parcourue.
\tcblower
\textbf{Étape 1.} La distance est la variation d'une primitive de $v$ entre $0$ et $2$ :
\[ D = \int_0^2 \bigl(10 + 2\sin(t)\bigr)\,\mathrm{d}t. \]
\textbf{Étape 2.} Primitive terme à terme : une primitive de $10$ est $10t$ ; une primitive de $2\sin(t)$ est $-2\cos(t)$ (car $(-2\cos t)' = 2\sin t$). Donc :
\[ V(t) = 10t - 2\cos(t). \]
\textbf{Étape 3.} Application des bornes :
\begin{align*}
D &= V(2) - V(0) = \bigl(10 \times 2 - 2\cos(2)\bigr) - \bigl(0 - 2\cos(0)\bigr) \\
&= 20 - 2\cos(2) + 2 = 22 - 2\cos(2).
\end{align*}
Avec $\cos(2) \approx -0{,}4161$ (calculatrice en mode \emph{radians}) :
\[ D \approx 22 - 2 \times (-0{,}4161) \approx 22{,}83 \text{ km}. \]
\textbf{Étape 4.} Interprétation : le moto-taxi parcourt environ \SI{22,83}{km} en 2 heures. Vérification de cohérence : sa vitesse oscille entre $8$ et $12$ km/h, donc la distance doit être comprise entre $8 \times 2 = 16$ km et $12 \times 2 = 24$ km. Notre résultat $22{,}83$ km est bien dans cet intervalle : il est plausible.
\end{exoresolu}

\begin{popfigure}
\begin{center}
\begin{tikzpicture}
\begin{axis}[
  width=0.9\linewidth, height=6.2cm,
  axis lines=middle,
  xlabel={$t$ (h)}, ylabel={$v$ (km/h)},
  xmin=-0.15, xmax=2.35, ymin=0, ymax=13.5,
  xtick={0,0.5,1,1.5,2}, ytick={2,4,6,8,10,12},
  domain=0:2, samples=200, clip=false]
  \addplot[popcurve, fill=popblueL, fill opacity=0.55, draw=popblue, very thick]
    {10 + 2*sin(deg(x))} \closedcycle;
  \addplot[popcurve, very thick] {10 + 2*sin(deg(x))};
  \draw[dashed, popmuted] (axis cs:2,0) -- (axis cs:2,10.83);
  \node[popdark] at (axis cs:1,5) {\textbf{distance} $= \displaystyle\int_0^2 v(t)\,\mathrm{d}t \approx 22{,}83$ km};
  \node[popblue, anchor=south west] at (axis cs:0.55,11.6) {$v(t) = 10 + 2\sin(t)$};
\end{axis}
\end{tikzpicture}
\end{center}
\legende{L'aire sous la courbe de la vitesse entre $t=0$ et $t=2$ représente la distance parcourue.}
\end{popfigure}

\begin{attention}[Erreur fréquente : confondre vitesse moyenne et distance totale]
Deux pièges classiques :
\begin{itemize}
\item Calculer la \emph{vitesse moyenne} $\dfrac{v(0)+v(2)}{2}$ et s'arrêter là : ce n'est \textbf{pas} la distance. La vitesse moyenne est un quotient $\dfrac{D}{b-a}$ ; la distance est l'accumulation $D = \displaystyle\int_a^b v(t)\,\mathrm{d}t$.
\item Multiplier $v(0)$ par la durée : cela suppose une vitesse constante, ce qui est faux ici.
\end{itemize}
Autre piège : laisser la calculatrice en mode \emph{degrés}. Ici $t$ est en heures et le sinus est pris en \textbf{radians} : $\cos(2\ \text{rad}) \approx -0{,}4161$, et non $\cos(2°) \approx 0{,}9994$.
\end{attention}

\courssub{Coût total à partir du coût marginal (économie et gestion)}

\begin{definition}[Coût marginal]
En économie, le \cle{coût marginal} $C_m(q)$ est le coût de production d'une unité supplémentaire lorsqu'on produit déjà $q$ unités. Mathématiquement, c'est la dérivée du coût total : $C_m(q) = C'(q)$. Réciproquement, le coût total est une primitive du coût marginal :
\[ C(q) = \int C_m(q)\,\mathrm{d}q + C_0, \]
où la constante $C_0$ représente les \emph{coûts fixes} (loyer de l'atelier, machines), c'est-à-dire le coût quand $q = 0$.
\end{definition}

\begin{exemplebox}[Menuiserie de Douala]
Un menuisier fabrique des chaises. Son coût marginal est $C_m(q) = 2000 + 20q$ (en FCFA par chaise), et ses coûts fixes sont de \SI{50000}{FCFA}. Une primitive de $C_m$ est $C(q) = 2000q + 10q^2 + K$. La condition initiale $C(0) = 50000$ donne $K = 50000$. Le coût total pour 30 chaises est :
\[ C(30) = 2000 \times 30 + 10 \times 30^2 + 50000 = 60000 + 9000 + 50000 = \SI{119000}{FCFA}. \]
Ici, la constante d'intégration a un sens économique précis : les coûts fixes.
\end{exemplebox}

\courssub{Aire sous une courbe (géométrie)}

\begin{propriete}[Aire et primitive]
Soit $f$ une fonction continue et \emph{positive} sur $[a\,;\,b]$, de courbe représentative $\mathcal{C}_f$ dans un repère orthogonal. L'\cle{aire} du domaine délimité par $\mathcal{C}_f$, l'axe des abscisses et les droites $x = a$ et $x = b$ vaut, en unités d'aire :
\[ \mathcal{A} = \int_a^b f(x)\,\mathrm{d}x = F(b) - F(a), \quad \text{où } F \text{ est une primitive de } f. \]
C'est exactement ce que montre la figure du moto-taxi : la distance est l'aire sous la courbe de la vitesse.
\end{propriete}

\begin{exoresolu}[Aire d'une tôle découpée]
Énoncé. Un tôlier doit découper une pièce dont le bord supérieur suit la courbe $f(x) = 3x^2 + 1$ (en dm) pour $x \in [0\,;\,2]$, le bord inférieur étant l'axe des abscisses. Calculer l'aire de la pièce.
\tcblower
$f$ est positive sur $[0\,;\,2]$ (somme de termes positifs). Une primitive de $f$ est $F(x) = x^3 + x$, car $F'(x) = 3x^2 + 1 = f(x)$. Donc :
\[ \mathcal{A} = F(2) - F(0) = (8 + 2) - 0 = 10 \text{ dm}^2. \]
La pièce a une aire de $10$ dm$^2$, soit \SI{0,10}{m^2}. Cette information permet au tôlier de chiffrer la matière première nécessaire.
\end{exoresolu}

\begin{savaistu}[De l'accélération à la position : la chaîne des primitives]
En physique, les grandeurs du mouvement s'enchaînent par dérivation : la \emph{vitesse} est la dérivée de la \emph{position}, et l'\emph{accélération} est la dérivée de la \emph{vitesse}. En remontant la chaîne par des primitives :
\[ \text{accélération} \xrightarrow{\ \text{primitive}\ } \text{vitesse} \xrightarrow{\ \text{primitive}\ } \text{position}. \]
C'est le principe des centrales inertielles embarquées dans les avions et les téléphones : on mesure l'accélération, puis on intègre deux fois pour retrouver la position. Chaque intégration fait apparaître une constante, fixée par les conditions initiales (vitesse initiale, position initiale) : sans elles, nulle navigation possible !
\end{savaistu}

\begin{exercice}[Pour s'entraîner]
Une pompe remplit un réservoir avec un débit $d(t) = 50 - 5t$ litres par minute, pour $t \in [0\,;\,6]$. Le réservoir contient initialement 100 litres.
\begin{enumerate}
\item Déterminer la quantité d'eau $Q(t)$ dans le réservoir à l'instant $t$ (primitive + condition initiale).
\item Quelle quantité d'eau a été ajoutée au bout de 6 minutes ? Quel est alors le contenu total du réservoir ?
\end{enumerate}
\end{exercice}

\begin{corrige}[Exercice : remplissage du réservoir]
1. Une primitive de $d(t) = 50 - 5t$ est $Q(t) = 50t - \dfrac{5}{2}t^2 + C$. La condition initiale $Q(0) = 100$ donne $C = 100$, donc $Q(t) = 50t - 2{,}5t^2 + 100$.

2. Quantité ajoutée : $Q(6) - Q(0) = (300 - 90 + 100) - 100 = 210$ litres. Contenu total : $Q(6) = 310$ litres.
\end{corrige}

\coursec{Synthèse, entraînement et évaluation formative}

Après avoir exploré les applications concrètes des primitives dans la modélisation de situations réelles (calcul de déplacements, d'énergies, de coûts marginaux), nous consolidons maintenant l'ensemble des acquis de cette leçon. Cette section vous propose une vue d'ensemble structurée, un entraînement progressif calibré sur les attendus du Probatoire, une grille d'auto-évaluation pour mesurer votre maîtrise, et les clés pour éviter les pièges classiques le jour de l'examen.

\courssub{Récapitulatif synthétique : la carte mentale de la leçon}

L'illustration ci-dessous résume les articulations logiques entre la définition, le théorème d'existence, les propriétés algébriques, les méthodes de calcul et les applications. Utilisez-la comme support de révision active : fermez le cours et essayez de la reconstruire de mémoire.

\begin{popfigure}
\begin{tikzpicture}[
    mindmap,
    concept color=popblue,
    level 1/.style={level distance=4.5cm, sibling angle=90},
    level 2/.style={level distance=3.2cm, sibling angle=45},
    every node/.style={font=\small},
    popblueL/.style={concept color=popblue!80!popcream},
    poporangeL/.style={concept color=poporange!80!popcream},
    popgreenL/.style={concept color=popgreen!80!popcream},
    poppurpleL/.style={concept color=poppurple!80!popcream},
    poppinkL/.style={concept color=poppink!80!popcream},
]
\node[concept, align=center]{Primitives\\Fonction continue}
    child[popblueL] { node[concept] {Définition}
        child { node[concept] {$F'(x)=f(x)$} }
        child { node[concept] {Intervalle $I$} }
        child { node[concept] {Constante $C$} }
    }
    child[poporangeL] { node[concept, align=center]{Théorème\\d'existence}
        child { node[concept] {$f$ continue sur $I$} }
        child { node[concept] {$\Rightarrow \exists F$ primitive} }
        child { node[concept] {Toutes primitives : $F+C$} }
    }
    child[popgreenL] { node[concept, align=center]{Propriétés\\algébriques}
        child { node[concept] {Linéarité} }
        child { node[concept] {$\lambda F + \mu G$} }
        child { node[concept] {Opérations usuelles} }
    }
    child[poppurpleL] { node[concept, align=center]{Méthodes\\de calcul}
        child { node[concept] {Table de référence} }
        child { node[concept] {Substitution $u=g(x)$} }
        child { node[concept] {Intégration par parties} }
        child { node[concept] {Décomposition rationnelle} }
    }
    child[poppinkL] { node[concept, align=center]{Applications\\concrètes}
        child { node[concept] {Mécanique : $v \to x$} }
        child { node[concept] {Électricité : $i \to q$} }
        child { node[concept] {Gestion : $C_m \to C$} }
        child { node[concept] {Conditions initiales} }
    };
\end{tikzpicture}
\legende{Carte mentale récapitulative : définition, théorème, propriétés, méthodes, applications.}
\end{popfigure}

\courssub{Exercices progressifs et résolus}

Nous passons maintenant à l'entraînement. Les exercices sont classés par niveau de difficulté croissante : identification directe, calculs standards, puis problèmes contextualisés type Probatoire.

\begin{exoresolu}[Exercice 1 : Identification et vérification]
\textbf{Énoncé.} Soit la fonction $f$ définie sur $\mathbb{R}$ par $f(x)=3x^2-4x+5$. 
Parmi les fonctions suivantes, laquelle est une primitive de $f$ sur $\mathbb{R}$ ?
\begin{enumerate}
    \item $F_1(x)=x^3-2x^2+5x$
    \item $F_2(x)=x^3-2x^2+5x+7$
    \item $F_3(x)=x^3-4x^2+5x$
\end{enumerate}
\textbf{Solution détaillée.}
On dérive chaque candidate et on compare à $f(x)$.
\begin{itemize}
    \item $F_1'(x)=3x^2-4x+5 = f(x)$. \textbf{C'est une primitive.}
    \item $F_2'(x)=3x^2-4x+5 = f(x)$. \textbf{C'est aussi une primitive.} (Diffère de $F_1$ d'une constante).
    \item $F_3'(x)=3x^2-8x+5 \neq f(x)$. \textbf{Pas une primitive.}
\end{itemize}
\textbf{Conclusion :} $F_1$ et $F_2$ sont des primitives de $f$. L'ensemble des primitives est $\{x^3-2x^2+5x+C \mid C \in \mathbb{R}\}$.
\end{exoresolu}

\begin{exoresolu}[Exercice 2 : Calcul de primitives usuelles et linéarité]
\textbf{Énoncé.} Déterminer une primitive sur l'intervalle indiqué :
\begin{enumerate}
    \item $f(x)=\dfrac{1}{x^2}$ sur $]0;+\infty[$.
    \item $g(x)=2\sqrt{x}+3\sin x$ sur $[0;+\infty[$.
    \item $h(x)=\dfrac{2x+1}{x^2+x}$ sur $]0;+\infty[$.
\end{enumerate}
\textbf{Solution détaillée.}
\begin{enumerate}
    \item $f(x)=x^{-2}$. Une primitive est $F(x)=\dfrac{x^{-1}}{-1}=-\dfrac{1}{x}$.
    \item $g(x)=2x^{1/2}+3\sin x$. 
    Primitive de $2x^{1/2}$ : $2 \cdot \dfrac{x^{3/2}}{3/2} = \dfrac{4}{3}x^{3/2}$.
    Primitive de $3\sin x$ : $-3\cos x$.
    Par linéarité : $G(x)=\dfrac{4}{3}x^{3/2}-3\cos x$.
    \item On remarque que le numérateur est la dérivée du dénominateur : $(x^2+x)'=2x+1$.
    Donc $h(x)=\dfrac{u'}{u}$ avec $u(x)=x^2+x > 0$ sur $]0;+\infty[$.
    Une primitive est $H(x)=\ln(x^2+x)$.
\end{enumerate}
\end{exoresolu}

\begin{exoresolu}[Exercice 3 : Substitution et problème contextualisé (Type Probatoire)]
\textbf{Énoncé.} 
\begin{enumerate}
    \item Calculer une primitive de la fonction $f(x)=\dfrac{x}{x^2+1}$ sur $\mathbb{R}$.
    \item Un technicien modélise le débit instantané (en L/min) d'une pompe par la fonction $D(t)=\dfrac{2t}{t^2+1}$ pour $t \in [0; 10]$ ($t$ en minutes). 
    Sachant que la cuve contient $50$ L au départ ($t=0$), déterminer le volume d'eau $V(t)$ dans la cuve à l'instant $t$. Calculer le volume à $t=3$ min.
\end{enumerate}
\textbf{Solution détaillée.}
\begin{enumerate}
    \item On pose $u(x)=x^2+1$, alors $u'(x)=2x$. On a $f(x)=\frac{1}{2}\frac{u'(x)}{u(x)}$.
    Une primitive est $F(x)=\frac{1}{2}\ln(u(x))=\frac{1}{2}\ln(x^2+1)$.
    (Vérification : $F'(x)=\frac{1}{2}\frac{2x}{x^2+1}=f(x)$).
    \item Le volume $V(t)$ est une primitive du débit $D(t)$.
    $D(t)=2 \cdot \frac{t}{t^2+1}$. D'après la question 1, une primitive de $\frac{t}{t^2+1}$ est $\frac{1}{2}\ln(t^2+1)$.
    Donc une primitive de $D$ est $V_0(t)=2 \times \frac{1}{2}\ln(t^2+1)=\ln(t^2+1)$.
    La primitive générale est $V(t)=\ln(t^2+1)+C$.
    Condition initiale : $V(0)=50 \Rightarrow \ln(1)+C=50 \Rightarrow 0+C=50 \Rightarrow C=50$.
    Expression finale : $\boxed{V(t)=\ln(t^2+1)+50}$.
    Volume à $t=3$ : $V(3)=\ln(10)+50 \approx 2{,}303+50 = \boxed{52{,}3 \text{ L (arrondi au dixième)}}$.
\end{enumerate}
\end{exoresolu}

\begin{exercice}[Exercice 4 : Intégration par parties (Entraînement)]
\textbf{Énoncé.} Déterminer une primitive sur $\mathbb{R}$ des fonctions suivantes :
\begin{enumerate}
    \item $f(x)=x e^x$
    \item $g(x)=x \cos x$
    \item $h(x)=\ln x$ (sur $]0;+\infty[$) \textit{Indication : écrire $\ln x = 1 \times \ln x$.}
\end{enumerate}
\end{exercice}

\begin{corrige}[Exercice 4]
\begin{enumerate}
    \item $u=x \Rightarrow u'=1$ ; $v'=e^x \Rightarrow v=e^x$. $\int x e^x = x e^x - \int e^x = x e^x - e^x = (x-1)e^x + C$.
    \item $u=x \Rightarrow u'=1$ ; $v'=\cos x \Rightarrow v=\sin x$. $\int x \cos x = x \sin x - \int \sin x = x \sin x + \cos x + C$.
    \item $u=\ln x \Rightarrow u'=1/x$ ; $v'=1 \Rightarrow v=x$. $\int \ln x = x \ln x - \int x \cdot \frac{1}{x} = x \ln x - x + C$.
\end{enumerate}
\end{corrige}

\courssub{Auto-évaluation : Checklist des compétences}

Utilisez cette grille pour valider votre préparation. Cochez les cases en toute honnêteté. Si une case reste non cochée, reportez-vous à la section correspondante du cours ou aux pistes de remédiation ci-dessous.

\begin{aretenir}[Checklist : Suis-je prêt pour le Probatoire ?]
\textbf{1. Savoirs définitoires et théoriques}
\begin{itemize}
    \item [ ] Je sais donner la définition rigoureuse d'une primitive sur un intervalle.
    \item [ ] Je sais énoncer le théorème d'existence des primitives pour une fonction continue.
    \item [ ] Je sais expliquer pourquoi deux primitives d'une même fonction ne diffèrent que d'une constante (théorème des fonctions de dérivée nulle).
    \item [ ] Je connais par cœur la table des primitives usuelles (puissances, $1/x$, $\exp$, $\sin$, $\cos$, $1/(1+x^2)$, $1/\sqrt{1-x^2}$).
\end{itemize}

\textbf{2. Savoir-faire calculatoires}
\begin{itemize}
    \item [ ] J'applique correctement la linéarité : primitive d'une somme / multiple par un scalaire.
    \item [ ] Je reconnais et traite la forme $\frac{u'}{u}$ (logarithme népérien).
    \item [ ] Je reconnais et traite la forme $\frac{u'}{2\sqrt{u}}$ (racine carrée) ou $u' u^n$ (puissance).
    \item [ ] Je sais mettre en œuvre la méthode de substitution (changement de variable) : choix de $u$, calcul de $du$, retour à $x$.
    \item [ ] Je sais mettre en œuvre l'intégration par parties : choix judicieux de $u$ et $v'$ (règle LIATE), calcul de $uv - \int u'v$.
    \item [ ] Je sais déterminer la constante $C$ à l'aide d'une condition initiale (problème de Cauchy).
\end{itemize}

\textbf{3. Compétences transverses (Rédaction, Modélisation, Vérification)}
\begin{itemize}
    \item [ ] Je rédige une solution complète : \emph{« Soit $F$ une primitive de $f$... »}, énoncé des formules utilisées, conclusion encadrée.
    \item [ ] Je sais modéliser un problème concret : identifier la grandeur dérivée (débit, vitesse, coût marginal) et la grandeur primitive (volume, position, coût total).
    \item [ ] Je vérifie systématiquement ma primitive en la dérivant mentalement ou sur brouillon.
    \item [ ] Je gère correctement les intervalles de définition (domaine de $f$, domaine de la primitive, valeurs absolues dans $\ln|u|$ si nécessaire).
\end{itemize}
\end{aretenir}

\courssub{Pistes de remédiation : difficultés récurrentes et erreurs à bannir}

Voici les points de blocage les plus fréquents observés chez les élèves de Terminale technique, avec la parade pédagogique pour chacun.

\begin{attention}[Erreur n°1 : Oubli de la constante $C$ ou de l'intervalle]
\textbf{Le piège :} Écrire « La primitive de $f$ est $F(x)=x^2$ » au lieu de « Les primitives de $f$ sur $\mathbb{R}$ sont les fonctions $F(x)=x^2+C$, $C \in \mathbb{R}$ ».
\textbf{La règle :} Une fonction continue admet \emph{une infinité} de primitives. L'oubli de $C$ ou de la mention de l'intervalle $I$ fait perdre des points (souvent 0,5 à 1 point sur 4). 
\textbf{Astuce :} Systématiser la phrase type : \emph{« L'ensemble des primitives de $f$ sur $I$ est $\{F(x)+C \mid C \in \mathbb{R}\}$ où $F$ est une primitive quelconque. »}
\end{attention}

\begin{attention}[Erreur n°2 : Mauvais choix dans l'intégration par parties (LIATE)]
\textbf{Le piège :} Pour $\int x \ln x \, dx$, poser $u=x$ et $v'=\ln x$ (impossible à intégrer directement). Ou pour $\int x e^x dx$, poser $u=e^x$.
\textbf{La règle LIATE (ordre de priorité pour choisir $u$) :}
\textbf{L}ogarithme ($\ln x$) $\to$ \textbf{I}nverse trigonométrique ($\arcsin$) $\to$ \textbf{A}lgebraique ($x, x^2$) $\to$ \textbf{T}rigonométrique ($\sin, \cos$) $\to$ \textbf{E}xponentielle ($e^x$).
\textbf{Exemple :} $\int x \ln x \, dx \Rightarrow u=\ln x$ (L avant A), $v'=x$.
\end{attention}

\begin{attention}[Erreur n°3 : Substitution incomplète (oubli du retour à $x$ ou du $dx$)]
\textbf{Le piège :} Poser $u=x^2+1$, écrire $\int \frac{1}{u} du = \ln u$ et s'arrêter là sans remplacer $u$ par $x^2+1$. Ou oublier de diviser par la dérivée intérieure (ex: $\int \cos(2x) dx = \sin(2x)$ au lieu de $\frac{1}{2}\sin(2x)$).
\textbf{La méthode sécurisée (3 étapes) :}
\begin{enumerate}
    \item Poser $u=g(x)$, calculer $du=g'(x)dx$, exprimer $dx$ (ou $x\,dx$).
    \item Transformer \textbf{intégralement} l'intégrale en variable $u$ (plus de $x$).
    \item Intégrer en $u$, puis \textbf{remplacer $u$ par $g(x)$} pour redonner le résultat en $x$.
\end{enumerate}
\end{attention}

\begin{attention}[Erreur n°4 : Domaine de définition et valeurs absolues]
\textbf{Le piège :} Écrire $\int \frac{1}{x} dx = \ln x + C$ sur $\mathbb{R}^*$ (faux car $\ln x$ n'est pas défini pour $x<0$).
\textbf{La rigueur :} $\int \frac{1}{x} dx = \ln|x| + C$ sur $\mathbb{R}^*$ (deux intervalles distincts : $]-\infty;0[$ et $]0;+\infty[$). Au Probatoire, l'intervalle est souvent donné (ex: $]0;+\infty[$), ce qui permet d'écrire $\ln x$ sans valeur absolue. \textbf{Lisez bien l'intervalle de l'énoncé.}
\end{attention}

\courssub{Préparation au Probatoire : types de questions attendues}

L'épreuve du Probatoire (Mathématiques, Enseignement Technique) teste la maîtrise des techniques standards et la capacité à les mobiliser dans un contexte technique. Voici les 3 profils de questions types.

\begin{exemplebox}[Exemple type Probatoire : Question « Classique » (Calcul pur)]
\textbf{Sujet :} Déterminer une primitive de la fonction $f$ définie par $f(x)=\dfrac{x^2+1}{x^3+3x}$ sur l'intervalle $]0;+\infty[$.
\textbf{Analyse :} 
\begin{itemize}
    \item Dénominateur $u(x)=x^3+3x$. Dérivée $u'(x)=3x^2+3 = 3(x^2+1)$.
    \item Numérateur $x^2+1 = \frac{1}{3}u'(x)$.
    \item Forme $\frac{u'}{u}$ identifiée avec un facteur constant.
\end{itemize}
\textbf{Rédaction attendue :}
Soit $u(x)=x^3+3x$. $u$ est dérivable et strictement positive sur $]0;+\infty[$.
$u'(x)=3x^2+3 = 3(x^2+1)$. Donc $f(x)=\frac{1}{3}\frac{u'(x)}{u(x)}$.
Une primitive de $f$ sur $]0;+\infty[$ est $F(x)=\frac{1}{3}\ln(u(x)) = \frac{1}{3}\ln(x^3+3x)$.
\textbf{Vérification :} $F'(x)=\frac{1}{3}\frac{3x^2+3}{x^3+3x}=\frac{x^2+1}{x^3+3x}=f(x)$.
\end{exemplebox}

\begin{exemplebox}[Exemple type Probatoire : Question « Condition initiale » (Problème de Cauchy)]
\textbf{Sujet :} La vitesse d'un mobile (en m/s) est donnée par $v(t)=3t^2-12t+9$ pour $t \in [0;5]$. 
Sachant que sa position initiale ($t=0$) est $x_0=5$ m, déterminer l'équation horaire $x(t)$ et l'instant où le mobile change de sens.
\textbf{Attendus :}
\begin{enumerate}
    \item Primitive de $v$ : $x(t)=t^3-6t^2+9t+C$.
    \item Condition $x(0)=5 \Rightarrow C=5$. Donc $x(t)=t^3-6t^2+9t+5$.
    \item Changement de sens $\Leftrightarrow v(t)=0 \Rightarrow 3(t^2-4t+3)=0 \Rightarrow (t-1)(t-3)=0 \Rightarrow t=1$ s ou $t=3$ s.
    \item Étude du signe de $v$ pour conclure : $v>0$ sur $[0;1[$, $v<0$ sur $]1;3[$, $v>0$ sur $]3;5]$.
\end{enumerate}
\end{exemplebox}

\begin{exemplebox}[Exemple type Probatoire : Question « Modélisation / Choix de méthode »]
\textbf{Sujet :} Le coût marginal (en milliers FCFA/unité) de production de $x$ pièces est $C_m(x)=0{,}02x+5-\frac{100}{x+1}$ pour $x \in [0;100]$. Le coût fixe est de $200$ milliers FCFA.
\begin{enumerate}
    \item Exprimer le coût total $C(x)$.
    \item Déterminer le coût de production de $50$ pièces.
\end{enumerate}
\textbf{Attendus :}
\begin{enumerate}
    \item $C(x)$ primitive de $C_m$. $C(x)=0{,}01x^2+5x-100\ln(x+1)+K$.
    \item $C(0)=200 \Rightarrow 0+0-100\ln(1)+K=200 \Rightarrow K=200$.
    \item $C(50)=0{,}01(2500)+250-100\ln(51)+200 = 25+250-100(3{,}93)+200 \approx 475-393 = 82$ milliers FCFA.
\end{enumerate}
\end{exemplebox}

\courssub{Le savais-tu ? : Les primitives, cœur de l'ingénierie moderne}

\begin{savaistu}[Au-delà du Probatoire : l'intégrale au service de la technique]
Les primitives que vous manipulez sont la face « algébrique » du \textbf{calcul intégral}, l'un des piliers des mathématiques appliquées.
\begin{itemize}
    \item \textbf{Génie civil / Mécanique :} Calculer le centre de gravité d'une poutre à section variable, la flèche d'une poutre sous charge (intégration double de l'équation de la déformée $EI y'' = M(x)$), le moment d'inertie de sections complexes.
    \item \textbf{Électrotechnique :} Déterminer la charge d'un condensateur $q(t)=\int i(t)dt$, l'énergie stockée $W=\frac{1}{2}\int L i^2 dt$, la réponse temporelle des circuits RLC (équations différentielles $\to$ primitives).
    \item \textbf{Économie / Gestion :} Reconstruire la fonction Coût Total à partir du Coût Marginal, la fonction Revenu à partir du Revenu Marginal, calculer le surplus du consommateur/producteur (aire sous courbe).
    \item \textbf{Statistiques / Qualité :} La loi normale (Gaussienne) repose sur l'intégrale $\int e^{-x^2}dx$ (fonction erreur), base des intervalles de confiance et des cartes de contrôle en industrie.
\end{itemize}
\textbf{En résumé :} Maîtriser la recherche de primitives, c'est s'ouvrir les portes de la modélisation quantitative de tout système physique, économique ou biologique évoluant dans le temps ou l'espace.
\end{savaistu}

\begin{methode}[Méthode « Zéro Faute » pour le jour de l'examen (Checklist rapide)]
\begin{enumerate}
    \item \textbf{Lire} l'énoncé 2 fois : intervalle $I$ ? Condition initiale ? Forme demandée (primitive générale / particulière) ?
    \item \textbf{Identifier} la structure : Somme ? $\frac{u'}{u}$ ? $u'u^n$ ? Produit $u v'$ ? Fraction rationnelle ?
    \item \textbf{Poser} l'intention : « Je cherche une primitive de $f$ sur $I$... » / « Je pose $u=...$ ».
    \item \textbf{Calculer} proprement, en alignant les égalités.
    \item \textbf{Vérifier} en dérivant mentalement le résultat obtenu (coût 10 secondes, gain : élimination des erreurs de signe/coefficient).
    \item \textbf{Conclure} : Encadrer la réponse finale, mentionner $+C$ et l'intervalle $I$. Si condition initiale : donner la fonction particulière (sans $C$).
\end{enumerate}
\end{methode}

\newpage

\coursec{Activité d'intégration}

\coursec{Primitives d'une fonction continue sur un intervalle}

\courssub{Objectifs de la leçon}

À l'issue de cette leçon, l'élève sera capable de :
\begin{itemize}
\item exploiter la \cle{définition d'une primitive} d'une fonction continue sur un intervalle ;
\item mobiliser le \cle{théorème d'existence} des primitives pour justifier la démarche ;
\item déterminer une primitive d'une fonction polynôme et d'une fonction affine par morceaux ;
\item interpréter concrètement une primitive : la distance parcourue comme primitive de la vitesse ;
\item rédiger et justifier une démarche complète, une réponse et une conclusion dans une situation de la vie courante (transport urbain, consommation de carburant).
\end{itemize}

\courssub{Rappels : dérivée et notion de primitive}

\begin{definition}[Primitive d'une fonction]
Soit $f$ une fonction définie sur un intervalle $I$. Une fonction $F$ dérivable sur $I$ est une \textbf{primitive} de $f$ sur $I$ si et seulement si :
\[ \forall x \in I,\quad F'(x) = f(x). \]
\end{definition}

\begin{propriete}[Unicité à constante près]
Si $F$ est une primitive de $f$ sur $I$, alors toutes les primitives de $f$ sur $I$ sont les fonctions de la forme $F + k$, avec $k \in \mathbb{R}$.
\end{propriete}

\begin{aretenir}[Notation]
On note souvent $\int f(x)\,dx$ l'ensemble des primitives de $f$. Une primitive particulière sera notée $F(x)$.
\end{aretenir}

\courssub{Théorème d'existence des primitives pour les fonctions continues}

\begin{propriete}[Théorème d'existence -- Exigible au Probatoire]
Toute fonction \cle{continue} sur un intervalle $I$ admet des primitives sur $I$.
\end{propriete}

\begin{attention}[Condition d'application]
La continuité sur un \textbf{intervalle} est une condition \textbf{suffisante} (mais non nécessaire). Une fonction discontinue peut avoir des primitives (ex. : fonction dérivée d'une fonction dérivable mais pas $C^1$), mais le programme de Terminale technique se concentre sur le cas continu.
\end{attention}

\courssub{Propriétés algébriques des primitives}

\begin{propriete}[Linéarité]
Soient $f$ et $g$ deux fonctions admettant des primitives sur $I$, et $\lambda \in \mathbb{R}$.
\begin{itemize}
\item Une primitive de $f+g$ est la somme d'une primitive de $f$ et d'une primitive de $g$.
\item Une primitive de $\lambda f$ est $\lambda$ fois une primitive de $f$.
\end{itemize}
\end{propriete}

\begin{aretenir}[Tableau de primitives usuelles (à connaître)]
\begin{center}
\begin{tabular}{|c|c|c|}\hline
$f(x)$ & Primitive $F(x)$ & Condition \\ \hline
$k$ (constante) & $kx$ & $k \in \mathbb{R}$ \\ \hline
$x^n$ ($n \in \mathbb{N}$) & $\dfrac{x^{n+1}}{n+1}$ & $n \neq -1$ \\ \hline
$\dfrac{1}{x}$ & $\ln|x|$ & $x \neq 0$ (hors programme Tle Tech) \\ \hline
\end{tabular}
\end{center}
\end{aretenir}

\courssub{Détermination de primitives : méthode et exercices types}

\begin{methode}[{Recherche d'une primitive d'une fonction polynôme sur $[a;b]$}]
\begin{enumerate}
\item Vérifier que la fonction est continue sur l'intervalle (polynôme $\Rightarrow$ continue sur $\mathbb{R}$).
\item Appliquer la linéarité : primitive de la somme = somme des primitives.
\item Pour chaque terme $a_n x^n$, écrire la primitive $\dfrac{a_n}{n+1}x^{n+1}$.
\item Ajouter une constante $k$.
\item Si une condition initiale $F(x_0)=y_0$ est donnée, déterminer $k$.
\item \textbf{Vérifier systématiquement} en dérivant : $F'(x) \stackrel{?}{=} f(x)$.
\end{enumerate}
\end{methode}

\begin{exoresolu}[Primitive d'un polynôme]
Déterminer la primitive $F$ de $f(x)=3x^2-4x+5$ sur $\mathbb{R}$ telle que $F(0)=2$.
\tcblower
$f$ est un polynôme, donc continue sur $\mathbb{R}$, elle admet des primitives.
\begin{align*}
F(x) &= 3\times\frac{x^3}{3} - 4\times\frac{x^2}{2} + 5x + k \\
&= x^3 - 2x^2 + 5x + k.
\end{align*}
$F(0)=2 \Rightarrow k=2$. Donc $F(x)=x^3-2x^2+5x+2$.
\textbf{Vérification :} $F'(x)=3x^2-4x+5=f(x)$. \checkmark
\end{exoresolu}

\begin{exoresolu}[Fonction affine par morceaux]
Soit $f$ définie sur $[0;3]$ par $f(x) = \begin{cases} 2x & \text{si } 0 \le x < 1 \\ 2 & \text{si } 1 \le x \le 3 \end{cases}$.
Déterminer la primitive $F$ telle que $F(0)=0$.
\tcblower
$f$ est continue sur $[0;1[$ et sur $]1;3]$, avec une discontinuité au point $x=1$ (saut). \textbf{Attention :} le théorème d'existence ne s'applique pas sur $[0;3]$ car $f$ n'est pas continue sur l'intervalle \textbf{fermé}. Cependant, on peut chercher une primitive \textbf{par morceaux} continue sur $[0;3]$.
\begin{itemize}
\item Sur $[0;1]$ : $F(x)=x^2+k_1$. $F(0)=0 \Rightarrow k_1=0$, donc $F(x)=x^2$.
\item Pour que $F$ soit continue en $1$, il faut $F(1)=1$.
\item Sur $[1;3]$ : $F(x)=2x+k_2$. $F(1)=1 \Rightarrow 2+k_2=1 \Rightarrow k_2=-1$.
\end{itemize}
Donc $F(x) = \begin{cases} x^2 & \text{si } 0 \le x \le 1 \\ 2x-1 & \text{si } 1 \le x \le 3 \end{cases}$.
Vérification : $F$ est continue sur $[0;3]$, dérivable sur $[0;1[$ et $]1;3]$, et $F'=f$ sur ces intervalles.
\end{exoresolu}

\courssub{Applications concrètes : modélisation de situations réelles}

\begin{experience}[Activité d'intégration : Le taxi-moto de Monsieur Etoundi]
\textbf{Situation.}
Monsieur Etoundi conduit un \emph{taxi-moto} à Douala, entre le rond-point Deïdo et le marché Mboppi. Pour préparer son budget hebdomadaire, il veut connaître la distance parcourue et le carburant consommé pendant ses deux heures de pointe du matin (6 h à 8 h).

Son compteur kilométrique étant en panne, son neveu relève la \cle{vitesse moyenne} $v(t)$ (en km/h) toutes les 15 minutes. Il obtient le tableau suivant, où $t$ est le temps écoulé en heures depuis 6 h :

\begin{center}
\begin{tabularx}{\linewidth}{|l|X|X|X|X|X|X|X|X|X|}
\hline
\rowcolor{popblueL}
$t$ (h) & $0$ & $0,25$ & $0,5$ & $0,75$ & $1$ & $1,25$ & $1,5$ & $1,75$ & $2$ \\
\hline
$v(t)$ (km/h) & $10$ & $25$ & $40$ & $45$ & $40$ & $30$ & $20$ & $15$ & $10$ \\
\hline
\end{tabularx}
\end{center}

Le taxi-moto consomme en moyenne \SI{2,5}{L} d'essence aux 100 km. Le litre d'essence coûte 730 FCFA.

\textbf{Problème :} quelle distance totale le taxi-moto a-t-il parcourue entre 6 h et 8 h ? Quelle quantité d'essence a-t-il consommée, et quel est le coût ?
\end{experience}

\courssub{Consignes de l'activité}

Travail en groupes de 3 ou 4 élèves. Chaque groupe rédige un rapport et le présente oralement en 5 minutes.

\begin{enumerate}
\item Tracer sur papier millimétré le nuage de points $(t\,;v(t))$. Échelle : 4 cm pour 1 h en abscisse, 1 cm pour 5 km/h en ordonnée.
\item Expliquer pourquoi on peut modéliser la vitesse par une fonction \cle{continue} sur $[0\,;2]$, et citer le théorème qui garantit l'existence de primitives.
\item On propose de modéliser la vitesse sur $[0\,;2]$ par la fonction polynôme qui passe par les points $(0;10)$, $(1;40)$ et $(2;10)$ :
\[ v(t) = -30t^2 + 60t + 10. \]
Vérifier que $v(0)=10$, $v(1)=40$ et $v(2)=10$. Comparer $v(0,5)$, $v(0,75)$, $v(1,5)$ avec les valeurs du tableau. Le modèle vous semble-t-il acceptable ?
\item Justifier que $v$ admet des primitives sur $[0\,;2]$, puis déterminer la primitive $D$ de $v$ qui vérifie $D(0)=0$. Que représente concrètement $D(t)$ ?
\item Calculer la distance totale parcourue entre 6 h et 8 h, c'est-à-dire $D(2)$.
\item Calculer la quantité d'essence consommée et la dépense en FCFA.
\item Rédiger une conclusion complète à l'attention de Monsieur Etoundi : distance, consommation, coût, et une remarque critique sur la fiabilité du modèle.
\end{enumerate}

\courssub{Ressources à mobiliser}

\begin{aretenir}[Boîte à outils de la leçon]
\begin{itemize}
\item \textbf{Définition :} $F$ est une primitive de $f$ sur un intervalle $I$ si $F$ est dérivable sur $I$ et $F'(x)=f(x)$ pour tout $x \in I$.
\item \textbf{Théorème d'existence :} toute fonction \cle{continue} sur un intervalle $I$ admet des primitives sur $I$.
\item \textbf{Unicité à constante près :} si $F$ est une primitive de $f$, toutes les primitives sont $F+k$, $k\in\mathbb{R}$. La condition $F(x_0)=y_0$ fixe une unique primitive.
\item \textbf{Primitives usuelles :} primitive de $x^n$ ($n\in\mathbb{N}$) est $\dfrac{x^{n+1}}{n+1}$ ; primitive de $a$ (constante) est $ax$.
\item \textbf{Linéarité :} primitive d'une somme = somme des primitives ; primitive de $\lambda f$ = $\lambda \times$ primitive de $f$.
\item \textbf{Lien physique :} la distance parcourue est une primitive de la vitesse par rapport au temps.
\end{itemize}
\end{aretenir}

\courssub{Démarche et corrigé commenté}

\begin{exoresolu}[Consommation de carburant d'un taxi-moto -- Corrigé détaillé]
Énoncé : à partir du relevé de vitesses et du modèle $v(t)=-30t^2+60t+10$ sur $[0\,;2]$, déterminer la distance parcourue, la consommation d'essence et le coût.
\tcblower
\textbf{Étape 1 : représentation graphique.}
On place les neuf points : $(0;10)$, $(0,25;25)$, $(0,5;40)$, $(0,75;45)$, $(1;40)$, $(1,25;30)$, $(1,5;20)$, $(1,75;15)$, $(2;10)$. Le nuage a une forme de « bosse » : la vitesse augmente (fluidité), atteint un maximum vers $t=0,75$ h, puis diminue (embouteillages).

\textbf{Étape 2 : continuité et existence des primitives.}
La vitesse d'un véhicule varie de façon continue (pas de saut instantané). Il est légitime de modéliser $v$ par une fonction continue sur $[0\,;2]$. D'après le \cle{théorème d'existence}, toute fonction continue sur un intervalle admet des primitives : la recherche d'une primitive de $v$ a un sens.

\textbf{Étape 3 : validation du modèle polynomial.}
Pour $v(t)=-30t^2+60t+10$ :
\begin{align*}
v(0) &= 10 \quad\text{(conforme)}\\
v(1) &= -30+60+10 = 40 \quad\text{(conforme)}\\
v(2) &= -30(4)+60(2)+10 = -120+120+10 = 10 \quad\text{(conforme)}.
\end{align*}
Comparaison aux autres points :
\begin{align*}
v(0,5) &= -30(0,25)+60(0,5)+10 = -7,5+30+10 = 32,5 \quad\text{(tableau : 40)}\\
v(0,75) &= -30(0,5625)+60(0,75)+10 = -16,875+45+10 = 38,125 \quad\text{(tableau : 45)}\\
v(1,5) &= -30(2,25)+60(1,5)+10 = -67,5+90+10 = 32,5 \quad\text{(tableau : 20)}.
\end{align*}
Le modèle passe exactement par les trois points choisis mais s'écarte notablement des autres mesures (surtout après 1 h). Il donne une \textbf{estimation approchée} de la distance. Un modèle affine par morceaux (méthode des trapèzes) serait plus fidèle aux données brutes.

\textbf{Étape 4 : détermination de la primitive $D$ telle que $D(0)=0$.}
$v$ est un polynôme, donc continue sur $[0\,;2]$, elle admet des primitives.
\begin{align*}
D(t) &= -30\times\frac{t^3}{3} + 60\times\frac{t^2}{2} + 10t + k \\
&= -10t^3 + 30t^2 + 10t + k.
\end{align*}
$D(0)=0 \Rightarrow k=0$. Donc :
\[ D(t) = -10t^3 + 30t^2 + 10t. \]
\textbf{Vérification (réflexe indispensable) :} $D'(t) = -30t^2+60t+10 = v(t)$. \checkmark
$D(t)$ représente la \textbf{distance parcourue en km} depuis 6 h.

\textbf{Étape 5 : distance totale $D(2)$.}
\begin{align*}
D(2) &= -10(8) + 30(4) + 10(2) \\
&= -80 + 120 + 20 = 60 \text{ km}.
\end{align*}
La distance totale parcourue est \textbf{60 km} (estimation par le modèle polynomial).

\textbf{Étape 6 : consommation et coût.}
Consommation : $\dfrac{2,5}{100} \times 60 = 1,5$ L.
Coût : $1,5 \times 730 = 1\,095$ FCFA.

\textbf{Étape 7 : conclusion rédigée.}
Entre 6 h et 8 h, le taxi-moto de Monsieur Etoundi a parcouru \textbf{environ 60 km} (selon le modèle polynomial), consommé \textbf{1,5 L} d'essence, pour une dépense de \textbf{1 095 FCFA}.
\textbf{Remarque critique :} le modèle polynomial lisse les variations réelles de la vitesse. Une estimation par la méthode des trapèzes sur les données brutes donne une distance d'environ 58,75 km (soit 1,47 L, 1 073 FCFA). L'erreur du modèle polynomial est d'environ 2 \%, acceptable pour un budget prévisionnel hebdomadaire.
\end{exoresolu}

\begin{attention}[Erreurs fréquentes à éviter]
\begin{itemize}
\item Oublier de \textbf{vérifier} que $D'(t)=v(t)$ après avoir calculé une primitive : c'est la seule garantie du résultat.
\item Se tromper dans la primitive de $t^2$ : c'est $\dfrac{t^3}{3}$ et non $3t^3$ ni $2t$.
\item Confondre les unités : $t$ est en \textbf{heures}, $v$ en km/h, donc $D$ est bien en km. Ne pas convertir 15 min en 15 (il faut $0,25$ h).
\item Oublier la condition initiale $D(0)=0$ : sans elle, la distance n'est pas déterminée de manière unique.
\item Appliquer le théorème d'existence à une fonction non continue sur l'intervalle \textbf{fermé} sans le signaler.
\end{itemize}
\end{attention}

\courssub{Bilan de la leçon}

Cette leçon a permis de mettre en évidence que :
\begin{itemize}
\item la notion de primitive n'est pas qu'un « calcul inverse de la dérivée » : elle possède un \textbf{sens concret} (distance à partir de la vitesse, quantité accumulée à partir d'un débit) ;
\item le \cle{théorème d'existence} joue un rôle clé : c'est la continuité du phénomène physique qui justifie la démarche de modélisation ;
\item la condition initiale $D(0)=0$ lève l'indétermination « à constante près » et donne une réponse unique ;
\item les mathématiques de la Terminale technique permettent de résoudre un problème budgétaire réel, en combinant lecture de tableau, modélisation, calcul de primitives et interprétation critique du résultat.
\end{itemize}

La restitution orale des groupes permet de comparer les modèles proposés (polynôme, affine par morceaux, méthode des trapèzes) et de discuter la robustesse des estimations obtenues.

\begin{savaistu}[Le saviez-vous ?]
Le lien entre vitesse et distance est au cœur du \textbf{calcul intégral}. Newton et Leibniz l'ont formalisé au XVII\textsuperscript{e} siècle. Aujourd'hui, les GPS de nos téléphones utilisent des méthodes numériques (similaires à la méthode des trapèzes) pour calculer la distance parcourue à partir des mesures de vitesse instantanée. Au Cameroun, les gestionnaires de flottes de transport (SOTUC, taxis, moto-taxis) utilisent ces principes pour optimiser la consommation de carburant et planifier la maintenance.
\end{savaistu}

\begin{exercice}[Entraînement -- Primitive et condition initiale]
Soit $f(x) = 4x^3 - 6x + 2$ sur $\mathbb{R}$.
\begin{enumerate}
\item Déterminer l'ensemble des primitives de $f$ sur $\mathbb{R}$.
\item Déterminer la primitive $F$ telle que $F(1)=5$.
\item Vérifier votre résultat.
\end{enumerate}
\end{exercice}

\begin{corrige}[Exercice 1]
\begin{enumerate}
\item $F(x) = x^4 - 3x^2 + 2x + k,\quad k \in \mathbb{R}$.
\item $F(1)=5 \Rightarrow 1 - 3 + 2 + k = 5 \Rightarrow k = 5$. Donc $F(x)=x^4-3x^2+2x+5$.
\item $F'(x)=4x^3-6x+2=f(x)$. \checkmark
\end{enumerate}
\end{corrige}

\begin{exercice}[Application -- Modélisation affine par morceaux]
Un artisan soudeur utilise une machine dont la puissance électrique $P(t)$ (en kW) varie ainsi pendant 4 heures :
\[ P(t) = \begin{cases} 2t & \text{si } 0 \le t < 1 \\ 2 & \text{si } 1 \le t < 3 \\ 8-2t & \text{si } 3 \le t \le 4 \end{cases} \]
L'énergie consommée $E(t)$ (en kWh) est une primitive de $P(t)$ avec $E(0)=0$.
\begin{enumerate}
\item Tracer la courbe de $P$ sur $[0;4]$.
\item Déterminer $E(t)$ sur chaque intervalle en assurant la continuité de $E$ sur $[0;4]$.
\item Calculer l'énergie totale consommée en 4 heures.
\item Si le kWh coûte 50 FCFA, quel est le coût de l'électricité ?
\end{enumerate}
\end{exercice}

\begin{corrige}[Exercice 2]
\begin{enumerate}
\item Courbe : triangle + rectangle + triangle.
\item Sur $[0;1]$ : $E(t)=t^2$ ($E(0)=0$). $E(1)=1$.
Sur $[1;3]$ : $E(t)=2t+k$. Continuité en 1 : $2+k=1 \Rightarrow k=-1$. $E(t)=2t-1$. $E(3)=5$.
Sur $[3;4]$ : $E(t)=8t-t^2+k$. Continuité en 3 : $24-9+k=5 \Rightarrow k=-10$. $E(t)=8t-t^2-10$.
\item $E(4)=32-16-10=6$ kWh.
\item Coût = $6 \times 50 = 300$ FCFA.
\end{enumerate}
\end{corrige}

\newpage

\coursec{Méthodes et exercices résolus}

\coursec{Rappels : Dérivée et notion de primitive}

\begin{definition}[Dérivée d'une fonction]
Soit $f$ une fonction définie sur un intervalle $I$ et $a \in I$. Si le taux d'accroissement $\frac{f(x)-f(a)}{x-a}$ admet une limite finie lorsque $x \to a$, on dit que $f$ est \textbf{dérivable} en $a$. Cette limite est appelée \textbf{dérivée de $f$ en $a$} et se note $f'(a)$.
\[ f'(a) = \lim_{x \to a} \frac{f(x)-f(a)}{x-a} = \lim_{h \to 0} \frac{f(a+h)-f(a)}{h} \]
Géométriquement, $f'(a)$ est le coefficient directeur de la tangente à la courbe $C_f$ au point d'abscisse $a$.
\end{definition}

\begin{definition}[Primitive d'une fonction]
Soit $f$ une fonction définie sur un intervalle $I$. On appelle \textbf{primitive de $f$ sur $I$} toute fonction $F$ dérivable sur $I$ telle que :
\[ \forall x \in I,\quad F'(x) = f(x) \]
\cle{Notation :} L'ensemble des primitives de $f$ sur $I$ se note parfois $\int f(x)\,dx$ (intégrale indéfinie).
\end{definition}

\begin{propriete}[Lien Dérivée -- Primitive]
$F$ est une primitive de $f$ sur $I$ $\iff$ $F$ est une antécédente de $f$ par l'application dérivation.
\cle{Conséquence immédiate :} Si $F$ est une primitive de $f$ sur $I$, alors pour toute constante $C \in \mathbb{R}$, la fonction $G(x) = F(x) + C$ est aussi une primitive de $f$ sur $I$.
\end{propriete}

\begin{popfigure}
\begin{tikzpicture}[scale=0.8, >=Stealth]
\begin{axis}[
    axis lines=middle, xlabel=$x$, ylabel=$y$,
    xmin=-0.5, xmax=4, ymin=-0.5, ymax=3.5,
    xtick={1,2,3}, ytick={1,2,3},
    tick label style={font=\small},
    clip=false
]
\addplot[popcurve, domain=0:3.5, samples=100] {0.5*x^2} node[right, pos=0.9] {$C_f : y = F(x)$};
\addplot[poporange, domain=0:3.5, samples=100] {x} node[right, pos=0.9] {$C_{f'} : y = f(x)$};
\draw[dashed, popmuted] (2,0) node[below] {$a$} -- (2,2) node[left] {$F(a)$};
\draw[dashed, popmuted] (2,2) -- (2,4) node[right] {$f(a)$};
\draw[popblue, thick] (1,1) -- (3,3) node[above right] {Tangente en $a$ (pente $f(a)$)};
\fill[popblue] (2,2) circle (2pt);
\fill[poporange] (2,2) circle (2pt); % same point for F
\fill[poporange] (2,2) circle (2pt); % f(a) is 2
\end{axis}
\end{tikzpicture}
\legende{Relation graphique : $F$ primitive de $f \iff$ pente de la tangente à $C_F$ en $a$ vaut $f(a)$.}
\end{popfigure}

\coursec{Théorème d'existence des primitives pour les fonctions continues}

\begin{propriete}[Théorème d'existence -- Exigible au Probatoire]
\textbf{Théorème :} Toute fonction continue sur un intervalle $I$ admet au moins une primitive sur $I$.
\par
\textbf{Théorème (Unicité à constante près) :} Si $F$ et $G$ sont deux primitives d'une même fonction $f$ sur un intervalle $I$, alors il existe une constante $C \in \mathbb{R}$ telle que $G(x) = F(x) + C$ pour tout $x \in I$.
\par
\cle{En pratique :} On note $\int f(x)\,dx = F(x) + C$ pour désigner l'ensemble de toutes les primitives de $f$ sur $I$.
\end{propriete}

\begin{aretenir}[Condition de continuité]
L'hypothèse de continuité sur un \textbf{intervalle} est essentielle.
\begin{itemize}
    \item $f(x) = \frac{1}{x}$ est continue sur $]-\infty,0[$ et sur $]0,+\infty[$ (deux intervalles distincts). Ses primitives sont $\ln|x| + C_1$ sur $]-\infty,0[$ et $\ln|x| + C_2$ sur $]0,+\infty[$ ($C_1$ et $C_2$ peuvent être différents).
    \item Une fonction à discontinuité (ex: partie entière) n'admet pas de primitive sur $\mathbb{R}$.
\end{itemize}
\end{aretenir}

\coursec{Propriétés algébriques des primitives}

\begin{propriete}[Linéarité de la primitive]
Soient $f$ et $g$ deux fonctions continues sur un intervalle $I$, et $\lambda \in \mathbb{R}$.
\begin{enumerate}
    \item \textbf{Somme :} $\int \bigl(f(x) + g(x)\bigr)\,dx = \int f(x)\,dx + \int g(x)\,dx$.
    \item \textbf{Multiple par un scalaire :} $\int \lambda f(x)\,dx = \lambda \int f(x)\,dx$.
\end{enumerate}
\cle{Attention :} On n'ajoute \textbf{qu'une seule constante $+C$} à la fin du calcul global.
\end{propriete}

\begin{aretenir}[Table de primitives usuelles (à connaître par cœur)]
Valables sur tout intervalle où les fonctions sont définies et continues ($C \in \mathbb{R}$).
\begin{center}
\begin{tabular}{|c|c|c|}\hline
$f(x)$ & Domaine $I$ & $\int f(x)\,dx = F(x) + C$ \\ \hline
$x^\alpha\ (\alpha \in \mathbb{Q},\ \alpha \neq -1)$ & Selon $\alpha$ & $\dfrac{x^{\alpha+1}}{\alpha+1}$ \\ \hline
$\dfrac{1}{x}$ & $\mathbb{R}^*$ & $\ln|x|$ \\ \hline
$e^x$ & $\mathbb{R}$ & $e^x$ \\ \hline
$\cos x$ & $\mathbb{R}$ & $\sin x$ \\ \hline
$\sin x$ & $\mathbb{R}$ & $-\cos x$ \\ \hline
$\dfrac{1}{1+x^2}$ & $\mathbb{R}$ & $\arctan x$ \\ \hline
$\dfrac{1}{\sqrt{1-x^2}}$ & $]-1,1[$ & $\arcsin x$ \\ \hline
\end{tabular}
\end{center}
\cle{Rappel :} $\sqrt{x} = x^{1/2}$, $\frac{1}{\sqrt{x}} = x^{-1/2}$, $\frac{1}{x^n} = x^{-n}$.
\end{aretenir}

\coursec{Détermination de primitives : méthodes et exercices types}

\begin{methode}[Reconnaître et écrire une primitive usuelle]
\emph{Objectif :} Retrouver une primitive $F$ d'une fonction $f$ en utilisant la table de référence.
\begin{enumerate}
    \item Identifier la forme de $f(x)$ : $x^\alpha$, $\frac{1}{x}$, $\sqrt{x}$, $\frac{1}{\sqrt{x}}$, $e^x$, $\ln x$, $\cos x$, $\sin x$, etc.
    \item Appliquer la formule inverse de la dérivation :
    \begin{itemize}
        \item $f(x)=x^\alpha\ (\alpha \neq -1) \implies F(x)=\frac{x^{\alpha+1}}{\alpha+1}$
        \item $f(x)=\frac{1}{x} \implies F(x)=\ln|x|$
        \item $f(x)=e^x \implies F(x)=e^x$
        \item $f(x)=\cos x \implies F(x)=\sin x$
        \item $f(x)=\sin x \implies F(x)=-\cos x$
    \end{itemize}
    \item \textbf{Nécessaire :} Ajouter la constante arbitraire $+C$ ($C \in \mathbb{R}$).
    \item Vérifier en dérivant $F$ pour retrouver $f$.
\end{enumerate}
\cle{Réflexe :} Toujours écrire $+C$ à la fin. Pour $\frac{1}{x}$, écrire $\ln|x|$ (valeurs absolues) car le domaine peut être $\mathbb{R}^*$.
\end{methode}

\begin{exoresolu}[Primitive d'une fonction polynôme et rationnelle simple]
\emph{Énoncé :} Déterminer l'ensemble des primitives sur l'intervalle de définition de la fonction $f$ définie par $f(x)=3x^2 - \frac{4}{x^2} + 5\sqrt{x}$.
\tcblower
\emph{Solution détaillée :}
La fonction $f$ est définie et continue sur $]0, +\infty[$ (car $\sqrt{x}$ impose $x \ge 0$ et $\frac{1}{x^2}$ impose $x \neq 0$). On cherche $F$ telle que $F'=f$ sur $]0, +\infty[$. On utilise la linéarité et la table des primitives usuelles.
\begin{itemize}
    \item Terme $3x^2$ : forme $x^\alpha$ avec $\alpha=2$. Primitive : $3 \times \frac{x^{3}}{3} = x^3$.
    \item Terme $-\frac{4}{x^2} = -4x^{-2}$ : forme $x^\alpha$ avec $\alpha=-2 \neq -1$. Primitive : $-4 \times \frac{x^{-1}}{-1} = \frac{4}{x}$.
    \item Terme $5\sqrt{x} = 5x^{1/2}$ : forme $x^\alpha$ avec $\alpha=1/2$. Primitive : $5 \times \frac{x^{3/2}}{3/2} = 5 \times \frac{2}{3}x^{3/2} = \frac{10}{3}x\sqrt{x}$.
\end{itemize}
En sommant et en ajoutant la constante $C$ :
\[ F(x) = x^3 + \frac{4}{x} + \frac{10}{3}x\sqrt{x} + C \quad (C \in \mathbb{R}) \]
\emph{Vérification :} $F'(x) = 3x^2 - \frac{4}{x^2} + \frac{10}{3}\bigl(\sqrt{x} + x \cdot \frac{1}{2\sqrt{x}}\bigr) = 3x^2 - \frac{4}{x^2} + \frac{10}{3}\frac{3\sqrt{x}}{2} = 3x^2 - \frac{4}{x^2} + 5\sqrt{x} = f(x)$.
\emph{Conclusion :} L'ensemble des primitives de $f$ sur $]0, +\infty[$ est $\{ F(x) = x^3 + \frac{4}{x} + \frac{10}{3}x\sqrt{x} + C \mid C \in \mathbb{R} \}$.
\end{exoresolu}

\begin{methode}[Utiliser la linéarité de la primitive (somme et multiple)]
\emph{Objectif :} Décomposer une fonction complexe en somme/différence de fonctions usuelles et factoriser les constantes pour primitiver terme à terme.
\begin{enumerate}
    \item Écrire $f(x)$ sous forme de somme algébrique : $f(x) = \sum k_i g_i(x)$ où $g_i$ sont des fonctions usuelles.
    \item Appliquer : $\int (u+v) = \int u + \int v$ et $\int k u = k \int u$ ($k$ constant).
    \item Primitiver chaque $g_i$ séparément.
    \item Regrouper les résultats et n'ajouter \textbf{qu'une seule constante $+C$} à la fin.
\end{enumerate}
\cle{Attention :} Ne pas mettre de $+C$ après chaque terme intermédiaire. Une seule constante pour toute la primitive.
\end{methode}

\begin{exoresolu}[Combinaison linéaire de formes usuelles]
\emph{Énoncé :} Trouver une primitive $F$ de $f(x) = \frac{2x^3 - 3x + 1}{x^2}$ sur $]0, +\infty[$.
\tcblower
\emph{Solution détaillée :}
On simplifie d'abord l'expression pour faire apparaître des puissances de $x$ (séparation du numérateur) :
\[ f(x) = \frac{2x^3}{x^2} - \frac{3x}{x^2} + \frac{1}{x^2} = 2x - \frac{3}{x} + x^{-2} \]
On primitivie terme à terme sur $]0, +\infty[$ (donc $|x|=x$) :
\begin{align*}
\int 2x \, dx &= 2 \cdot \frac{x^2}{2} = x^2 \\
\int -\frac{3}{x} \, dx &= -3 \ln x \\
\int x^{-2} \, dx &= \frac{x^{-1}}{-1} = -\frac{1}{x}
\end{align*}
En sommant et ajoutant $C$ :
\[ F(x) = x^2 - 3\ln x - \frac{1}{x} + C \quad (C \in \mathbb{R}) \]
\emph{Vérification :} $F'(x) = 2x - \frac{3}{x} + \frac{1}{x^2} = f(x)$.
\emph{Conclusion :} Les primitives de $f$ sur $]0,+\infty[$ sont les fonctions $F(x) = x^2 - 3\ln x - \frac{1}{x} + C$, $C \in \mathbb{R}$.
\end{exoresolu}

\begin{methode}[Primitiver une fonction composée de la forme $u'(x) \cdot g(u(x))$ (Règle de la chaîne inverse)]
\emph{Objectif :} Reconnaître la dérivée d'une fonction composée $G(u(x))$ où $G'=g$. On utilise la formule : $\int u'(x) g(u(x)) \, dx = G(u(x)) + C$.
\begin{enumerate}
    \item Identifier une fonction intérieure $u(x)$ dont la dérivée $u'(x)$ apparaît (ou un multiple constant) dans $f(x)$.
    \item Écrire $f(x) = k \cdot u'(x) \cdot g(u(x))$ ($k$ constant).
    \item Trouver une primitive $G$ de $g$ (sur l'ensemble des valeurs de $u$).
    \item La primitive cherchée est $F(x) = k \cdot G(u(x)) + C$.
    \item \textbf{Vérifier systématiquement en dérivant} (règle de la chaîne).
\end{enumerate}
\cle{Cas fréquents au Probatoire :}
\begin{itemize}
    \item $\int u'(x) u(x)^\alpha dx = \frac{u(x)^{\alpha+1}}{\alpha+1}+C$ ($\alpha \neq -1$)
    \item $\int \frac{u'(x)}{u(x)} dx = \ln|u(x)|+C$
    \item $\int u'(x) e^{u(x)} dx = e^{u(x)}+C$
    \item $\int u'(x) \cos(u(x)) dx = \sin(u(x))+C$
    \item $\int u'(x) \sin(u(x)) dx = -\cos(u(x))+C$
\end{itemize}
\end{methode}

\begin{exoresolu}[Fonction composée : forme $u'/u$ et forme puissance]
\emph{Énoncé :} Déterminer les primitives sur $]1, +\infty[$ de $f(x) = \frac{2x}{x^2-1} + 3x(x^2+1)^2$.
\tcblower
\emph{Solution détaillée :}
On traite les deux termes séparément.
\textbf{1er terme :} $f_1(x) = \frac{2x}{x^2-1}$.
On pose $u(x) = x^2-1$. Alors $u'(x) = 2x$. On a bien la forme $\frac{u'(x)}{u(x)}$.
Une primitive de $\frac{1}{u}$ est $\ln|u|$. Comme sur $]1,+\infty[$, $u(x)=x^2-1 > 0$, $|u|=u$.
Primitive : $F_1(x) = \ln(x^2-1)$.

\textbf{2e terme :} $f_2(x) = 3x(x^2+1)^2$.
On pose $v(x) = x^2+1$. Alors $v'(x) = 2x$. On a $3x = \frac{3}{2} \cdot 2x = \frac{3}{2}v'(x)$.
$f_2(x) = \frac{3}{2} v'(x) \cdot v(x)^2$. Forme $u' \cdot u^\alpha$ avec $\alpha=2$.
Primitive de $v^2$ est $\frac{v^3}{3}$.
Primitive : $F_2(x) = \frac{3}{2} \cdot \frac{v(x)^3}{3} = \frac{1}{2}(x^2+1)^3$.

\textbf{Somme :} $F(x) = \ln(x^2-1) + \frac{1}{2}(x^2+1)^3 + C$, $C \in \mathbb{R}$.
\emph{Vérification :}
$F'(x) = \frac{2x}{x^2-1} + \frac{1}{2} \cdot 3(x^2+1)^2 \cdot 2x = \frac{2x}{x^2-1} + 3x(x^2+1)^2 = f(x)$.
\emph{Conclusion :} L'ensemble des primitives est $\{ \ln(x^2-1) + \frac{1}{2}(x^2+1)^3 + C \mid C \in \mathbb{R} \}$.
\end{exoresolu}

\begin{methode}[Déterminer la constante $C$ par une condition initiale (Problème de Cauchy)]
\emph{Objectif :} Trouver \textbf{la} primitive particulière qui satisfait une condition $F(x_0) = y_0$ (passage par un point, valeur initiale).
\begin{enumerate}
    \item Calculer l'expression générale de la primitive $F(x) = F_0(x) + C$ (où $F_0$ est une primitive quelconque, par exemple la somme des primitives usuelles sans constante).
    \item Remplacer $x$ par $x_0$ et $F(x)$ par $y_0$ dans l'équation : $y_0 = F_0(x_0) + C$.
    \item Résoudre pour $C$ : $C = y_0 - F_0(x_0)$.
    \item Écrire l'expression finale de $F(x)$ en remplaçant $C$ par sa valeur numérique.
    \item Vérifier que $F(x_0)=y_0$ et $F'=f$.
\end{enumerate}
\cle{Rédaction type :} « Soit $F$ une primitive de $f$ sur $I$. $F(x) = \dots + C$. La condition $F(x_0)=y_0$ donne $C = \dots$. Donc $F(x) = \dots$ »
\end{methode}

\begin{exoresolu}[Problème de Cauchy : condition initiale]
\emph{Énoncé :} La fonction $f$ est définie sur $]0, +\infty[$ par $f(x) = 3x^2 - \frac{2}{x}$. Déterminer la primitive $F$ de $f$ qui prend la valeur $4$ en $1$ (c'est-à-dire $F(1)=4$).
\tcblower
\emph{Solution détaillée :}
\textbf{Étape 1 : Primitive générale.}
$f(x) = 3x^2 - 2x^{-1}$.
Une primitive sur $]0,+\infty[$ est :
$F_0(x) = 3 \cdot \frac{x^3}{3} - 2 \ln x = x^3 - 2\ln x$.
L'ensemble des primitives est $F(x) = x^3 - 2\ln x + C$, $C \in \mathbb{R}$.

\textbf{Étape 2 : Utilisation de la condition $F(1)=4$.}
$F(1) = 1^3 - 2\ln 1 + C = 1 - 0 + C = 1 + C$.
On impose $1 + C = 4 \implies C = 3$.

\textbf{Étape 3 : Expression finale.}
$F(x) = x^3 - 2\ln x + 3$.

\textbf{Vérification :} $F'(x) = 3x^2 - \frac{2}{x} = f(x)$ et $F(1) = 1 - 0 + 3 = 4$. \checkmark
\emph{Conclusion :} La primitive cherchée est $F(x) = x^3 - 2\ln x + 3$.
\end{exoresolu}

\begin{methode}[Réécrire l'expression avant de primitiver (division, factorisation, changement de variable simple)]
\emph{Objectif :} Transformer une expression algébrique (fraction rationnelle, produit, racine imbriquée) en une somme de formes usuelles ou en forme composée $u' g(u)$.
\begin{enumerate}
    \item \textbf{Fraction rationnelle :} Effectuer la division euclidienne si $\deg(\text{num}) \ge \deg(\text{den})$, ou séparer les termes : $\frac{A+B}{C} = \frac{A}{C} + \frac{B}{C}$.
    \item \textbf{Racines / Puissances fractionnaires :} Réécrire $\sqrt[n]{x^m} = x^{m/n}$.
    \item \textbf{Produit :} Développer si possible (ex: $(x+1)^2 = x^2+2x+1$).
    \item \textbf{Changement de variable explicite :} Poser $u = \varphi(x)$, exprimer $dx = \frac{du}{\varphi'(x)}$, primitiver en $u$, revenir à $x$.
\end{enumerate}
\cle{Astuce :} Pour $\int \frac{P(x)}{Q(x)}dx$ avec $\deg P < \deg Q$, essayer d'écrire le numérateur comme dérivée du dénominateur (ou combinaison) pour utiliser $\int \frac{u'}{u} = \ln|u|$.
\end{methode}

\begin{exoresolu}[Transformation algébrique : Division et décomposition]
\emph{Énoncé :} Calculer une primitive sur $]0, +\infty[$ de $f(x) = \frac{x^3 + 2x^2 - x + 1}{x^2}$.
\tcblower
\emph{Solution détaillée :}
Le degré du numérateur (3) est strictement supérieur à celui du dénominateur (2). On effectue la division euclidienne (ou on sépare les termes) :
\[ \frac{x^3 + 2x^2 - x + 1}{x^2} = \frac{x^3}{x^2} + \frac{2x^2}{x^2} - \frac{x}{x^2} + \frac{1}{x^2} = x + 2 - \frac{1}{x} + x^{-2} \]
On primitivie terme à terme sur $]0,+\infty[$ :
\begin{align*}
\int x \, dx &= \frac{x^2}{2} \\
\int 2 \, dx &= 2x \\
\int -\frac{1}{x} \, dx &= -\ln x \\
\int x^{-2} \, dx &= -x^{-1} = -\frac{1}{x}
\end{align*}
En sommant et ajoutant $C$ :
\[ F(x) = \frac{x^2}{2} + 2x - \ln x - \frac{1}{x} + C \quad (C \in \mathbb{R}) \]
\emph{Vérification :} $F'(x) = x + 2 - \frac{1}{x} + \frac{1}{x^2} = \frac{x^3 + 2x^2 - x + 1}{x^2} = f(x)$.
\emph{Conclusion :} Les primitives de $f$ sur $]0,+\infty[$ sont $F(x) = \frac{x^2}{2} + 2x - \ln x - \frac{1}{x} + C$.
\end{exoresolu}

\coursec{Applications concrètes : modélisation de situations réelles}

\begin{methode}[Reconstituer une grandeur à partir de sa variation (Vitesse $\to$ Position, Coût marginal $\to$ Coût total)]
\emph{Objectif :} Modéliser une situation réelle où l'on connaît le taux de variation instantané (dérivée) et l'on cherche la grandeur totale (primitive), en utilisant une condition initiale (valeur au départ).
\begin{enumerate}
    \item Identifier la fonction dérivée $f(t)$ (vitesse, coût marginal, débit, etc.) et la grandeur primitive $F(t)$ (position, coût total, volume).
    \item Noter la condition initiale : $F(t_0) = F_0$ (position initiale, coût fixe, volume initial).
    \item Calculer une primitive $F_0(t)$ de $f(t)$.
    \item Déterminer $C$ grâce à $F(t_0)=F_0$.
    \item Répondre à la question posée (position à l'instant $t$, coût pour $q$ unités, etc.).
\end{enumerate}
\cle{Unités :} Vérifier la cohérence des unités (ex: $v$ en m/s, $t$ en s $\implies$ $x$ en m). $C$ a l'unité de $F$.
\end{methode}

\begin{popfigure}
\begin{tikzpicture}[scale=0.7, >=Stealth]
\begin{axis}[
    axis lines=middle, xlabel=$t$ (h), ylabel=$v$ (km/h),
    xmin=0, xmax=4.5, ymin=0, ymax=100,
    xtick={0,1,2,3,4}, ytick={0,20,40,60,80,100},
    tick label style={font=\small},
    fill=popcream
]
\addplot[popcurve, domain=0:4, samples=100, thick] {60 + 20*sqrt(max(0,x))};
\addplot[poporange, thick, domain=0:3.52, samples=60] {60 + 20*sqrt(max(0,x))};
\draw[dashed, popdark] (3.52,0) node[below] {$t_1$} -- (3.52,97.5);
\draw[dashed, popdark] (0,97.5) node[left] {$v(t_1)$} -- (3.52,97.5);
\node[poporange, align=center] at (2, 40) {Aire = Distance\\ $\int_0^{t_1} v(t)dt = 300$ km};
\end{axis}
\end{tikzpicture}
\legende{Modélisation vitesse/position : l'aire sous la courbe de vitesse donne le déplacement.}
\end{popfigure}

\begin{exoresolu}[Modélisation : Distance parcourue à partir de la vitesse]
\emph{Énoncé :} Un véhicule démarre d'un péage (abscisse $0$ km) à l'instant $t=0$ h. Sa vitesse (en km/h) est modélisée par $v(t) = 60 + 20\sqrt{t}$ pour $t \in [0, 4]$. Déterminer l'heure (en heures et minutes) à laquelle il se trouve à $300$ km du péage.
\tcblower
\emph{Solution détaillée :}
\textbf{Modélisation :} La position $x(t)$ (en km) est une primitive de la vitesse $v(t)$. $x'(t) = v(t)$. Condition initiale : $x(0) = 0$.

\textbf{Primitive de $v$ :}
$v(t) = 60 + 20 t^{1/2}$.
$X_0(t) = \int 60 \, dt + \int 20 t^{1/2} \, dt = 60t + 20 \cdot \frac{t^{3/2}}{3/2} = 60t + \frac{40}{3} t\sqrt{t}$.
Primitive générale : $x(t) = 60t + \frac{40}{3} t\sqrt{t} + C$.

\textbf{Condition initiale :} $x(0) = 0 + 0 + C = 0 \implies C = 0$.
Donc $x(t) = 60t + \frac{40}{3} t\sqrt{t}$.

\textbf{Résolution de l'équation $x(t) = 300$ :}
$60t + \frac{40}{3} t\sqrt{t} = 300$.
On divise par $20$ : $3t + \frac{2}{3} t\sqrt{t} = 15$.
Multiplier par $3$ : $9t + 2t\sqrt{t} = 45$.
Posons $u = \sqrt{t} \ge 0$, alors $t = u^2$.
$9u^2 + 2u^3 = 45 \iff 2u^3 + 9u^2 - 45 = 0$.
On étudie $g(u) = 2u^3 + 9u^2 - 45$ sur $[0, 2]$ (car $t \in [0,4] \implies u \in [0,2]$).
$g(1.5) = 2(3,375) + 9(2,25) - 45 = 6,75 + 20,25 - 45 = -18 < 0$.
$g(2) = 2(8) + 9(4) - 45 = 16 + 36 - 45 = 7 > 0$.
La solution $u_0$ est dans $]1,5 ; 2[$. Par approximation (calculatrice ou dichotomie) : $u_0 \approx 1,875$.
Alors $t_0 = u_0^2 \approx 3,515$ h.
Conversion en minutes : $0,515 \times 60 \approx 31$ min.
\emph{Conclusion :} Le véhicule se trouve à 300 km du péage vers \textbf{3 h 31 min} (arrondi à la minute).
\end{exoresolu}

\coursec{Synthèse, entraînement et évaluation formative}

\begin{aretenir}[Résumé des étapes pour primitiver]
\begin{enumerate}
    \item \textbf{Domaine :} Préciser l'intervalle $I$ de continuité de $f$.
    \item \textbf{Simplification :} Réécrire $f$ (division, développement, puissances, valeurs absolues).
    \item \textbf{Linéarité :} Séparer les termes, factoriser les constantes.
    \item \textbf{Formes usuelles / Composée :} Appliquer la table ou la règle de la chaîne inverse ($u' g(u)$).
    \item \textbf{Constante :} Ajouter $+C$ ($C \in \mathbb{R}$).
    \item \textbf{Vérification :} Dériver le résultat trouvé.
    \item \textbf{Condition initiale (si problème de Cauchy) :} Calculer $C$ numérique.
\end{enumerate}
\end{aretenir}

\begin{attention}[Les 5 erreurs qui coûtent des points au Probatoire]
\begin{enumerate}
    \item \textbf{Oublier la constante $+C$ (ou écrire $+c$ minuscule) :} Une primitive n'est pas unique ; l'ensemble des primitives est une famille paramétrée par $C \in \mathbb{R}$. Sans $+C$, la réponse est incomplète (perte de 1 à 2 points).
    \item \textbf{Écrire $\ln x$ au lieu de $\ln|x|$ pour $\int \frac{1}{x} dx$ :} Si l'intervalle n'est pas explicitement $]0,+\infty[$, les valeurs absolues sont obligatoires (domaine $\mathbb{R}^*$). Oubli sanctionné.
    \item \textbf{Mauvaise gestion des puissances :} $\int x^\alpha dx = \frac{x^{\alpha+1}}{\alpha+1}$ \textbf{uniquement si $\alpha \neq -1$}. Erreur classique : $\int x^{-1} dx = \frac{x^0}{0}$ (absurde) au lieu de $\ln|x|$. Aussi : $\int \sqrt{x} dx = \int x^{1/2} dx = \frac{2}{3}x^{3/2}$ (ne pas écrire $\frac{3}{2}x^{1/2}$).
    \item \textbf{Oublier la dérivée intérieure dans la règle de la chaîne inverse :} Pour $\int (2x+1)^3 dx$, on ne peut pas écrire $\frac{(2x+1)^4}{4}$. Il faut le facteur $2$ (dérivée de $2x+1$). La primitive est $\frac{1}{2} \cdot \frac{(2x+1)^4}{4} = \frac{(2x+1)^4}{8}$. Vérifier en dérivant !
    \item \textbf{Confondre Primitive et Intégrale (notation) :} Au Probatoire, on demande « Déterminer les primitives de $f$ ». La réponse attendue est une \textbf{fonction} $F(x) = \dots + C$. Ne pas écrire $\int f(x) dx = \dots$ sans expliquer que c'est une primitive, et surtout ne pas calculer une intégrale définie (avec bornes) si ce n'est pas demandé.
\end{enumerate}
\end{attention}

\begin{exercice}[Entraînement 1 : Primitives usuelles et linéarité]
Déterminer l'ensemble des primitives sur l'intervalle de définition de chaque fonction :
\begin{enumerate}
    \item $f(x) = 4x^3 - 5x + \frac{2}{x}$ sur $]0, +\infty[$.
    \item $g(x) = \frac{3}{\sqrt{x}} - 2x^2 + 7$ sur $]0, +\infty[$.
    \item $h(x) = (x+2)(x-3)$ sur $\mathbb{R}$.
    \item $k(x) = \frac{x^4 - 2x^2 + 1}{x^2}$ sur $]-\infty, 0[ \cup ]0, +\infty[$.
\end{enumerate}
\end{exercice}

\begin{exercice}[Entraînement 2 : Formes composées $u' g(u)$]
Déterminer les primitives des fonctions suivantes sur les intervalles indiqués :
\begin{enumerate}
    \item $f(x) = 2x \cos(x^2)$ sur $\mathbb{R}$.
    \item $g(x) = \frac{3x^2}{x^3+1}$ sur $]-1, +\infty[$.
    \item $h(x) = (2x-1)e^{x^2-x}$ sur $\mathbb{R}$.
    \item $k(x) = 4x^3 \sqrt{x^4+5}$ sur $\mathbb{R}$.
\end{enumerate}
\end{exercice}

\begin{exercice}[Entraînement 3 : Problèmes de Cauchy et Applications]
\begin{enumerate}
    \item Trouver la primitive $F$ de $f(x) = 6x^2 - \frac{4}{x^3}$ sur $]0, +\infty[$ telle que $F(1) = 2$.
    \item \textbf{Électricité :} L'intensité d'un courant dans un circuit est donnée par $i(t) = 2 + 3t^2$ (en Ampères) pour $t \ge 0$. La charge $q(t)$ (en Coulombs) satisfait $q'(t) = i(t)$. Sachant que $q(0) = 5$ C, exprimer $q(t)$ et calculer la charge passée à $t=2$ s.
    \item \textbf{Gestion :} Le coût marginal (en milliers de FCFA/unité) de production de $x$ unités est $C_m(x) = 10 - \frac{100}{(x+10)^2}$. Les coûts fixes sont de 500 mille FCFA. Déterminer la fonction coût total $C(x)$ et le coût de production de 90 unités.
\end{enumerate}
\end{exercice}

\begin{corrige}[Exercice 1]
\begin{enumerate}
    \item $F(x) = x^4 - \frac{5}{2}x^2 + 2\ln x + C,\ C\in\mathbb{R}$.
    \item $G(x) = 6\sqrt{x} - \frac{2}{3}x^3 + 7x + C,\ C\in\mathbb{R}$.
    \item $H(x) = \frac{x^3}{3} - \frac{x^2}{2} - 6x + C,\ C\in\mathbb{R}$.
    \item $K(x) = \frac{x^3}{3} - 2x - \frac{1}{x} + C,\ C\in\mathbb{R}$ (sur chaque intervalle $]-\infty,0[$ et $]0,+\infty[$).
\end{enumerate}
\end{corrige}

\begin{corrige}[Exercice 2]
\begin{enumerate}
    \item $u=x^2, u'=2x \implies F(x) = \sin(x^2) + C$.
    \item $u=x^3+1, u'=3x^2 \implies G(x) = \ln(x^3+1) + C$ ($x^3+1>0$ sur $]-1,+\infty[$).
    \item $u=x^2-x, u'=2x-1 \implies H(x) = e^{x^2-x} + C$.
    \item $u=x^4+5, u'=4x^3 \implies K(x) = \frac{2}{3}(x^4+5)^{3/2} + C$.
\end{enumerate}
\end{corrige}

\begin{corrige}[Exercice 3]
\begin{enumerate}
    \item $F_0(x) = 2x^3 + \frac{2}{x^2}$. $F(1)=2+2+C=4+C=2 \implies C=-2$. $F(x) = 2x^3 + \frac{2}{x^2} - 2$.
    \item $q(t) = \int (2+3t^2)dt = 2t + t^3 + C$. $q(0)=C=5$. $q(t)=t^3+2t+5$. $q(2)=8+4+5=17$ C.
    \item $C(x) = \int (10 - 100(x+10)^{-2})dx = 10x + \frac{100}{x+10} + C$. $C(0)=0+10+C=500 \implies C=490$. $C(x)=10x + \frac{100}{x+10} + 490$. $C(90)=900 + 1 + 490 = 1391$ mille FCFA.
\end{enumerate}
\end{corrige}

\newpage

\coursec{Exercices}

\coursec{Primitives d'une fonction continue sur un intervalle}

\courssub{Rappels : dérivée et notion de primitive}

\begin{definition}[Dérivée d'une fonction en un point]
Soit $f$ une fonction définie sur un intervalle $I$ et $a \in I$. Si le taux d'accroissement $\dfrac{f(x)-f(a)}{x-a}$ admet une limite finie lorsque $x \to a$, on dit que $f$ est \cle{dérivable} en $a$. Cette limite est appelée \cle{dérivée} de $f$ en $a$ et se note $f'(a)$.
\[
f'(a) = \lim_{x \to a} \frac{f(x)-f(a)}{x-a}
\]
La fonction $f'$ qui associe $f'(x)$ à tout $x$ où $f$ est dérivable est la \cle{fonction dérivée} de $f$.
\end{definition}

\begin{definition}[Primitive d'une fonction sur un intervalle]
Soit $f$ une fonction définie sur un intervalle $I$. On appelle \cle{primitive} de $f$ sur $I$ toute fonction $F$ dérivable sur $I$ telle que :
\[
\forall x \in I,\quad F'(x) = f(x)
\]
On note souvent $F' = f$ sur $I$.
\end{definition}

\begin{exemplebox}[Exemples simples]
\begin{itemize}
    \item $f(x)=2x$ sur $\mathbb{R}$. Une primitive est $F(x)=x^2$ car $F'(x)=2x=f(x)$.
    \item $f(x)=\cos x$ sur $\mathbb{R}$. Une primitive est $F(x)=\sin x$ car $F'(x)=\cos x=f(x)$.
    \item $f(x)=e^x$ sur $\mathbb{R}$. Une primitive est $F(x)=e^x$ car $F'(x)=e^x=f(x)$.
\end{itemize}
\end{exemplebox}

\begin{attention}[Notation et vocabulaire]
\begin{itemize}
    \item On ne dit pas « \emph{la} primitive » mais « \emph{une} primitive » car il en existe une infinité.
    \item L'opération qui consiste à trouver une primitive s'appelle l'\cle{intégration} (ou recherche de primitive). C'est l'opération inverse de la dérivation.
    \item La condition « $f$ définie sur un \textbf{intervalle} » est essentielle (voir théorème ci-dessous).
\end{itemize}
\end{attention}

\courssub{Théorème d'existence des primitives pour les fonctions continues}

\begin{propriete}[Théorème fondamental : existence des primitives]
\emph{(Exigible au Probatoire : énoncé et utilisation)}
\begin{center}
\textbf{Toute fonction continue sur un intervalle $I$ admet au moins une primitive sur $I$.}
\end{center}
Plus précisément, pour tout $a \in I$, la fonction $F_a$ définie par $F_a(x) = \int_a^x f(t)\,dt$ est une primitive de $f$ sur $I$ qui s'annule en $a$.
\end{propriete}

\begin{aretenir}[Conséquences immédiates]
\begin{enumerate}
    \item Les fonctions polynomiales, rationnelles (sur leur ensemble de définition qui est une réunion d'intervalles), trigonométriques usuelles, exponentielles, logarithmes (sur $]0;+\infty[$) sont continues sur leurs intervalles de définition, donc admettent des primitives.
    \item Une fonction qui n'est pas continue sur un intervalle (ex: fonction partie entière, fonction avec asymptote verticale à l'intérieur de l'intervalle) n'admet pas nécessairement de primitive sur cet intervalle.
\end{enumerate}
\end{aretenir}

\begin{experience}[Découverte : pourquoi l'intervalle est nécessaire ?]
Considérons $f(x) = \frac{1}{x}$ sur $D_f = ]-\infty;0[ \cup ]0;+\infty[$.
\begin{enumerate}
    \item Montrer que $F(x) = \ln|x|$ est une primitive de $f$ sur $]-\infty;0[$ et sur $]0;+\infty[$.
    \item Peut-on définir une primitive de $f$ sur $D_f$ tout entier (qui n'est pas un intervalle) ?
    \item Essayer de définir $G(x) = \ln|x| + C_1$ pour $x<0$ et $G(x) = \ln|x| + C_2$ pour $x>0$. $G$ est-elle dérivable sur $D_f$ ? $G'=f$ ?
\end{enumerate}
\emph{Conclusion :} Sur un ensemble qui n'est pas un intervalle, une fonction continue admet des primitives \emph{par morceaux}, mais pas nécessairement une primitive globale unique (les constantes peuvent différer sur chaque composante connexe).
\end{experience}

\courssub{Propriétés algébriques des primitives}

\begin{propriete}[Linéarité et famille des primitives]
Soit $f$ une fonction continue sur un intervalle $I$.
\begin{enumerate}
    \item \textbf{Famille des primitives :} Si $F$ est une primitive de $f$ sur $I$, alors pour toute constante $k \in \mathbb{R}$, la fonction $x \mapsto F(x) + k$ est aussi une primitive de $f$ sur $I$. Réciproquement, si $F_1$ et $F_2$ sont deux primitives de $f$ sur $I$, il existe une constante $k$ telle que $F_1(x) = F_2(x) + k$ pour tout $x \in I$.
    \item \textbf{Somme :} Si $F$ est une primitive de $f$ et $G$ une primitive de $g$ sur $I$, alors $F+G$ est une primitive de $f+g$ sur $I$.
    \item \textbf{Multiple par un scalaire :} Si $F$ est une primitive de $f$ sur $I$ et $\lambda \in \mathbb{R}$, alors $\lambda F$ est une primitive de $\lambda f$ sur $I$.
\end{enumerate}
\end{propriete}

\begin{attention}[Pièges classiques]
\begin{itemize}
    \item \textbf{Mauvaise linéarité :} Une primitive de $f \times g$ n'est \textbf{pas} le produit d'une primitive de $f$ par une primitive de $g$. De même pour le quotient $f/g$ (sauf cas particuliers $u'/u$, $u'/u^2$).
    \item \textbf{Domaine de définition :} Vérifier toujours que les fonctions sont continues sur un \textbf{même intervalle} $I$ avant d'appliquer ces propriétés.
    \item \textbf{Constante d'intégration :} Ne pas oublier le $+k$ (ou $+C$) dans la réponse finale d'une primitive indéfinie.
\end{itemize}
\end{attention}

\courssub{Détermination de primitives : table usuelles et méthodes}

\begin{aretenir}[Table des primitives usuelles (à connaître par cœur)]
\begin{center}
\begin{tabular}{|c|c|c|}\hline
Fonction $f(x)$ & Intervalle $I$ & Primitive $F(x)$ \\ \hline
$x^n$ ($n \in \mathbb{N}$) & $\mathbb{R}$ & $\dfrac{x^{n+1}}{n+1}$ \\ \hline
$x^\alpha$ ($\alpha \in \mathbb{Q}\setminus\{-1\}$) & $]0;+\infty[$ ($[0;+\infty[$ si $\alpha \ge 0$) & $\dfrac{x^{\alpha+1}}{\alpha+1}$ \\ \hline
$\dfrac{1}{x}$ & $]-\infty;0[$ et $]0;+\infty[$ & $\ln|x|$ \\ \hline
$e^x$ & $\mathbb{R}$ & $e^x$ \\ \hline
$\cos x$ & $\mathbb{R}$ & $\sin x$ \\ \hline
$\sin x$ & $\mathbb{R}$ & $-\cos x$ \\ \hline
$\dfrac{1}{1+x^2}$ & $\mathbb{R}$ & $\arctan x$ \\ \hline
$\dfrac{1}{\sqrt{1-x^2}}$ & $]-1;1[$ & $\arcsin x$ \\ \hline
\end{tabular}
\end{center}
\end{aretenir}

\begin{methode}[Recherche d'une primitive d'une fonction continue sur $I$]
\begin{enumerate}
    \item \textbf{Simplifier / Développer} l'expression de $f(x)$ (mettre sous forme somme de termes usuels).
    \item \textbf{Identifier} chaque terme comme une forme usuelle ou une forme composée $u'(x) \times g(u(x))$.
    \item \textbf{Appliquer} la linéarité : primitive de la somme = somme des primitives.
    \item \textbf{Utiliser} les formules de primitives composées (chaîne) :
    \begin{itemize}
        \item $\int u'(x) u(x)^n \,dx = \frac{u(x)^{n+1}}{n+1}$ ($n \neq -1$)
        \item $\int \frac{u'(x)}{u(x)} \,dx = \ln|u(x)|$
        \item $\int u'(x) e^{u(x)} \,dx = e^{u(x)}$
        \item $\int u'(x) \cos(u(x)) \,dx = \sin(u(x))$
        \item $\int u'(x) \sin(u(x)) \,dx = -\cos(u(x))$
    \end{itemize}
    \item \textbf{Ajouter} la constante $+k$ (ou $+C$).
    \item \textbf{Vérifier} en dérivant le résultat trouvé.
\end{enumerate}
\end{methode}

\begin{exemplebox}[Primitive composée : cas $u'/u$]
Soit $f(x) = \frac{3x^2+2x}{x^3+x^2+5}$ sur $\mathbb{R}$.
Le dénominateur $u(x)=x^3+x^2+5$ a pour dérivée $u'(x)=3x^2+2x$ qui est le numérateur.
De plus $u(x) > 0$ sur $\mathbb{R}$ (discriminant $<0$ et coefficient directeur $>0$).
Une primitive est $F(x) = \ln(x^3+x^2+5)$ (pas de valeur absolue nécessaire car $u>0$).
\end{exemplebox}

\begin{exemplebox}[Primitive avec condition initiale]
Trouver la primitive $F$ de $f(x)=3x^2-4x+2$ sur $\mathbb{R}$ telle que $F(1)=5$.
\begin{enumerate}
    \item Primitive générale : $F(x) = x^3 - 2x^2 + 2x + k$.
    \item Condition : $F(1) = 1 - 2 + 2 + k = 1 + k = 5 \Rightarrow k = 4$.
    \item Réponse : $F(x) = x^3 - 2x^2 + 2x + 4$.
\end{enumerate}
\end{exemplebox}

\courssub{Applications concrètes : modélisation de situations réelles}

\begin{savaistu}[Du coût marginal au coût total -- Économie]
En économie, le \cle{coût marginal} $C'_m(x)$ est la dérivée de la fonction \cle{coût total} $C_t(x)$. $C_t(x) = C_f + C_v(x)$ où $C_f$ sont les coûts fixes (indépendants de $x$) et $C_v(x)$ les coûts variables.
Connaître $C'_m$ permet de retrouver $C_t$ par intégration : $C_t(x) = \int C'_m(x)\,dx$. La constante d'intégration est exactement le coût fixe $C_f = C_t(0)$.
\end{savaistu}

\begin{savaistu}[Vitesse et position -- Mécanique / Physique]
Si $v(t)$ est la vitesse (dérivée de la position), alors la position $x(t)$ est une primitive de $v(t)$.
La distance parcourue entre $t_1$ et $t_2$ est $\int_{t_1}^{t_2} |v(t)|\,dt$ (valeur absolue si changement de sens).
Si $v(t) \ge 0$ sur l'intervalle, distance $= x(t_2) - x(t_1)$.
\end{savaistu}

\begin{savaistu}[Courant et charge -- Électricité]
L'intensité du courant $i(t)$ est la dérivée de la charge électrique $q(t)$ : $i(t) = q'(t)$.
La charge qui traverse une section entre $t_1$ et $t_2$ est $q(t_2)-q(t_1) = \int_{t_1}^{t_2} i(t)\,dt$.
\end{savaistu}

\courssub{Je vérifie mes connaissances}

\begin{exercice}[QCM : Primitive d'une fonction polynôme]
Parmi les fonctions suivantes, laquelle est une primitive de $f(x)=3x^2$ sur $\mathbb{R}$ ?
\begin{enumerate}
    \item $F(x)=x^3+2x$
    \item $F(x)=x^3$
    \item $F(x)=6x$
    \item $F(x)=x^3+x$
\end{enumerate}
\emph{Justifiez votre choix en dérivant chaque proposition.}
\end{exercice}

\begin{exercice}[Vrai ou Faux : Propriétés algébriques]
Soit $f$ une fonction continue sur un intervalle $I$ et $F$ une primitive de $f$ sur $I$.
\begin{enumerate}
    \item La fonction $G(x)=F(x)+5$ est une primitive de $f$ sur $I$.
    \item La fonction $H(x)=2F(x)$ est une primitive de $f$ sur $I$.
    \item Si $F_1$ et $F_2$ sont deux primitives de $f$ sur $I$, alors $F_1-F_2$ est constante sur $I$.
\end{enumerate}
Répondez par Vrai ou Faux et justifiez chaque réponse en utilisant la définition de la primitive.
\end{exercice}

\begin{exercice}[Définition et vocabulaire]
Complétez la phrase suivante en utilisant les termes appropriés : \cle{primitive}, \cle{dérivée}, \cle{continue}, \cle{intervalle}, \cle{constante}.
\bigskip
\noindent « Soit $f$ une fonction \underline{\hspace{3cm}} sur un \underline{\hspace{3cm}} $I$. On appelle \underline{\hspace{3cm}} de $f$ sur $I$ toute fonction $F$ dérivable sur $I$ telle que $F'=f$ sur $I$. Si $F$ est une primitive de $f$, alors pour tout réel $k$, la fonction $x\mapsto F(x)+k$ est aussi une primitive de $f$. Réciproquement, deux primitives de $f$ sur $I$ ne diffèrent que d'une \underline{\hspace{3cm}}. »
\end{exercice}

\begin{exercice}[Calcul direct de primitives]
Déterminer une primitive sur l'intervalle indiqué des fonctions suivantes :
\begin{enumerate}
    \item $f(x)=4x^3-2x+5$ sur $\mathbb{R}$
    \item $g(x)=\dfrac{1}{x^2}$ sur $]-\infty;0[$ et $]0;+\infty[$
    \item $h(x)=\sqrt{x}$ sur $[0;+\infty[$
    \item $k(x)=\cos(2x)$ sur $\mathbb{R}$
\end{enumerate}
\end{exercice}

\courssub{Je m'entraîne}

\begin{exercice}[Coût total de production]
Dans une usine de transformation agroalimentaire, le coût marginal (en FCFA par unité produite) est modélisé par la fonction $C'_m(x)=0,01x^2-0,5x+10$, où $x$ est le nombre d'unités produites. Le coût fixe (loyer, assurances) est de $5\,000$ FCFA.
\begin{enumerate}
    \item Exprimer le coût total $C_t(x)$ en fonction de $x$, en choisissant la primitive qui convient.
    \item Calculer le coût total de production de $100$ unités.
    \item Interpréter le terme $0,01x^2$ dans l'expression du coût marginal.
\end{enumerate}
\end{exercice}

\begin{exercice}[Charge électrique traversant un conducteur]
L'intensité du courant électrique (en ampères) traversant un dipôle est donnée par $i(t)=2e^{-0,5t}$ pour $t\ge 0$ ($t$ en secondes). La charge électrique $q(t)$ (en coulombs) est une primitive de $i(t)$ telle que $q(0)=0$.
\begin{enumerate}
    \item Déterminer l'expression de $q(t)$.
    \item Calculer la charge qui a traversé la section du conducteur entre $t=0$ et $t=2\,\text{s}$.
    \item Quelle est la limite de $q(t)$ quand $t\to +\infty$ ? Interpréter physiquement.
\end{enumerate}
\end{exercice}

\begin{exercice}[Distance parcourue par un mobile]
Un véhicule démarre au temps $t=0$ avec une vitesse (en m/s) donnée par $v(t)=10-t$ pour $t\in[0;10]$.
\begin{enumerate}
    \item Montrer que $v$ est continue sur $[0;10]$ et en déduire l'existence d'une primitive.
    \item Déterminer l'équation horaire $x(t)$ (position en mètres) sachant que $x(0)=0$.
    \item Calculer la distance totale parcourue avant l'arrêt du véhicule.
\end{enumerate}
\end{exercice}

\begin{exercice}[Primitive d'un quotient de type $u'/u$]
Soit la fonction $f$ définie sur $\mathbb{R}$ par $f(x)=\dfrac{2x+1}{x^2+x+3}$.
\begin{enumerate}
    \item Vérifier que le dénominateur ne s'annule jamais sur $\mathbb{R}$.
    \item Montrer que le numérateur est la dérivée du dénominateur.
    \item En déduire une primitive de $f$ sur $\mathbb{R}$.
    \item Calculer $\displaystyle\int_0^1 f(x)\,dx$.
\end{enumerate}
\end{exercice}

\courssub{Je me prépare au Probatoire}

\begin{exercice}[Vidange d'un réservoir -- Modélisation intégrale]
Un réservoir d'eau cubique de côté $1\,\text{m}$ (volume total $1\,000\,\text{L}$) se vide par un orifice au fond. Le débit instantané de sortie (en L/min) est modélisé par la fonction $v(t)=100-5t$ pour $t\in[0;20]$, où $t$ est le temps en minutes. On note $V(t)$ le volume d'eau restant dans le réservoir à l'instant $t$ (en litres), avec $V(0)=1000$.
\begin{enumerate}
    \item Justifier que la fonction $v$ est continue sur $[0;20]$.
    \item Exprimer $V(t)$ à l'aide d'une primitive de $v$. Déterminer l'expression explicite de $V(t)$.
    \item Calculer le volume total d'eau écoulé entre $t=0$ et $t=20\,\text{min}$.
    \item Déterminer l'instant $t_0$ auquel le réservoir est à moitié vide.
    \item Vérifier que le réservoir est vide à $t=20\,\text{min}$. Commenter la validité du modèle au-delà de $t=20$.
\end{enumerate}
\end{exercice}

\begin{exercice}[Fonction exponentielle composée -- Existence et expression intégrale]
Soit la fonction $f$ définie sur $[0;1]$ par $f(x)=e^{x^2}$.
\begin{enumerate}
    \item Montrer que $f$ est continue sur $[0;1]$.
    \item En déduire, en citant le théorème exact, que $f$ admet des primitives sur $[0;1]$.
    \item Donner une expression intégrale de la primitive $F$ de $f$ sur $[0;1]$ qui s'annule en $0$.
    \item On admet qu'on ne peut pas exprimer $F$ en fonctions élémentaires usuelles. Proposer une méthode numérique approchée pour calculer $F(1)$ (méthode des rectangles, trapèzes ou Simpson) et donner la formule pour un pas $h=0,25$.
\end{enumerate}
\end{exercice}

\begin{exercice}[Optimisation du coût moyen -- Synthèse]
Une entreprise de menuiserie industrielle a un coût marginal (en FCFA/unité) donné par $C'_m(x)=0,03x^2-0,4x+15$ pour une production de $x$ meubles ($x\ge 0$). Les coûts fixes s'élèvent à $2\,000$ FCFA.
\begin{enumerate}
    \item Déterminer la fonction coût total $C_t(x)$.
    \item Définir la fonction coût moyen unitaire $C_m(x)=\dfrac{C_t(x)}{x}$ pour $x>0$. Étudier les variations de $C_m$ sur $]0;+\infty[$.
    \item Déterminer le niveau de production $x_0$ qui minimise le coût moyen unitaire. Calculer ce coût minimal.
    \item Interpréter économiquement le résultat : que se passe-t-il si l'on produit moins de $x_0$ unités ? Plus de $x_0$ unités ?
    \item Tracer (schématiquement) les courbes de $C'_m$ et $C_m$ dans le même repère et montrer graphiquement le point d'intersection correspondant au minimum de $C_m$.
\end{enumerate}
\end{exercice}

\newpage

\coursec{Corrigés des exercices}

\coursec{Primitives d'une fonction continue sur un intervalle — Exercices corrigés et commentés}

\courssub{Exercices sur la définition et les propriétés algébriques}

\begin{corrige}[Exercice 1 — QCM : Primitive d'une fonction polynôme]
La bonne réponse est la \textbf{proposition 2} : $F(x)=x^3$.

\emph{Justification par dérivation de chaque proposition :}
\begin{enumerate}
    \item $F(x)=x^3+2x \Rightarrow F'(x)=3x^2+2 \neq f(x)$. Faux.
    \item $F(x)=x^3 \Rightarrow F'(x)=3x^2 = f(x)$. \textbf{Vrai.}
    \item $F(x)=6x \Rightarrow F'(x)=6 \neq f(x)$. Faux (c'est la dérivée de $f$, pas une primitive).
    \item $F(x)=x^3+x \Rightarrow F'(x)=3x^2+1 \neq f(x)$. Faux.
\end{enumerate}
\emph{Point de vigilance :} ne pas confondre dérivation et intégration ; pour vérifier qu'une fonction $F$ est une primitive de $f$, on \textbf{dérive} $F$ et on compare à $f$.
\end{corrige}

\begin{corrige}[Exercice 2 — Vrai ou Faux : Propriétés algébriques des primitives]
\begin{enumerate}
    \item \textbf{Vrai.} $G'(x) = F'(x) + 0 = f(x)$ car la dérivée d'une constante est nulle. Donc $G$ est une primitive de $f$ sur $I$.
    \item \textbf{Faux.} $H'(x) = 2F'(x) = 2f(x) \neq f(x)$ (sauf si $f$ est la fonction nulle). $H$ est une primitive de $2f$, pas de $f$.
    \item \textbf{Vrai.} $(F_1 - F_2)'(x) = F_1'(x) - F_2'(x) = f(x) - f(x) = 0$ pour tout $x \in I$. Une fonction de dérivée nulle sur un \emph{intervalle} est constante : il existe $k\in\mathbb{R}$ tel que $F_1(x)-F_2(x)=k$ pour tout $x\in I$.
\end{enumerate}
\emph{Point de vigilance :} pour la question 3, l'hypothèse « $I$ est un intervalle » est indispensable à citer : sur une réunion d'intervalles disjoints, la conclusion serait fausse.
\end{corrige}

\begin{corrige}[Exercice 3 — Définition et vocabulaire]
« Soit $f$ une fonction \textbf{continue} sur un \textbf{intervalle} $I$. On appelle \textbf{primitive} de $f$ sur $I$ toute fonction $F$ dérivable sur $I$ telle que $F'=f$ sur $I$. Si $F$ est une primitive de $f$, alors pour tout réel $k$, la fonction $x\mapsto F(x)+k$ est aussi une primitive de $f$. Réciproquement, deux primitives de $f$ sur $I$ ne diffèrent que d'une \textbf{constante}. »

\emph{Remarque :} le mot \cle{dérivée} n'apparaît pas dans les blancs ; il sert à définir la relation $F'=f$.
\end{corrige}

\courssub{Calcul de primitives : cas classiques et formes composées}

\begin{corrige}[Exercice 4 — Calcul direct de primitives]
\begin{enumerate}
    \item $f(x)=4x^3-2x+5$ sur $\mathbb{R}$. Par linéarité et avec $\int x^n dx = \frac{x^{n+1}}{n+1}$ ($n \neq -1$) :
    \[ F(x) = 4\cdot\frac{x^4}{4} - 2\cdot\frac{x^2}{2} + 5x + k = x^4 - x^2 + 5x + k,\quad k\in\mathbb{R}. \]
    \emph{Vérification :} $F'(x)=4x^3-2x+5=f(x)$. \checkmark
    \item $g(x)=\dfrac{1}{x^2}=x^{-2}$ sur $]-\infty;0[$ et sur $]0;+\infty[$. Avec $n=-2\neq -1$ :
    \[ G(x) = \frac{x^{-1}}{-1} + k = -\frac{1}{x} + k. \]
    \emph{Vérification :} $G'(x)=\dfrac{1}{x^2}=g(x)$. \checkmark\ (Une primitive sur chaque intervalle séparément.)
    \item $h(x)=\sqrt{x}=x^{1/2}$ sur $[0;+\infty[$. Avec $\alpha=\tfrac12$ :
    \[ H(x) = \frac{x^{3/2}}{3/2} + k = \frac{2}{3}x\sqrt{x} + k. \]
    \emph{Vérification :} $H'(x)=\frac{2}{3}\cdot\frac{3}{2}x^{1/2}=\sqrt{x}=h(x)$. \checkmark
    \item $k(x)=\cos(2x)$ sur $\mathbb{R}$. Forme $\cos(u(x))$ avec $u(x)=2x$, $u'(x)=2$. On écrit $k(x)=\frac{1}{2}\cdot 2\cos(2x)$ :
    \[ K(x) = \frac{1}{2}\sin(2x) + k. \]
    \emph{Vérification :} $K'(x)=\frac{1}{2}\cdot 2\cos(2x)=\cos(2x)$. \checkmark
\end{enumerate}
\emph{Point de vigilance :} pour les formes composées, ne pas oublier de diviser par le coefficient de $x$ (ici $2$) ; la vérification par dérivation permet de détecter l'erreur.
\end{corrige}

\begin{corrige}[Exercice 8 — Primitive d'un quotient de type $u'/u$]
\begin{enumerate}
    \item Soit $u(x)=x^2+x+3$. Discriminant : $\Delta = 1-12 = -11 < 0$. Le trinôme ne s'annule pas sur $\mathbb{R}$ et garde le signe de $a=1>0$ : $u(x)>0$ pour tout $x\in\mathbb{R}$. Donc $f$ est définie et continue sur $\mathbb{R}$.
    \item $u'(x)=2x+1$, qui est exactement le numérateur de $f$. Donc $f=\dfrac{u'}{u}$.
    \item Une primitive de $\dfrac{u'}{u}$ est $\ln|u|$. Comme $u(x)>0$ sur $\mathbb{R}$ :
    \[ \boxed{F(x)=\ln(x^2+x+3)+k,\quad k\in\mathbb{R}} \]
    \item 
    \begin{align*}
    \int_0^1 f(x)\,dx &= \left[\ln(x^2+x+3)\right]_0^1 = \ln(1+1+3)-\ln(3) = \ln 5 - \ln 3 = \ln\frac{5}{3}.
    \end{align*}
    Valeur approchée : $\ln\frac{5}{3}\approx 0{,}51$.
\end{enumerate}
\emph{Point de vigilance :} avant d'écrire $\ln(u)$ sans valeur absolue, il faut avoir justifié que $u>0$ sur l'intervalle (question 1).
\end{corrige}

\begin{corrige}[Exercice 10 — Fonction exponentielle composée : existence et approximation numérique]
\begin{enumerate}
    \item $f(x)=e^{x^2}$ est la composée de $x\mapsto x^2$ (continue sur $\mathbb{R}$) et de $x\mapsto e^x$ (continue sur $\mathbb{R}$) : $f$ est continue sur $\mathbb{R}$, donc sur $[0;1]$.
    \item D'après le \textbf{théorème fondamental} : \emph{toute fonction continue sur un intervalle admet au moins une primitive sur cet intervalle}. Comme $f$ est continue sur l'intervalle $[0;1]$, $f$ admet des primitives sur $[0;1]$.
    \item La primitive de $f$ qui s'annule en $0$ est :
    \[ F(x) = \int_0^x e^{t^2}\,dt. \]
    On a bien $F(0)=0$ et $F'(x)=e^{x^2}=f(x)$.
    \item \textbf{Méthode des rectangles} (rectangles à droite) avec pas $h=0{,}25$ : on subdivise $[0;1]$ en $4$ intervalles $[0;0{,}25]$, $[0{,}25;0{,}5]$, $[0{,}5;0{,}75]$, $[0{,}75;1]$. Formule des rectangles à droite :
    \[ F(1) = \int_0^1 e^{x^2}dx \approx h\sum_{k=1}^{4} f(kh) = 0{,}25\left[f(0{,}25)+f(0{,}5)+f(0{,}75)+f(1)\right]. \]
    Calcul : $f(0{,}25)=e^{0{,}0625}\approx 1{,}0645$ ; $f(0{,}5)=e^{0{,}25}\approx 1{,}2840$ ; $f(0{,}75)=e^{0{,}5625}\approx 1{,}7551$ ; $f(1)=e\approx 2{,}7183$.
    \[ F(1)\approx 0{,}25\times(1{,}0645+1{,}2840+1{,}7551+2{,}7183) = 0{,}25\times 6{,}8219 \approx 1{,}71. \]
    (Valeur exacte approchée : $F(1)\approx 1{,}4627$ ; la méthode des rectangles à droite surestime car $f$ est croissante ; un pas plus fin ou la méthode des trapèzes améliore la précision.)
\end{enumerate}
\textbf{Barème indicatif (sur 10) :} Q1 : 2 pts ; Q2 : 2 pts (citation exacte du théorème exigée) ; Q3 : 2 pts ; Q4 : 4 pts (formule 2 pts, application numérique 2 pts).

\textbf{Points de vigilance :} en Q2, le théorème doit être cité avec ses hypothèses (continuité \emph{sur un intervalle}) ; en Q3, la borne variable est $x$, la variable d'intégration doit être renommée $t$ ; en Q4, bien préciser la méthode choisie et le découpage.
\end{corrige}

\courssub{Applications concrètes : modélisation technique et économique}

\begin{corrige}[Exercice 5 — Coût total de production]
\begin{enumerate}
    \item $C'_m(x)=0{,}01x^2-0{,}5x+10$ est une fonction polynomiale, donc continue sur $[0;+\infty[$ : elle admet des primitives. Une primitive générale est :
    \[ C_t(x) = 0{,}01\cdot\frac{x^3}{3} - 0{,}5\cdot\frac{x^2}{2} + 10x + k = \frac{0{,}01}{3}x^3 - 0{,}25x^2 + 10x + k. \]
    Le coût fixe est $C_t(0)=5\,000$, donc $k=5\,000$ :
    \[ \boxed{C_t(x) = \frac{0{,}01}{3}x^3 - 0{,}25x^2 + 10x + 5\,000} \]
    \item Pour $x=100$ :
    \begin{align*}
    C_t(100) &= \frac{0{,}01}{3}(100)^3 - 0{,}25(100)^2 + 10(100) + 5\,000 \\
    &= \frac{10\,000}{3} - 2\,500 + 1\,000 + 5\,000 \\
    &\approx 3\,333{,}33 + 3\,500 = 6\,833{,}33 \text{ FCFA}.
    \end{align*}
    \item Le terme $0{,}01x^2$ traduit un \textbf{surcoût marginal croissant} : au-delà d'un certain niveau de production, produire une unité supplémentaire coûte de plus en plus cher (saturation des machines, heures supplémentaires, maintenance accrue).
\end{enumerate}
\end{corrige}

\begin{corrige}[Exercice 6 — Charge électrique traversant un conducteur]
\begin{enumerate}
    \item $i(t)=2e^{-0{,}5t}$ est de la forme $u'e^{u}$ à un coefficient près : $e^{-0{,}5t}$ a pour primitive $\dfrac{e^{-0{,}5t}}{-0{,}5}=-2e^{-0{,}5t}$. Donc :
    \[ q(t) = 2\times\left(-2e^{-0{,}5t}\right) + k = -4e^{-0{,}5t} + k. \]
    Condition $q(0)=0$ : $-4+k=0 \Rightarrow k=4$. D'où :
    \[ \boxed{q(t) = 4\left(1 - e^{-0{,}5t}\right)} \]
    \item Charge traversée entre $t=0$ et $t=2$ s :
    \[ q(2)-q(0) = 4\left(1-e^{-1}\right) \approx 4(1-0{,}368) \approx 4\times 0{,}632 \approx 2{,}53 \text{ C}. \]
    \item $\displaystyle\lim_{t\to+\infty} e^{-0{,}5t}=0$, donc $\displaystyle\lim_{t\to+\infty} q(t)=4$ C. \emph{Interprétation :} la charge totale qui peut traverser le conducteur est bornée ($4$ coulombs) ; le courant décroît exponentiellement vers $0$ (régime transitoire qui s'éteint, par exemple décharge d'un condensateur).
\end{enumerate}
\end{corrige}

\begin{corrige}[Exercice 7 — Distance parcourue par un mobile]
\begin{enumerate}
    \item $v(t)=10-t$ est une fonction affine, donc continue sur $\mathbb{R}$, en particulier sur $[0;10]$. D'après le théorème fondamental (\emph{toute fonction continue sur un intervalle admet une primitive}), $v$ admet des primitives sur $[0;10]$.
    \item Une primitive générale : $x(t)=10t-\dfrac{t^2}{2}+k$. Condition $x(0)=0 \Rightarrow k=0$ :
    \[ \boxed{x(t)=10t-\frac{t^2}{2}} \]
    \item Sur $[0;10]$, $v(t)=10-t\ge 0$ : le véhicule avance sans reculer. La distance parcourue est :
    \[ x(10)-x(0) = 10(10)-\frac{100}{2} = 100-50 = 50 \text{ m}. \]
    Le véhicule parcourt $50$ m avant de s'arrêter ($v(10)=0$).
\end{enumerate}
\emph{Point de vigilance :} la distance parcourue n'est égale à $x(t_2)-x(t_1)$ que si la vitesse garde un signe constant ; il faut le justifier (ici $v\ge 0$ sur $[0;10]$).
\end{corrige}

\begin{corrige}[Exercice 9 — Vidange d'un réservoir]
\begin{enumerate}
    \item $v(t)=100-5t$ est une fonction affine, donc continue sur $\mathbb{R}$, en particulier sur $[0;20]$.
    \item Le volume écoulé entre $0$ et $t$ est $\displaystyle\int_0^t v(s)\,ds$. Une primitive de $v$ est $100t-\dfrac{5t^2}{2}$. Le volume restant est donc :
    \[ V(t) = V(0) - \int_0^t v(s)\,ds = 1\,000 - \left(100t-\frac{5t^2}{2}\right) = 1\,000 - 100t + \frac{5}{2}t^2. \]
    \emph{Vérification :} $V'(t)=-100+5t=-v(t)$, cohérent : le volume diminue au rythme du débit de sortie.
    \item Volume écoulé entre $t=0$ et $t=20$ min :
    \[ \int_0^{20}(100-5t)\,dt = \left[100t-\frac{5t^2}{2}\right]_0^{20} = 2\,000 - \frac{5\times 400}{2} = 2\,000 - 1\,000 = 1\,000 \text{ L}. \]
    \item Réservoir à moitié vide : $V(t_0)=500$ :
    \[ 1\,000 - 100t_0 + \frac{5}{2}t_0^2 = 500 \iff \frac{5}{2}t_0^2 - 100t_0 + 500 = 0 \iff t_0^2 - 40t_0 + 200 = 0. \]
    $\Delta = 1\,600 - 800 = 800$, $\sqrt{800}=20\sqrt{2}\approx 28{,}28$. Solutions : $t_0 = \dfrac{40\pm 20\sqrt{2}}{2}=20\pm 10\sqrt{2}$. Seule la solution dans $[0;20]$ convient :
    \[ t_0 = 20 - 10\sqrt{2} \approx 5{,}86 \text{ min} \approx 5 \text{ min } 52 \text{ s}. \]
    \item $V(20)=1\,000-2\,000+\frac{5}{2}(400)=1\,000-2\,000+1\,000=0$ : le réservoir est vide à $t=20$ min. Au-delà, le modèle donnerait $v(t)<0$ (débit négatif, absurde) et $V(t)>0$ qui croîtrait : le modèle n'est \textbf{valable que sur $[0;20]$}.
\end{enumerate}
\textbf{Barème indicatif (sur 10) :} Q1 : 1 pt ; Q2 : 3 pts (primitive 1,5 pt + condition initiale 1,5 pt) ; Q3 : 2 pts ; Q4 : 2,5 pts (équation 1 pt, résolution et sélection de la racine 1,5 pt) ; Q5 : 1,5 pt.

\textbf{Points de vigilance :} citer la continuité avant d'intégrer ; bien distinguer volume \emph{écoulé} et volume \emph{restant} ; rejeter la racine hors de $[0;20]$ en le justifiant ; commenter le domaine de validité du modèle.
\end{corrige}

\begin{corrige}[Exercice 11 — Optimisation du coût moyen]
\begin{enumerate}
    \item $C'_m(x)=0{,}03x^2-0{,}4x+15$ est polynomiale, donc continue sur $[0;+\infty[$. Primitive générale :
    \[ C_t(x) = 0{,}01x^3 - 0{,}2x^2 + 15x + k. \]
    $C_t(0)=2\,000 \Rightarrow k=2\,000$ :
    \[ \boxed{C_t(x) = 0{,}01x^3 - 0{,}2x^2 + 15x + 2\,000} \]
    \item Pour $x>0$ : $C_m(x)=\dfrac{C_t(x)}{x} = 0{,}01x^2 - 0{,}2x + 15 + \dfrac{2\,000}{x}$.
    Dérivée : $C_m'(x) = 0{,}02x - 0{,}2 - \dfrac{2\,000}{x^2} = \dfrac{0{,}02x^3 - 0{,}2x^2 - 2\,000}{x^2}$.
    Posons $\varphi(x)=0{,}02x^3-0{,}2x^2-2\,000$ sur $]0;+\infty[$. $\varphi'(x)=0{,}06x^2-0{,}4x=0{,}02x(3x-20)$ : $\varphi$ décroît sur $]0;\tfrac{20}{3}]$ puis croît sur $[\tfrac{20}{3};+\infty[$. Or $\varphi(\tfrac{20}{3})<0$ et $\varphi(x)\to+\infty$ en $+\infty$, donc $\varphi$ s'annule une seule fois. Testons $x=50$ : $\varphi(50)=0{,}02(125\,000)-0{,}2(2\,500)-2\,000=2\,500-500-2\,000=0$. Donc $x_0=50$.
    $C_m'(x)<0$ sur $]0;50[$ et $C_m'(x)>0$ sur $]50;+\infty[$ : $C_m$ est décroissante puis croissante.
    \item Le minimum est atteint en $x_0=50$ meubles :
    \[ C_m(50) = 0{,}01(2\,500) - 0{,}2(50) + 15 + \frac{2\,000}{50} = 25 - 10 + 15 + 40 = 70 \text{ FCFA/unité}. \]
    \item \emph{Interprétation :} si l'on produit moins de $50$ meubles, les coûts fixes ($2\,000$ FCFA) sont mal amortis et le coût unitaire est supérieur à $70$ FCFA ; si l'on produit plus de $50$ meubles, le surcoût marginal (termes en $x^2$) fait remonter le coût unitaire. Le niveau $x_0=50$ est la \textbf{taille optimale de production}.
    \item Propriété classique : au minimum du coût moyen, $C_m(x_0)=C'_m(x_0)$. Vérification : $C'_m(50)=0{,}03(2\,500)-0{,}4(50)+15=75-20+15=70=C_m(50)$. \checkmark
    \begin{popfigure}
    \begin{center}
    \begin{tikzpicture}[scale=0.9]
    \draw[->,popdark] (0,0) -- (8.5,0) node[right]{$x$ (dizaines d'unités)};
    \draw[->,popdark] (0,0) -- (0,5.5) node[above]{Coût (dizaines de FCFA)};
    % C_m(x) = 0.01x^2 - 0.2x + 15 + 2000/x. En dizaines : x_tikz = x/10, y_tikz = y/10.
    % C_m_tikz(x) = (0.01*(10x)^2 - 0.2*(10x) + 15 + 2000/(10x)) / 10 = (x^2 - 2x + 15 + 200/x)/10 = 0.1x^2 - 0.2x + 1.5 + 20/x
    \draw[popblue,thick,domain=0.5:8,samples=200] plot (\x,{0.1*\x*\x - 0.2*\x + 1.5 + 20/\x});
    % C'_m(x) = 0.03x^2 - 0.4x + 15. C'_m_tikz(x) = (0.03*(10x)^2 - 0.4*(10x) + 15)/10 = (3x^2 - 4x + 15)/10 = 0.3x^2 - 0.4x + 1.5
    \draw[poporange,thick,domain=0:8,samples=100] plot (\x,{0.3*\x*\x - 0.4*\x + 1.5});
    \draw[dashed,popmuted] (5,0) node[below]{$50$} -- (5,3.5);
    \draw[dashed,popmuted] (0,3.5) node[left]{$70$} -- (5,3.5);
    \fill[popink] (5,3.5) circle (2pt);
    \node[popblue] at (7.5,4.8) {$C_m$};
    \node[poporange] at (7.5,2.5) {$C'_m$};
    \end{tikzpicture}
    \end{center}
    \legende{Le coût marginal coupe le coût moyen en son minimum ($x_0=50$, $70$ FCFA).}
    \end{popfigure}
\end{enumerate}
\textbf{Barème indicatif (sur 20) :} Q1 : 3 pts ; Q2 : 6 pts (expression 2 pts, dérivée 2 pts, sens de variation 2 pts) ; Q3 : 4 pts ; Q4 : 3 pts ; Q5 : 4 pts (tracé 2 pts, intersection au minimum 2 pts).

\textbf{Points de vigilance :} ne pas oublier les coûts fixes dans $C_t$ ; le domaine de $C_m$ est $]0;+\infty[$ (division par $x$) ; la justification du signe de $C_m'$ doit passer par l'étude du numérateur ; au Probatoire, la vérification $C_m(x_0)=C'_m(x_0)$ est un argument attendu pour le graphique.
\end{corrige}

\newpage

\coursec{Fiche bilan}

\begin{popfigure}
\begin{tikzpicture}[
  mindmap,
  grow cyclic,
  every node/.style={concept, execute at begin node=\small\bfseries},
  concept color=popblue,
  level 1/.append style={level distance=4.5cm, sibling angle=360/6},
  level 2/.append style={level distance=3cm, sibling angle=360/8},
  branch/.style={concept color=#1}
]
\node[concept, text width=3.2cm, align=center, font=\large\bfseries] {Primitives\\Fonction continue}
  child[branch=poporange] { node[concept, text width=2.5cm, align=center] {Définition\\$F'=f$ sur $I$} }
  child[branch=popgreen] { node[concept, text width=2.5cm, align=center] {Existence\\Théorème fondamental} }
  child[branch=poppurple] { node[concept, text width=2.5cm, align=center] {Propriétés\\Linéarité, $+C$} }
  child[branch=poppink] { node[concept, text width=2.5cm, align=center] {Table usuelles\\$x^n$, $1/x$, $\mathrm{e}^x$, $\sin$, $\cos$} }
  child[branch=popgold] { node[concept, text width=2.5cm, align=center] {Calcul\\Linéarité, $u^n u'$, $u'/u$} }
  child[branch=popdark] { node[concept, text width=2.5cm, align=center] {Applications\\Aire, Physique, Technique} };
\end{tikzpicture}
\legende{Carte mentale : Primitives d'une fonction continue sur un intervalle (Terminale Technique)}
\end{popfigure}

\begin{bilan}

\courssub{Définitions et théorème fondamental}

\begin{definition}[Primitive]
Soit $f$ une fonction définie sur un intervalle $I$. Une fonction $F$ dérivable sur $I$ est une \textbf{primitive} de $f$ sur $I$ si et seulement si :
\[
\forall x \in I,\quad F'(x) = f(x).
\]
\end{definition}

\begin{propriete}[Théorème d'existence -- Exigible au Probatoire]
Toute fonction \textbf{continue} sur un intervalle $I$ admet au moins une primitive sur $I$.
\end{propriete}

\begin{propriete}[Famille des primitives]
Si $F$ est une primitive de $f$ sur $I$, alors l'ensemble des primitives de $f$ sur $I$ est :
\[
\{ F + C \mid C \in \mathbb{R} \}.
\]
Graphiquement, ce sont des courbes obtenues par translation verticale.
\end{propriete}

\courssub{Propriétés algébriques (Linéarité)}

\begin{propriete}
Soient $f,g$ continues sur $I$, $F,G$ primitives respectives, $\lambda, \mu \in \mathbb{R}$.
\begin{itemize}
  \item $\lambda F + \mu G$ est une primitive de $\lambda f + \mu g$ sur $I$.
  \item Si $u$ est dérivable sur $I$ et $n \in \mathbb{N}^*$, une primitive de $u^n u'$ est $\dfrac{u^{n+1}}{n+1}$.
  \item Si $u$ est dérivable et strictement positive (ou négative) sur $I$, une primitive de $\dfrac{u'}{u}$ est $\ln|u|$.
\end{itemize}
\end{propriete}

\courssub{Table des primitives usuelles (à connaître par cœur)}

\begin{center}
\begin{tabularx}{\linewidth}{|>{\columncolor{popblueL}}c|X|>{\columncolor{popblueL}}c|X|}\hline
\textbf{Fonction $f(x)$} & \textbf{Primitive $F(x)$ (sur $I$ adapté)} & \textbf{Fonction $f(x)$} & \textbf{Primitive $F(x)$ (sur $I$ adapté)} \\\hline
$k$ ($k \in \mathbb{R}$) & $kx$ & $\cos x$ & $\sin x$ \\\hline
$x^n$ ($n \in \mathbb{N}$, $n \neq -1$) & $\dfrac{x^{n+1}}{n+1}$ & $\sin x$ & $-\cos x$ \\\hline
$x^\alpha$ ($\alpha \in \mathbb{Q}\setminus\{-1\}$) & $\dfrac{x^{\alpha+1}}{\alpha+1}$ ($x>0$) & $\dfrac{1}{1+x^2}$ & $\arctan x$ \\\hline
$\dfrac{1}{x}$ & $\ln|x|$ ($x \neq 0$) & $\dfrac{1}{\sqrt{1-x^2}}$ & $\arcsin x$ ($|x|<1$) \\\hline
$\mathrm{e}^x$ & $\mathrm{e}^x$ & $\dfrac{1}{\sqrt{1-x^2}}$ & $-\arccos x$ ($|x|<1$) \\\hline
$a^x$ ($a>0, a\neq 1$) & $\dfrac{a^x}{\ln a}$ & $\dfrac{u'}{u}$ & $\ln|u|$ \\\hline
$\cosh x$ & $\sinh x$ & $u^n u'$ ($n\neq -1$) & $\dfrac{u^{n+1}}{n+1}$ \\\hline
$\sinh x$ & $\cosh x$ & & \\\hline
\end{tabularx}
\end{center}

\courssub{Méthodes express}

\begin{methode}[Déterminer une primitive sur un intervalle]
\begin{enumerate}
  \item \textbf{Vérifier la continuité} de $f$ sur l'intervalle $I$ considéré.
  \item \textbf{Décomposer} $f$ en somme/différence de termes usuels (linéarité).
  \item \textbf{Identifier} les formes composées : $u^n u'$ ou $u'/u$ (chercher $u$ et $u'$).
  \item \textbf{Écrire} la primitive $F(x)$ en ajoutant \textbf{$+ C$} ($C \in \mathbb{R}$).
  \item \textbf{Vérifier} en dérivant $F$ : on doit retrouver $f$.
  \item Si condition initiale ($F(x_0)=y_0$) : \textbf{calculer $C$} et donner l'unique primitive.
\end{enumerate}
\end{methode}

\begin{attention}[Pièges fréquents]
\begin{itemize}
  \item Oublier la constante $C$ (famille infinie de primitives).
  \item Appliquer la formule $x^n \to x^{n+1}/(n+1)$ pour $n=-1$ (donner $\ln|x|$).
  \item Négliger le domaine de définition (ex: $\ln|x|$ pour $1/x$ sur $\mathbb{R}^*$, pas sur $\mathbb{R}$).
  \item Confondre primitive et intégrale définie (aire = différence de primitives).
  \item Mauvaise identification de $u$ dans $u^n u'$ : vérifier que le facteur restant est bien $u'$ (ou $k u'$).
\end{itemize}
\end{attention}

\end{bilan}

\begin{aretenir}[Checklist avant le Probatoire]
\begin{itemize}
  \item $\square$ Je sais donner la définition rigoureuse d'une primitive ($F'=f$ sur $I$).
  \item $\square$ Je connais et sais énoncer le théorème d'existence pour les fonctions continues.
  \item $\square$ Je sais que deux primitives d'une même fonction ne diffèrent que d'une constante.
  \item $\square$ Je connais par cœur la table des primitives usuelles (puissances, $1/x$, exp, trig, ln, arctan).
  \item $\square$ J'applique correctement la linéarité : primitive de somme = somme des primitives.
  \item $\square$ Je reconnais et traite la forme $u^n u'$ ($n \neq -1$) $\to \frac{u^{n+1}}{n+1}$.
  \item $\square$ Je reconnais et traite la forme $u'/u$ $\to \ln|u|$.
  \item $\square$ Je sais déterminer la constante $C$ à l'aide d'une condition initiale ($F(x_0)=y_0$).
  \item $\square$ Je vérifie systématiquement ma primitive en la dérivant.
  \item $\square$ Je précise toujours l'intervalle $I$ sur lequel la primitive est définie.
  \item $\square$ Je sais modéliser un problème technique (vitesse $\to$ position, courant $\to$ charge, débit $\to$ volume).
  \item $\square$ Je distingue primitive (famille de fonctions) et intégrale définie (nombre/aire).
\end{itemize}
\end{aretenir}

\findecours

\end{document}
